% Full manuscript: reworking-v3, 29 October 2026.
% Polyhedral transfer; fixed 2p faces; one smoother and one path in support coordinates.
% Existing headline bounds and interpolation identities retained.
% Source: preamble
\documentclass[11pt,letterpaper]{article}
\usepackage[margin=1.05in]{geometry}
\usepackage[T1]{fontenc}
\usepackage{lmodern,microtype}
\usepackage{amsmath,amssymb,amsthm,mathtools}
\usepackage{longtable}
\usepackage[hidelinks,hypertexnames=false,linktoc=all]{hyperref}
\newcommand{\Ac}{\mathsf A}
\usepackage{fancyhdr}
\pagestyle{fancy}
\fancyhf{}
\fancyhead[L]{\small Gaussian approximation over rectangles}
\fancyhead[R]{\small\thepage}
\setlength{\headheight}{14pt}
\hypersetup{pdftitle={High-Dimensional Gaussian Approximations over Rectangles without Covariance Restrictions},pdfsubject={Bounded independent summands; rates optimal up to logarithmic factors},pdfauthor={Victor Chernozhukov, Denis Chetverikov, Kengo Kato, Yuta Koike}}
\setlength{\parindent}{1.25em}
\setlength{\parskip}{0pt}
\emergencystretch=2em
\allowdisplaybreaks[2]
\raggedbottom
\clubpenalty=10000
\widowpenalty=10000
\newtheorem{theorem}{Theorem}
\newtheorem{proposition}{Proposition}
\newtheorem{lemma}{Lemma}
\newtheorem{corollary}{Corollary}
\theoremstyle{remark}
\newtheorem{remark}{Remark}
\newcommand{\E}{\mathbb E}
\newcommand{\Pp}{\mathbb P}
\newcommand{\R}{\mathbb R}
\newcommand{\cR}{\mathcal R}
\newcommand{\cF}{\mathcal F}
\newcommand{\cP}{\mathcal P}
\newcommand{\cD}{\mathcal D}
\newcommand{\Cov}{\operatorname{Cov}}
\newcommand{\Var}{\operatorname{Var}}
\newcommand{\tr}{\operatorname{tr}}
\newcommand{\op}{\mathrm{op}}
\newcommand{\norm}[1]{\left\lVert#1\right\rVert}
\newcommand{\rh}{\rho_{\mathcal R}}
\newcommand{\ind}{\mathbf 1}
\newcommand{\dd}{\,d}
\newcommand{\Hh}{H_2}
\newcommand{\surf}{\mathcal J}

\newcommand{\logp}{\log_+}

% Source: front
\begin{document}
\begin{center}
{\Large High-Dimensional Gaussian Approximations\\[0.15em]
over Rectangles without Covariance Restrictions}\\[0.85em]
{\normalsize Victor Chernozhukov, Denis Chetverikov, Kengo Kato, Yuta
Koike\renewcommand{\thefootnote}{\fnsymbol{footnote}}\footnote{The paper
was developed through collaboration with the AI systems GPT-6 Astra,
GPT-6 Sol, and Claude Opus.}}\\[0.65em]
{\normalsize 29 October 2026}
\end{center}
\setcounter{footnote}{0}
\thispagestyle{plain}

\begin{abstract}
 This paper studies the Gaussian approximation of sums of independent bounded random
vectors over rectangles, with no assumption on the correlation matrix,
which may be singular. Let \(X_1,\ldots,X_n\) be independent centered
random vectors in \(\R^p\) whose coordinates \(X_{ij}\) are bounded in
absolute value by a constant~\(B\), and suppose that each coordinate of
\(W_n\coloneqq n^{-1/2}(X_1+\cdots+X_n)\) has unit variance. Let \(Z_n\) be a
centered Gaussian vector with the covariance matrix of \(W_n\), and put
\(L\coloneqq \log(ep)\) and \(\log_+x\coloneqq \max\{0,\log x\}\). The largest difference
between \(\Pp(W_n\in A)\) and \(\Pp(Z_n\in A)\) over rectangles \(A\) is
at most
\[
 \frac{BL^{3/2}}{\sqrt n}\sqrt{1+\log_+\sqrt{n/L}},
\]
up to a numerical constant. If, for a fixed \(c_0>0\), no two
coordinates have correlation larger than \(1-c_0\) in absolute value, the
bound improves slightly. When \(p\) grows
polynomially in~\(n\), the two bounds are of order \(n^{-1/2}(\log n)^2\) and
\(n^{-1/2}(\log n)^{3/2}\), respectively. A further refinement holds if the
third moments cancel in aggregate, that is, if
\(\sum_i\E[X_{ij}X_{ik}X_{i\ell}]=0\) for all coordinates \(j,k,\ell\).
The factor \(BL^{3/2}/\sqrt n\) cannot be improved in general, so the
second rate is then optimal and the first is optimal up to the factor
\((\log n)^{1/2}\).

\end{abstract}

\clearpage
\tableofcontents
\clearpage

\section{\texorpdfstring{Introduction and main results}{Introduction and main results}}\label{sec:results}
Gaussian approximation over rectangles provides a common foundation for
approximating maxima and constructing simultaneous confidence regions, with
applications to a wide variety of statistical inference problems reviewed in~\cite{CCKK2023}.
In many of these applications the correlation matrix may be nearly or
exactly singular, resulting in arbitrarily small eigenvalues. This motivates the
Gaussian approximations of this paper, whose quality does not depend on the
eigenvalues.

Consider a sum of independent centered bounded random vectors in
\(\R^p\) with marginal unit variance and covariance \(\Sigma\), and let
\(Z\sim N(0,\Sigma)\) be the ``approximating'' normal vector that shares the
mean and covariance of the sum:
\begin{equation}
 \begin{gathered}
 W\coloneqq (W_j)_{j=1}^p=\sum_{i=1}^n\xi_i,\qquad
 \E\xi_i=0,\qquad \norm{\xi_i}_\infty\le\delta\quad\text{a.s.},\\
 \Sigma\coloneqq \Cov(W),\qquad\Sigma_{jj}=1\ (j\le p),
 \qquad Z\sim N(0,\Sigma).
 \end{gathered}
 \label{eq:setup}
\end{equation}
The summands need not have the same distribution, and coordinates within
a summand may be dependent. The unit marginal variances fix the scale
of each coordinate; they can be obtained by coordinatewise rescaling
whenever all marginal variances are positive; there is no
restriction on the rank of \(\Sigma\).

The parameter \(\delta\) bounds the contribution of any one summand
to any coordinate of \(W\), and such a uniform bound is called an
\emph{envelope}. To state the
bound compactly, put
\[
 L\coloneqq \log(ep),\qquad
 D_L\coloneqq 1+\log_+\frac1{\delta\sqrt L},\qquad
 \log_+x\coloneqq \max\{0,\log x\}.
\]
Here \(L\) encodes the dependence on dimension, while \(D_L\)
contributes an additional logarithmic factor. Constants denoted by
\(C\) are absolute unless a subscript records their dependence.

\indent
Let \(\mathcal R\) denote the class of axis-aligned hyper-rectangles in
\(\R^p\). We measure the approximation error by the Kolmogorov distance
between the distributions of the random vectors \(W\) and \(Z\):
\[
 \rh(W,Z)\coloneqq \sup_{A\in\mathcal R}
       |\Pp(W\in A)-\Pp(Z\in A)|.
\]
The same notation is used for any two random vectors in \(\R^p\). For
example, taking \(A\coloneqq (-\infty,s]^p\) shows that this distance controls the
error in approximating the distribution of \(\max_{j\le p}W_j\), the maximum
coordinate of \(W\).

The first bound of this paper depends only on the summand envelope and the
dimension, so it holds uniformly as covariance eigenvalues approach or
reach zero.

\begin{theorem}[Unrestricted covariance]\label{thm:unrestricted}
Under~\eqref{eq:setup},
\begin{equation}
 \rh(W,Z)\le C\min\{1,\delta L^{3/2}\sqrt{D_L}\}.
 \label{eq:unrestricted}
\end{equation}
\end{theorem}

Theorem~\ref{thm:unrestricted} allows any rank of \(\Sigma\), perfect
correlations, and singular principal submatrices, and for normalized
observations it gives the following dependence on the sample size.

\begin{corollary}[Normalized observations]\label{cor:observations}
Let \(X_1,\ldots,X_n\in\R^p\) be independent and centered. Suppose
\(\norm{X_i}_\infty\le B\) almost surely and
\(n^{-1}\sum_i\E X_{ij}^2=1\) for every \(j\). Then, with
\(W_n\coloneqq n^{-1/2}\sum_iX_i\) and \(Z_n\sim N(0,\Cov(W_n))\),
\begin{equation}
 \rh(W_n,Z_n)\le C\min\!\left\{1,
 \frac{B L^{3/2}}{\sqrt n}
 \sqrt{1+\log_+\frac{\sqrt n}{B\sqrt L}}\right\}.
 \label{eq:observations-unrestricted}
\end{equation}
\end{corollary}

Corollary~\ref{cor:observations} follows from Theorem~\ref{thm:unrestricted}
by setting \(\xi_i\coloneqq X_i/\sqrt n\), which has envelope \(\delta=B/\sqrt n\). For
fixed \(B\) and \(p\) growing polynomially in~\(n\), its bound has
order \(n^{-1/2}(\log n)^2\), so the power \(n^{-1/2}\) of the
sample size is reached up to logarithmic factors, however small the
eigenvalues of \(\Sigma\) are.

The factor \(BL^{3/2}/\sqrt n\) is sharp. For \(B\ge2\), \(p\to\infty\) and
\(BL^{3/2}/\sqrt n\to0\), Theorem~\ref{opt:skew-lower} gives i.i.d.\ observations
with independent coordinates whose error is at least of order
\(BL^{3/2}/\sqrt n\). The example is inspired by that of Chernozhukov,
Chetverikov, and Koike~\cite[Proposition~2.1]{CCK}, and it drops some of their
side conditions on \(B\) and \(p\). Thus, in this regime, the
bound~\eqref{eq:observations-unrestricted} is optimal up to the logarithmic
factor \(\sqrt{1+\log_+(\sqrt n/(B\sqrt L))}\); when \(p\) grows polynomially
in~\(n\) and \(B\) is fixed, this factor is of order \((\log n)^{1/2}\). Whether
this factor can be improved remains an open question.

For unequal marginal variances, let
\[
 \sigma_j^2\coloneqq n^{-1}\sum_i\E X_{ij}^2,
 \qquad
 \underline\sigma\coloneqq \min_j\sigma_j>0.
\]
We then apply Theorem~\ref{thm:unrestricted} to the summands
\(\xi_{ij}\coloneqq X_{ij}/(\sigma_j\sqrt n)\). If
\(\max_i\norm{X_i}_\infty\le K_n\), these summands have
envelope \(\delta=K_n/(\underline\sigma\sqrt n)\).
Positive diagonal scaling maps rectangles to rectangles and so preserves
the distance \(\rh\), and the resulting bound also holds in the
original coordinates.

\subsection{Correlation gaps and third-moment cancellation}
We can obtain a better bound when no two distinct coordinates are nearly
equal or opposite. To measure this separation, define
\[
 c\coloneqq 1-\max_{j\ne k}|\Sigma_{jk}|,
 \qquad c\coloneqq 1\ \text{when }p=1,
\]
and put
\begin{equation}
 \kappa\coloneqq \begin{cases}\min\{L,c^{-1}\},&c>0,\\L,&c=0,\end{cases}
 \qquad E_\kappa\coloneqq 1+\log_+\frac{\sqrt\kappa}{\delta L}.
 \label{eq:parameters}
\end{equation}
A positive gap excludes equal or opposite coordinates; it still permits
singular covariance matrices.

\begin{theorem}[Gap-sensitive bound]\label{thm:optimized}
Under~\eqref{eq:setup},
\begin{equation}
 \rh(W,Z)\le C\min\!\left\{1,
 \delta[L^{3/2}+L\sqrt{\kappa E_\kappa}]\right\}.
 \tag{R*}\label{eq:optimized}
\end{equation}
\end{theorem}

Since \(\kappa\le L\) and \(E_\kappa\le D_L\),
Theorem~\ref{thm:optimized} implies Theorem~\ref{thm:unrestricted};
the term \(\delta L^{3/2}\) in~\eqref{eq:optimized} comes from the
third moments of the summands,
and it disappears when the full third-moment tensor cancels.

\begin{theorem}[Cancellation of third moments]\label{thm:zero-skew}
Under~\eqref{eq:setup}, suppose
\begin{equation}
 \sum_i\E\xi_i^{\otimes3}=0,
 \quad\text{equivalently,}\quad
 \sum_i\E[\xi_{ij}\xi_{ik}\xi_{i\ell}]=0
 \quad(j,k,\ell\le p).
 \label{eq:zero-third}
\end{equation}
Then
\begin{equation}
 \rh(W,Z)\le C\min\{1,\delta L\sqrt{\kappa E_\kappa}\}.
 \tag{ZS}\label{eq:zero-skew}
\end{equation}
\end{theorem}

Condition~\eqref{eq:zero-third} holds when each summand is centrally
symmetric, that is, when \(\xi_i\) and \(-\xi_i\) have the same law. Because it
involves the mixed third moments \(\E[\xi_{ij}\xi_{ik}\xi_{i\ell}]\), vanishing
marginal third moments do not suffice for this cancellation.

For a fixed standardized envelope and \(c\ge c_0>0\), the sample-size
bounds become
\begin{align}
 \rh(W_n,Z_n)&\le C_{B,c_0}\min\!\left\{1,
 \frac{L^{3/2}}{\sqrt n}
 +\frac L{\sqrt n}\sqrt{1+\log_+\frac{\sqrt n}{L}}\right\},
 \label{eq:observations-fixed}\\
 \rh(W_n,Z_n)&\le C_{B,c_0}\min\!\left\{1,
 \frac L{\sqrt n}\sqrt{1+\log_+\frac{\sqrt n}{L}}\right\}
 \quad\text{under~\eqref{eq:zero-third}}.
 \label{eq:zero-fixed}
\end{align}
In polynomial dimension the bound~\eqref{eq:observations-fixed}
has order \(n^{-1/2}(\log n)^{3/2}\), and it is optimal there, since the
lower-bound examples of~\cite[Proposition~2.1]{CCK} and
Theorem~\ref{opt:skew-lower} have independent coordinates and hence
correlation gap one. These bounds give uniform convergence at a fixed
positive gap when \(\log p=o(n^{1/3})\) in general and \(\log p=o(\sqrt n)\)
under cancellation. Theorem~\ref{opt:consistency} shows that these
conditions are sharp over fixed-envelope classes, with envelope at least
two for the first assertion and at least one for the second. In
the second case the sharpness concerns only the dimension threshold, and
below the threshold the error rate in~\eqref{eq:zero-fixed} is not known
to be sharp.

\subsection{Polyhedral sets}\label{sec:polyhedra}
A polyhedron is an intersection of halfspaces, so membership can be
tested by collecting the scalar products with its face normals into
a vector. The resulting event is a rectangle event in this new space.
To apply the rectangle bounds, we divide each scalar product by its
standard deviation. The condition below ensures that this normalization
leaves the transformed summands uniformly bounded; it does not require
the covariance matrix to be invertible.

\begin{corollary}[Gaussian approximation over polyhedra]\label{cor:polyhedra}
Let \(W\coloneqq\sum_{i=1}^n\xi_i\), where the vectors
\(\xi_i\in\R^p\) are independent and centered, with
\(\norm{\xi_i}_\infty\le\delta\) almost surely, for \(\delta>0\).
Set \(\Sigma\coloneqq\Cov(W)\) and let \(Z\sim N(0,\Sigma)\).
For an integer \(m\ge1\) and a constant \(K>0\), let
\(\mathcal P_m(K;\Sigma)\) be the class of polyhedra admitting a
representation
\[
 P=\bigcap_{\alpha=1}^q
 \{x\in\R^p:\nu_\alpha^\top x\le a_\alpha\},
 \qquad 1\le q\le m,
\]
whose normals satisfy
\[
 \sigma_\alpha\coloneqq
 \sqrt{\nu_\alpha^\top\Sigma\nu_\alpha}>0,
 \qquad \norm{\nu_\alpha}_1\le K\sigma_\alpha
 \quad(\alpha=1,\ldots,q).
\]
Put \(\Delta\coloneqq K\delta\) and \(L_m\coloneqq\log(em)\).
Then
\[
 \sup_{P\in\mathcal P_m(K;\Sigma)}
 |\Pp(W\in P)-\Pp(Z\in P)|
 \le C\min\!\left\{1,\,
 \Delta L_m^{3/2}
 \sqrt{1+\log_+\frac1{\Delta\sqrt{L_m}}}\right\},
\]
where \(C\) is an absolute constant. We include the empty set and
\(\R^p\) in the class. No restriction is imposed on the
rank of \(\Sigma\).
\end{corollary}

This transfer uses the reduction to face projections employed by
Chernozhukov, Chetverikov, and Kato~\cite[Section~3.1 and the proof of
Proposition~3.1]{CCK2017} in their extension from rectangles to simple
convex sets. Subsection~\ref{app:polyhedral-transfer} gives the proof
by applying Theorem~\ref{thm:unrestricted} to the standardized
projections.

The condition on the normals is unchanged if a defining inequality
is multiplied by a positive constant. Its role is to bound the size
of each standardized summand:
\[
 \left|\frac{\nu_\alpha^\top\xi_i}{\sigma_\alpha}\right|
 \le\frac{\norm{\nu_\alpha}_1\delta}{\sigma_\alpha}
 \le K\delta=\Delta.
\]
In particular, a uniform bound \(\norm{\nu_\alpha}_1\le B\),
together with \(\sigma_\alpha\ge s_0>0\), gives \(K=B/s_0\).
Thus the number of inequalities enters the bound through \(L_m\),
while the size of the standardized summands enters through \(\Delta\).
Zero-variance directions impose deterministic conditions and are
handled separately in Subsection~\ref{app:polyhedral-transfer}.

Theorems~\ref{thm:optimized} and~\ref{thm:zero-skew} also apply to
these standardized projections. For a representation with \(q\)
inequalities, their correlation matrix is
\[
 R_{\alpha\beta}\coloneqq
 \frac{\nu_\alpha^\top\Sigma\nu_\beta}
      {\sigma_\alpha\sigma_\beta}.
\]
The refinements use dimension \(q\), envelope \(\Delta\), and the
correlation gap of \(R\), which can differ from the gap of the
original coordinates. Aggregate third-moment cancellation is
preserved by the linear transformation. The same appendix subsection
verifies these assertions.

\subsection{Sub-exponential transfer}\label{sec:subexponential}
Truncation provides a second transfer: it extends the bounded theorems
to observations with sub-exponential tails. We first state the rectangle
bounds, then transfer them to polyhedra by the standardized face
projections of Subsection~\ref{sec:polyhedra}. The tail condition is
measured by the Orlicz norm
\[
 \|V\|_{\psi_1}\coloneqq \inf\{s>0:\E e^{|V|/s}\le2\}.
\]
If the norm is finite, \(V\) is said to be sub-exponential.

\begin{corollary}[Sub-exponential transfer]\label{cor:subexponential}
Let \(X_1,\ldots,X_n\in\R^p\) be independent and centered, with
\(n^{-1}\sum_i\E X_{ij}^2=1\) and
\(\max_{i,j}\|X_{ij}\|_{\psi_1}\le K\), where \(K\ge1\).
For \(W_n\coloneqq n^{-1/2}\sum_iX_i\) and a centered Gaussian \(Z_n\)
with the same covariance, put
\[
 \ell_n\coloneqq \log(en),\qquad b_n\coloneqq \frac{K\ell_n}{\sqrt n},\qquad
 D_n\coloneqq 1+\log_+\frac1{b_n\sqrt L}.
\]
Then
\begin{equation}
 \rh(W_n,Z_n)\le C\min\{1,b_nL^{3/2}\sqrt{D_n}\}.
 \label{eq:subexponential}
\end{equation}
If the correlation gap of \(W_n\) is at least a fixed \(c_0>0\), then
\begin{equation}
 \rh(W_n,Z_n)\le C_{c_0}\min\!\left\{1,
 b_n\left[L^{3/2}+L\sqrt{1+\log_+\frac1{b_nL}}\right]\right\}.
 \label{eq:subexponential-gap}
\end{equation}

For polyhedra, set \(\Sigma\coloneqq\Cov(W_n)\), fix \(J>0\), and
let \(\mathcal P\coloneqq\mathcal P_m(J;\Sigma)\) be the class in
Corollary~\ref{cor:polyhedra}, with at most \(m\) defining inequalities
and \(\norm{\nu_\alpha}_1\le J\sigma_\alpha\), where
\(\sigma_\alpha^2=\nu_\alpha^\top\Sigma\nu_\alpha>0\). Define
\[
 \rho_{\mathcal P}(W_n,Z_n)\coloneqq
 \sup_{P\in\mathcal P}|\Pp(W_n\in P)-\Pp(Z_n\in P)|,
 \qquad K_{\mathcal P}\coloneqq\max\{1,JK\}.
\]
Then~\eqref{eq:subexponential} also holds with
\(\rh\), \(K\), and \(L\) replaced throughout by
\(\rho_{\mathcal P}\), \(K_{\mathcal P}\), and
\(L_m=\log(em)\), respectively. The same replacement
in~\eqref{eq:subexponential-gap} holds for any subclass whose members
admit representations with
\[
 \left|\frac{\nu_\alpha^\top\Sigma\nu_\beta}
 {\sigma_\alpha\sigma_\beta}\right|\le1-c_0
 \qquad(\alpha\ne\beta).
\]
Here the gap condition is imposed on the standardized face projections.
\end{corollary}

The bound~\eqref{eq:subexponential} requires neither a lower bound on a
local determinant nor a lower bound on the positive eigenvalues of the
covariance matrix of \(W_n\). For unequal marginal variances, the tail
condition is imposed on \(X_{ij}/\sigma_j\), as in the standardization after
Corollary~\ref{cor:observations}. The polyhedral transfer follows from
\[
 \left\|\frac{\nu_\alpha^\top X_i}{\sigma_\alpha}\right\|_{\psi_1}
 \le \frac{K\norm{\nu_\alpha}_1}{\sigma_\alpha}
 \le JK\le K_{\mathcal P}.
\]
Thus the same tail argument applies after projection, with the number
of defining inequalities replacing the ambient dimension. A correlation
gap in the original coordinates does not by itself ensure a gap among
these projections.

For fixed \(K\) and \(p\asymp n^a\) with \(a>0\) fixed, the
bound~\eqref{eq:subexponential} is of order \(n^{-1/2}(\log n)^3\), and the
bound~\eqref{eq:subexponential-gap} under a fixed correlation gap is of order
\(n^{-1/2}(\log n)^{5/2}\), so the power \(n^{-1/2}\) of the sample size is
again optimal up to logarithmic factors. By~\eqref{eq:subexponential}, the error
tends to zero for fixed \(K\) when \(\log p=o(n^{1/3}/\log(en))\).
Appendix~\ref{app:subexponential} proves the rectangle bounds by
clipping each coordinate at a level of order \(K\log(en)\).
Exponential moments control the discarded sum, and a joint Gaussian
coupling restores the original covariance. For polyhedra, the proof
first forms the standardized face projections and applies this same
argument at tail scale \(K_{\mathcal P}\); no additional truncation
or Gaussian comparison is needed.

\subsection{Relation to recent work}\label{sec:literature}
High-dimensional Gaussian approximation for maxima and rectangles was
developed by Chernozhukov, Chetverikov, and Kato~\cite{CCK2013,CCK2017}.
Without covariance restrictions beyond positive marginal variances,
Chernozhukov, Chetverikov, Kato, and Koike~\cite[Theorem~2.1]{CCKK2022}
obtained an \(n^{-1/4}\) rate, up to logarithms, under coordinatewise
sub-exponential tails and an averaged fourth-moment condition.
The cube-root rectangle bound reported in~\cite[Theorem~1 and Section~5]{GaoZhang} improves the sample-size power to \(n^{-1/3}\). Under the tail assumption of Corollary~\ref{cor:subexponential}, in the unit-variance notation of this paper, it is
\begin{equation*}
 C\min\left\{1,
 \left(\frac{K^2}{n}\right)^{1/3}\{\log(2np)\}^{7/3}
 +\frac K{\sqrt n}\{\log(2np)\}^{5/2}\right\}.
\end{equation*}
No additional fourth-moment condition is needed for this bound.
For fixed \(K\) and \(p\asymp n^a\), with fixed \(a>0\), it has
order \(n^{-1/3}(\log n)^{7/3}\). In this dimension regime,
Corollaries~\ref{cor:observations} and~\ref{cor:subexponential}, with
fixed \(B\) or \(K\), give the sharp sample-size power
\(n^{-1/2}\), with logarithmic factors \((\log n)^2\) for bounded
observations and \((\log n)^3\) for sub-exponential coordinates.
These comparisons concern fixed envelopes or tail scales; the full
bounds retain their different dependence on these parameters.

A second line of work obtains rates close to \(n^{-1/2}\) under
covariance nondegeneracy. For bounded sums with positive definite
\(\Sigma\), the nondegenerate \(n^{-1/2}\) bound reported in~\cite[Theorem~1.1]{AL} removes the
\(\log(en)\) factor from the bound of Chernozhukov, Chetverikov,
and Koike~\cite[Corollary~2.1]{CCK}, and reads
\begin{equation*}
 \rh(W_n,Z_n)\le C\min\left\{1,
 \frac{BL^{3/2}}{\lambda_{\min}(\Sigma)\sqrt n}\right\}.
\end{equation*}
A fixed positive minimum eigenvalue therefore gives the sharp
\(BL^{3/2}/\sqrt n\) scale in the regimes of matching lower
bounds~\cite[Proposition~2.1]{CCK}.
The unrestricted bound of this paper does not involve \(\lambda_{\min}(\Sigma)\),
and it exceeds this scale only by the factor \(\sqrt{D_L}\).

Other results replace full nondegeneracy by conditions on small
groups of coordinates. Fang, Koike, Liu, and Zhao~\cite[Theorem~1.1]{FKLZ} obtain nearly
\(n^{-1/2}\) rates under sub-exponential tails and quantitative
two- and three-coordinate covariance conditions, allowing the full
covariance matrix to be singular. The companion paper~\cite{Local}
assumes uniformly positive eigenvalues for principal submatrices of
order at most four. For bounded observations it gives
\(C_aBL^{3/2}/\sqrt n\), and \(C_aBL/\sqrt n\) under aggregate
third-moment cancellation. The present results also cover singular
three- and four-coordinate submatrices. The two papers therefore
complement each other, and without the local condition the bound loses
only the factor \(\sqrt{D_L}\).

The bounds of this paper reach the rate \(n^{-1/2}\), up to logarithmic factors, with
no condition on the covariance matrix, and they lose only the factor
\(\sqrt{D_L}\) against the sharp bounds under nondegeneracy. They arise
from an interpolation between the sum and its Gaussian counterpart along
which the derivative of a smoothed rectangle probability is bounded
without reference to the eigenvalues of the covariance matrix.

Section~\ref{sec:smoothing} outlines the proof, derives
the rates from its main estimates, and explains why they reach the
power \(n^{-1/2}\), comparing the argument with those of \cite{AL} and~\cite{GaoZhang}; it also introduces the
notation used in the appendices. Section~\ref{opt:section}
discusses sharpness, and Subsection~\ref{sec:lean-certification}
records the formal verification in Lean.
Appendices~\ref{sec:main-proof}--\ref{app:subexponential}
contain all proofs of the upper bounds. Appendix~\ref{app:symmetric-lower}
proves the centrally symmetric lower bound. The Supplement~\cite{Supplement} collects the lower-bound calculations,
exact coordinate reductions, and further investigations.

\subsection{\texorpdfstring{Formal verification in Lean}{Formal verification in Lean}}\label{sec:lean-certification}
\leavevmode Theorems~\ref{thm:unrestricted}--\ref{opt:consistency} and the two rectangle bounds of
Corollary~\ref{cor:subexponential} were certified on 20 and 21~October 2026
in Lean~4.19.0, a proof assistant that checks formal proofs by computer,
with its library mathlib~v4.19.0, and the proofs use only Lean's standard
axioms. The Lean sources, the axiom audits, and the reviews that compare the
formal statements with the manuscript at that time are in the
\href{https://github.com/VC2015/HDCLTs/blob/bbf1ab99c733a398c181ca8c9f5a243a41f258c6/verification/2026-10-21-companion-results/README.md}{verification record}
of the repository, for source revision
\href{https://github.com/VC2015/HDCLTs/commit/6548201a5c5d9c132a53f1d504f7ef70346aa782}{\texttt{6548201a5c5d}}.
The certification refers to these Lean proofs, not to the written
derivations, and does not cover the Supplement as a whole.
Theorem~\ref{opt:symmetric-lower} has not been certified.
The revised proofs in \texttt{reworking-v2} and \texttt{reworking-v3}, and the polyhedral extensions in Corollaries~\ref{cor:polyhedra}
and~\ref{cor:subexponential}, have not been re-certified.
Updating the Lean development and checking its correspondence with this
version of the manuscript remain to be done.

\section{Outline of the proof}
\label{sec:smoothing}

We replace the indicator of a rectangle by a logistic approximation and
compare its expectations along a path that gradually replaces the bounded
summands with Gaussian noise while preserving the covariance. We regard
the interpolation parameter as time: the path starts at the bounded sum
at time zero and ends at the Gaussian vector at time one. We estimate
the derivatives using the sum of the summands' covariance matrices.
The resulting bound does not depend on the eigenvalues or rank of this
matrix and therefore applies even when it is singular.

Differentiating these expectations with respect to time produces
derivatives of the logistic approximation, which we estimate in two
ways. At early times, we bound these derivatives directly. At later
times, when the path contains enough Gaussian noise, we use Gaussian
integration by parts. This rewrites expectations involving third- and
fourth-order derivatives of the logistic approximation in terms of its
first- and second-order derivatives, multiplied by explicit functions
of the Gaussian noise. The lower derivatives grow more slowly as the
approximation becomes sharper, but we must also control the additional
Gaussian factors. In both estimates, we also use anti-concentration
inequalities, which control the probability of falling in a narrow
region near the boundary of a rectangle. These probability bounds
provide an additional small factor. The more delicate part of the
proof is to preserve this gain when applying integration by parts.

Subsection~\ref{sec:outline-smoothing} reduces the proof to bounded
rectangles, introduces the support coordinates, and bounds the error
of the logistic approximation.
Subsection~\ref{sec:outline-path} introduces the calibrated
interpolation path and states the derivative bounds.
Subsections~\ref{sec:outline-early} and~\ref{sec:outline-late} explain
how to bound the derivatives directly at early times and how to use
Gaussian integration by parts at later times. In
Subsection~\ref{sec:rate}, we integrate these bounds and choose the
smoothing scale to balance the smoothing and comparison errors.

The proofs are organized by these dependencies. Appendix~\ref{sec:main-proof}
derives the theorems from the smoothing and derivative bounds, then
transfers them to polyhedral sets.
Appendix~\ref{appx-sec:smoothing} proves the properties of the logistic
approximation, and Appendix~\ref{appx-sec:support} proves the derivative
bounds. Appendix~\ref{appx-app:concentration} gives the anti-concentration
argument, and Appendix~\ref{app:subexponential} transfers the bounds
to sub-exponential observations over rectangles and polyhedra.

\subsection{Support coordinates and the logistic approximation}
\label{sec:outline-smoothing}

We use bounded closed rectangles, allowing an interval with its lower
endpoint above its upper endpoint to make the rectangle empty.
Each coordinate contributes a lower and an upper inequality, giving
\(m\coloneqq2p\) inequalities in all. This restriction loses no
generality: for any possibly unbounded rectangle \(B\), the sets
\(B\cap[-M,M]^p\) increase to \(B\).
Continuity of probability from below therefore extends any uniform
bound for bounded rectangles to all rectangles.
Open or half-open intervals can likewise be exhausted by closed
intervals, so their endpoint conventions cause no additional restriction.

We now make a single change of coordinates for the proof. Its purpose
is to represent the Gaussian target by a standard normal vector and
to obtain a covariance identity that controls the sum of derivative
contributions. Let \(r\coloneqq\operatorname{rank}(\Sigma)\), and
choose a \(p\times r\) matrix \(\Gamma\) with linearly independent
columns such that \(\Sigma=\Gamma\Gamma^\top\). Each summand lies
almost surely in its range: if \(z\) is in the null space of
\(\Sigma\), independence and centering give
\(\sum_i\E(z^\top\xi_i)^2=z^\top\Sigma z=0\), so every
nonnegative term vanishes. Define
\[
 \zeta_i^0\coloneqq(\Gamma^\top\Gamma)^{-1}\Gamma^\top\xi_i,
 \qquad S\coloneqq\sum_i\zeta_i^0,
 \qquad G\sim N(0,I_r),
\]
where \(G\) is independent of the summands. The vectors
\(\zeta_i^0\in\R^r\) remain independent and centered, and
\[
 \xi_i=\Gamma\zeta_i^0,\qquad
 W=\Gamma S,\qquad Z=\Gamma G,\qquad
 \sum_i\E\zeta_i^0(\zeta_i^0)^\top=I_r.
\]
The last identity follows by applying the coordinate map to
\(\Sigma=\Gamma\Gamma^\top\). The superscript \(0\) distinguishes
these summands from their independent copies used later. From this
point on, the smoothing and interpolation arguments take place in
\(\R^r\).

To describe the sets in this space, let \(e_j\) be the \(j\)th
standard basis vector in \(\R^p\), and set
\[
 u_j\coloneqq\Gamma^\top e_j,\qquad
 u_{p+j}\coloneqq-\Gamma^\top e_j,
 \qquad j=1,\ldots,p.
\]
These \(m=2p\) vectors are the \emph{face normals}: they are the
outward normals of the defining inequalities after the coordinate
change. They are unit vectors because
\(\|\Gamma^\top e_j\|^2=\Sigma_{jj}=1\); repeated normals are
allowed. For a threshold vector \(a\in\R^m\), define
\[
 A_a\coloneqq\{x\in\R^r:u_\alpha\cdot x\le a_\alpha
                       \text{ for }\alpha=1,\ldots,m\},
\]
where \(\cdot\) denotes the usual inner product. A bounded rectangle
\(B=\prod_j[\ell_j,b_j]\), including the empty case when some
\(\ell_j>b_j\), corresponds to
\(a_j=b_j\) and \(a_{p+j}=-\ell_j\): then
\(\Gamma x\in B\) if and only if \(x\in A_a\). We allow all
finite threshold vectors; inconsistent inequalities give
\(A_a=\varnothing\). Thus the approximation error to be bounded is
\[
 \rho\coloneqq\sup_{a\in\R^m}
 |\Pp(S\in A_a)-\Pp(G\in A_a)|=\rh(W,Z).
\]
Each set is indexed by its threshold vector, which also specifies
the smoothing formula when some inequalities are redundant.

The summand bound survives along every face normal, since
\(u_j\cdot\zeta_i^0=\xi_{ij}\) and
\(u_{p+j}\cdot\zeta_i^0=-\xi_{ij}\). We measure these components
together by the \emph{face norm}
\[
 \norm x_U\coloneqq\max_{\alpha\le m}|u_\alpha\cdot x|,
 \qquad \norm{\zeta_i^0}_U\le\delta.
\]
This is a norm because \(\norm x_U=\norm{\Gamma x}_\infty\)
and \(\Gamma\) is injective. Its bound does not depend on how small
the positive eigenvalues of \(\Sigma\) are. This is the coordinate
reduction used in~\cite[Section~2.3]{GaoZhang}. The original
correlation gap \(c\) gives
\(u_\alpha\cdot u_\beta\le1-c\) for distinct face indices:
opposite faces have inner product \(-1\), and the other inner
products are signed correlations. The cancellation assumption
\eqref{eq:zero-third} is equivalently
\(\sum_i\E(\zeta_i^0)^{\otimes3}=0\), because the coordinate
map is invertible on the covariance support.

We now smooth these face inequalities. Let \(\ind_{A_a}\) denote
the indicator of \(A_a\), equal to one on the set and zero outside.
For \(h>0\), define its logistic approximation by
\begin{equation}
 F_{a,h}(x)\coloneqq
 \left(1+\sum_{\alpha=1}^m
 e^{(u_\alpha\cdot x-a_\alpha)/h}\right)^{-1},
 \qquad x\in\R^r.
 \label{eq:selector}
\end{equation}
For each fixed \(x\), this formula defines a probability distribution
over the \(m\) faces and an additional baseline option. Face
\(\alpha\) receives weight \(e^{(u_\alpha\cdot x-a_\alpha)/h}\),
while the baseline option receives weight one. After normalization,
the baseline probability is \(F_{a,h}(x)\), and the face probabilities
sum to \(1-F_{a,h}(x)\). We use these probabilities below to express
and bound the derivatives of the logistic approximation.

We compare the sum \(S\) and the standard Gaussian vector \(G\)
through the expectations of these smooth functions. Their largest
discrepancy is
\begin{equation}
 M_h\coloneqq\sup_{a\in\R^m}
 \left|\E F_{a,h}(S)-\E F_{a,h}(G)\right|.
 \label{eq:smoothed-discrepancy}
\end{equation}
The interpolation argument will bound \(M_h\). Combining this
bound with the error from replacing indicators by their logistic
approximations will then bound \(\rho\).

Membership in \(A_a\) requires the largest face-inequality violation
to be nonpositive. The logistic approximation smooths both the maximum
and the indicator that tests this condition. We can write
\[
 F_{a,h}(x)=g_h\!\left(h\log\sum_{\alpha=1}^m
 e^{(u_\alpha\cdot x-a_\alpha)/h}\right),
 \qquad g_h(s)\coloneqq\frac1{1+e^{s/h}}.
\]
The inner function smoothly approximates the maximum violation, while
\(g_h\) approximates the indicator of \(({-\infty},0]\). We use the
same scale \(h\) for both steps; the construction of Chernozhukov,
Chetverikov, and Kato~\cite{CCK2013} allows the cutoff function and
its scale to be chosen separately from the smooth maximum. The
resulting multinomial logistic form, with a baseline score of zero,
was used by McFadden~\cite[p.~1001]{McFadden1989} to smooth indicators
of simultaneous inequalities and also appears in~\cite[eq.~(5)]{GaoZhang}.
For every \(h>0\), the approximation is smooth, allowing us to
differentiate along the interpolation path. As \(h\) decreases to
zero, it approaches the indicator away from the boundary, but the
uniform bounds on its derivatives increase.

An anti-concentration inequality controls the probability of falling
in the boundary layer created by moving the thresholds. Write
\[
 A_a^s\coloneqq\{x:u_\alpha\cdot x\le a_\alpha+s
                         \text{ for all }\alpha\le m\}.
\]
The next lemma is the rectangle anti-concentration bound expressed
through the \(m=2p\) face inequalities. Recall that \(L=\log(ep)\),
so \(\log(em)\le(1+\log2)L\).
\begin{lemma}[Rectangle anti-concentration]\label{lem:concentration}
Let \(X_1,\ldots,X_n\in\R^r\) be independent centered random
vectors, and put \(V\coloneqq\sum_iX_i\). Suppose that
\[
 \norm{X_i}_U\le\eta\quad\text{almost surely for every }i,
 \qquad \Var(u_\alpha\cdot V)=v>0
                   \quad\text{for every }\alpha\le m.
\]
Let \(R\in\R^r\) be a random vector independent of the summands.
Then, for every \(a\in\R^m\) and \(s>0\),
\[
 \Pp(V+R\in A_a^s\setminus A_a^{-s})
 \le C\min\{1,(s+\eta)\sqrt{L/v}\}.
\]
If instead \(V\) is a centered Gaussian vector with
\(\Var(u_\alpha\cdot V)=v>0\) for every \(\alpha\), the same
bound holds with \(\eta=0\) for every \(R\) independent of \(V\).
The constant \(C\) is absolute.
\end{lemma}
Appendix~\ref{appx-app:concentration} proves the bounded-sum assertion
by conditioning on pairs of independent copies and applying the
Rademacher anti-concentration inequality of O'Donnell, Servedio,
and Tan~\cite{OST}, in the form recorded
in~\cite[Lemma~J.4]{CCKK2022}. This recovers the bounded-sum estimate
in~\cite[Lemma~2.1]{AL}, obtained using a bootstrap argument and
Kane's~\cite{Kane} noise-sensitivity inequality. The Gaussian assertion
is Nazarov's inequality applied to the vector
\((u_\alpha\cdot V)_{\alpha\le m}\);
see~\cite[Lemma~A.1]{CCK2017}.

Recall that \(\kappa=\min\{L,c^{-1}\}\), interpreted as
\(\kappa=L\) when \(c=0\). The following lemma bounds the error
of the logistic approximation and connects the smoothed discrepancy
\(M_h\) to the desired probability bound.
\begin{lemma}[Smoothing error]\label{lem:endpoint}
Under~\eqref{eq:setup}, for \(h\ge\delta\), uniformly in
\(a\in\R^m\),
\begin{equation}
 \E|F_{a,h}(S)-\ind_{A_a}(S)|+
 \E|F_{a,h}(G)-\ind_{A_a}(G)|\le Ch\sqrt L\,\kappa.
 \label{eq:endpoint-error}
\end{equation}
Consequently,
\begin{equation}
 \rho\le M_h+Ch\sqrt L\,\kappa.
 \label{eq:endpoint-min}
\end{equation}
\end{lemma}

To bound the smoothing error without assuming a positive correlation
gap, let \(\omega_0,\omega_1,\ldots,\omega_m\) be independent
exponential random variables with mean one. The logistic approximation
has the exact representation
\begin{equation}
 \begin{split}
 F_{a,h}(x)&=\Pp_\omega(x\in A_{a,\omega}),\\
 A_{a,\omega}&\coloneqq\{x:u_\alpha\cdot x\le
 a_\alpha+h\log(\omega_\alpha/\omega_0)
                   \text{ for all }\alpha\le m\}.
 \end{split}
 \label{eq:mixture}
\end{equation}
Thus the thresholds are shifted randomly while the face normals
remain fixed. Let
\[
 R_h\coloneqq\max_{\alpha\le m}
             |h\log(\omega_\alpha/\omega_0)|
\]
be the largest absolute change in a threshold. Each normalized change
has tails bounded by \(2e^{-s}\), so a union bound over the
\(m=2p\) faces and integration give \(\E R_h\le ChL\).
For fixed exponentials, the two indicators can differ only in
\(A_a^{R_h}\setminus A_a^{-R_h}\). Applying
Lemma~\ref{lem:concentration} and averaging over the exponentials
bounds the smoothing error by
\(C(\E R_h+\delta)\sqrt L\le ChL^{3/2}\), since \(h\ge\delta\).

When \(c>0\), separation of the face normals gives another bound
through the second derivatives of the potential
\[
 \Psi(x)\coloneqq h\log\!\left(1+\sum_{\alpha=1}^m
                         e^{(u_\alpha\cdot x-a_\alpha)/h}\right).
\]
The pointwise error satisfies
\( |F_{a,h}(x)-\ind_{A_a}(x)|\le(2h/c)\tr D^2\Psi(x)\),
where the trace is the sum of the second derivatives in the
\(r\) support coordinates. Lemma~\ref{appx-lem:endpoint-gap}
uses the covariance identity and Gaussian integration by parts to
bound the expected trace at \(S\) and \(G\) by \(C\sqrt L\).
This gives the alternative error bound \(Ch\sqrt L/c\).
Taking the smaller bound gives the factor \(\kappa\) in
Lemma~\ref{lem:endpoint}. Lemmas~\ref{appx-lem:endpoint-free}
and~\ref{appx-lem:endpoint-gap} give the details.

The comparison along the interpolation path defined in
Subsection~\ref{sec:outline-path} uses third- and fourth-order
derivatives of this same function. Abbreviate \(F_{a,h}\) to
\(F\) when the thresholds and width are fixed, and write
\[
 w_\alpha(x)\coloneqq e^{(u_\alpha\cdot x-a_\alpha)/h},\qquad
 \pi_\alpha\coloneqq\frac{w_\alpha}{1+\sum_\beta w_\beta},\qquad
 \theta\coloneqq1-F=\sum_\alpha\pi_\alpha.
\]
Thus \(\pi_\alpha\) is the logistic probability of face \(\alpha\)
and \(\theta\) is the total probability of the faces.

For a direction \(z\in\R^r\), the chain rule gives
\[
 DF(x)[z]= -\frac{F(x)}h
           \sum_{\alpha=1}^m\pi_\alpha(x)(u_\alpha\cdot z).
\]
Each component \(|u_\alpha\cdot z|\) is at most \(\norm z_U\).
Summing the face probabilities therefore yields
\[
 |DF(x)[z]|\le h^{-1}F(x)\theta(x)\norm z_U.
\]
Repeated differentiation introduces one factor \(h^{-1}\) at each
step. The remaining factors are \(F\), face probabilities, and
components of the derivative directions along face normals. Products
of \(k\) face probabilities sum to \(\theta^k\le\theta\), so
the same argument applies at every fixed derivative order. The
coefficients are organized by Fa\`a di Bruno's higher-order chain
rule~\cite{Comtet1974}.

We also need a bound at nearby points. If \(\norm v_U\le h\),
every exponential weight at \(x+v\) differs from its value at
\(x\) by a factor between \(e^{-1}\) and \(e\). Comparing the
positive weight products term by term gives the \emph{derivative
envelope}: for \(q\ge1\) and directions \(z_1,\ldots,z_q\in\R^r\),
\[
 \sup_{\norm v_U\le h}|D^qF(x+v)[z_1,\ldots,z_q]|
 \le C_qh^{-q}F(x)\theta(x)\prod_{j=1}^q\norm{z_j}_U.
\]
Here \(C_q\) depends only on \(q\). Appendix~\ref{appx-sec:smoothing}
proves this bound as~\eqref{appx-eq:local-envelope}, retaining the
individual face factors that the covariance argument will use.
The derivative formula is also given in~\cite[Lemma~13]{GaoZhang}.

The product \(F\theta\) is at most \(1/4\), but anti-concentration
gives a smaller bound for its expectation. Fix \(x\), put
\(T\coloneqq\sum_\alpha w_\alpha(x)\), and let \(F_s\) denote
the smoother with every threshold increased by \(s\).
Shifting by \(s h\) replaces \(T\) by \(e^{-s}T\), so
\[
 F_h(x)-F_{-h}(x)
 =\int_{-1}^{1}\frac{T}{(e^{s/2}+Te^{-s/2})^2}\,ds
 \ge\frac2e F(x)\theta(x).
\]
Both exponential factors in the denominator are at most
\(e^{1/2}\) throughout the integration interval, which explains
the constant. By the random-threshold representation, the difference
on the left is an average of indicators of layers
\(A_{a,\omega}^h\setminus A_{a,\omega}^{-h}\).
Lemma~\ref{lem:concentration} therefore gives
\[
 \E[F(U)\theta(U)]\le Ch\sqrt L
\]
whenever \(U\in\R^r\) has an independent component satisfying
that lemma with face variance \(v\ge1/2\) and summand face bound
\(\eta\le h/2\). Lemma~\ref{appx-lem:envelope} states the precise
conditions. For fixed directions, combining this estimate with the
derivative envelope gives
\[
 \E\sup_{\norm v_U\le h}|D^qF(U+v)[z_1,\ldots,z_q]|
 \le C_qh^{1-q}\sqrt L\prod_{j=1}^q\norm{z_j}_U.
\]
Thus averaging supplies a factor \(h\sqrt L\) beyond the uniform
derivative bound, a useful gain when \(h\sqrt L\) is small.

\subsection{The Bernoulli--Gauss interpolation}
\label{sec:outline-path}

To compare the smoothed expectations at \(S\) and \(G\), we connect
their distributions by randomly deleting summands and adding Gaussian
noise. The retained summands are multiplied by a deterministic factor,
called the amplitude, which helps preserve the covariance along the
path. Let \(t\in[0,1]\) denote interpolation time, and set
\[
 \gamma\coloneqq 1+t,\qquad \varpi\coloneqq \frac{1-t}{(1+t)^2}.
\]
For each fixed \(t\), let \(\mathbb B_{1,t},\ldots,\mathbb B_{n,t}\)
be independent Bernoulli variables with success probability
\(\varpi\), independent of the summands and of \(G\). Define
\begin{equation}
 U_t\coloneqq\gamma\sum_{i=1}^n\mathbb B_{i,t}\zeta_i^0+\sqrt t\,G.
 \label{eq:path}
\end{equation}
Here \(\gamma\) is the amplitude of each retained summand, and
\(\varpi\) is the probability of keeping a summand. These Bernoulli
variables are called \emph{masks}: \(\mathbb B_{i,t}=1\) retains
summand \(i\), whereas \(\mathbb B_{i,t}=0\) deletes it. This random
selection is called Bernoulli thinning. At time zero every summand is retained, so
\(U_0=S\). At time one every summand is deleted, so \(U_1=G\).
This is the Bernoulli--Gauss interpolation of~\cite[Section~4]{AL}.

Bernoulli thinning is also the mechanism underlying dropout
regularization in statistical learning, where randomly selected
activations or features are set to zero~\cite{Dropout2014}. In the
usual mean-preserving convention, retained activations are divided by
their retention probability, so their expectation over the random
masks is unchanged. Here the summands are already centered, and their
rescaling, together with the added Gaussian noise, preserves the total
covariance \(I_r\) along the interpolation.

The amplitude and retention probability are chosen to preserve covariance
while controlling how the third moments change with time. They satisfy
\[
 \varpi\gamma^2=1-t,\qquad
 \varpi\gamma^3=1-t^2.
\]
Because the masks are independent of the centered summands, the covariance
of the retained sum is \(\varpi\gamma^2I_r=(1-t)I_r\). The independent
Gaussian term contributes \(tI_r\), so \(\Cov(U_t)=I_r\) throughout
the interpolation. Similarly, multiplying a summand by
\(\gamma\mathbb B_{i,t}\) multiplies each of its third-order coordinate
moments by \(\varpi\gamma^3=1-t^2\). The derivative of this coefficient
is \(-2t\), which vanishes at the bounded endpoint \(t=0\). As the
calculation below shows, this vanishing factor also appears in the term
involving third derivatives of the logistic approximation. This is the
purpose of the calibration: it weakens that term near the endpoint where
little Gaussian noise is available~\cite[eq.~(17)]{AL}.

To track the smoothed probability for a fixed threshold vector
\(a\in\R^m\), define
\[
 P_a(t)\coloneqq\E F_{a,h}(U_t).
\]
Its endpoint values are \(\E F_{a,h}(S)\) and \(\E F_{a,h}(G)\), so
\[
 M_h=\sup_a|P_a(1)-P_a(0)|.
\]
The bounds below justify integrating \(P_a'\) up to both endpoints.
The fundamental theorem of calculus, together with the smoothing lemma,
then gives
\begin{equation}
 \rho\le M_h+Ch\sqrt L\,\kappa
 \le\sup_a\int_0^1|P_a'(t)|\dd t+Ch\sqrt L\,\kappa.
 \label{eq:roadmap}
\end{equation}
The final term accounts for replacing rectangle indicators by their
logistic approximations. The remaining task is therefore to bound the
rate of change \(P_a'(t)\), uniformly over threshold vectors, and integrate
that bound over time.

The derivative separates into a third-moment term and a fourth-order
remainder. Lemma~\ref{appx-an:lem:identity} gives, for \(0<t<1\),
\begin{equation}
 P_a'(t)=-\frac t3
       \E D^3F_{a,h}(U_t):\sum_i\E(\zeta_i^0)^{\otimes3}
       +R_{4,a,h}(t).
 \label{eq:path-summary}
\end{equation}
The tensor \((\zeta_i^0)^{\otimes3}\) has entries
\(\zeta_{ij}^0\zeta_{ik}^0\zeta_{i\ell}^0\), for \(j,k,\ell\le r\).
The colon denotes the sum of products of corresponding tensor entries. The coefficient \(t\) makes
the third-moment contribution small near the bounded endpoint.
Moreover, every third derivative is evaluated at the same point
\(U_t\), allowing the summands' third moments to cancel before an
absolute value is taken. Under condition~\eqref{eq:zero-third}, therefore,
\begin{equation}
 P_a'(t)=R_{4,a,h}(t).
 \label{zs:identity}
\end{equation}
The fourth-order remainder has a similar tensor representation, with
derivatives evaluated along short segments near \(U_t\). To describe
these segments, first remove the \(i\)th bounded summand from the
interpolation path and set
\[
 H_i\coloneqq U_t-\gamma\mathbb B_{i,t}\zeta_i^0,
 \qquad \gamma=1+t.
\]
For each \(i\), introduce independent copies \(\zeta_i\) and
\(\widetilde\zeta_i\) of \(\zeta_i^0\). Take all copies to be
mutually independent and independent of the original summands,
the Bernoulli masks, and the Gaussian vector. The points
\(H_i+s\gamma\zeta_i\), \(0\le s\le1\), then trace a segment that
inserts an independent copy of the deleted summand with increasing
amplitude. The fourth-order tensor associated with this segment is
\[
 \begin{aligned}
 \mathcal T_{i,t}(s)\coloneqq{}&a_t(s)(\zeta_i)^{\otimes4}\\
 &+b_t(s)(\widetilde\zeta_i)^{\otimes2}\otimes(\zeta_i)^{\otimes2}\\
 &+c_t(s)(\widetilde\zeta_i)^{\otimes3}\otimes\zeta_i.
 \end{aligned}
\]
The scalar kernels \(a_t,b_t,c_t\) are specified in
Lemma~\ref{appx-an:lem:identity}. Their absolute integrals over
\(0\le s\le1\) are bounded by an absolute constant, uniformly in
\(0<t<1\).
With the colon now contracting all four tensor indices, the exact
remainder is
\[
 R_{4,a,h}(t)=\int_0^1\sum_i\E\left[
 D^4F_{a,h}(H_i+s\gamma\zeta_i):\mathcal T_{i,t}(s)
 \right]\dd s.
\]
Thus the first term in~\eqref{eq:path-summary} contracts a third
derivative with the aggregate third-moment tensor, while the remainder
contracts fourth derivatives with weighted products of four summand
components. Every tensor product in \(\mathcal T_{i,t}(s)\) contains
two factors from the same independent copy. This is the pair that the
covariance inequality in Subsection~\ref{sec:outline-early} will control.
The fourth derivatives are evaluated at points depending on \(i\),
\(s\), and \(\zeta_i\), whose distance from the common interpolation
point satisfies
\[
 \norm{H_i+s\gamma\zeta_i-U_t}_U
 =\norm{s\gamma\zeta_i-\gamma\mathbb B_{i,t}\zeta_i^0}_U
 \le4\delta.
\]
These shifts explain the need for the local derivative envelope.
Lemma~\ref{appx-an:lem:identity} proves the identity. The following
proposition bounds both the full derivative and the remainder.

\begin{proposition}[Two time ranges]\label{an:prop:derivative}
There are sufficiently large absolute constants \(\Ac,d_0\) such that,
for \(h\ge\Ac\delta\), uniformly over the thresholds,
\begin{align}
 |P_a'(t)|&\le C\sqrt L\left(
 \frac{t\delta}{h^2}+\frac{\delta^2}{h^3}\right),
 &&0<t<1,\label{an:early}\\
 |P_a'(t)|&\le CL^{3/2}\left(
 \delta+\frac{\delta^2}{t h}\right),
 &&d_0\delta^2L\le t<1.\label{an:late}
\end{align}
The fourth-order remainder separately satisfies
\begin{align}
 |R_{4,a,h}(t)|&\le C\delta^2\sqrt L/h^3,
 &&0<t<1,\label{eq:remainder-early}\\
 |R_{4,a,h}(t)|&\le C\delta^2L^{3/2}/(t h),
 &&d_0\delta^2L\le t<1.\label{an:remainder}
\end{align}
\end{proposition}

For fixed \(h>0\), the first bound is integrable over \(0<t<1\).
Moreover, boundedness and continuity of the logistic approximation give
\(P_a(t)\to P_a(0)\) as \(t\downarrow0\) and
\(P_a(t)\to P_a(1)\) as \(t\uparrow1\). These facts justify integration
up to both endpoints, as used in~\eqref{eq:roadmap}. The next two
subsections explain the early and late estimates.

\subsection{The early bounds and the covariance identity}
\label{sec:outline-early}

We now explain the derivative bounds~\eqref{an:early}
and~\eqref{eq:remainder-early}, which remain finite near the bounded
endpoint \(t=0\). The argument combines the covariance identity,
the summand face bound, and the derivative envelope from
Subsection~\ref{sec:outline-smoothing}. The independent copies
\(\zeta_i,\widetilde\zeta_i\) introduced in
Subsection~\ref{sec:outline-path} each have aggregate covariance
\(I_r\). Consequently, for any
unit vectors \(u,v\in\R^r\),
Cauchy--Schwarz gives
\[
 \begin{aligned}
 \sum_i\E|(u^\top\zeta_i)(v^\top\zeta_i)|
 &\le\Bigl(\sum_i\E(u^\top\zeta_i)^2\Bigr)^{1/2}
       \Bigl(\sum_i\E(v^\top\zeta_i)^2\Bigr)^{1/2}\\
 &=1.
 \end{aligned}
\]
This inequality controls the sum of the products contributed by a pair
of derivative directions. Appendix~\ref{appx-sec:support} records its
conditional version, which will also be used in the late bounds.

Each third derivative produces products of three components of the
summands along face normals; each fourth derivative produces products
of four. The covariance inequality above controls the sum of the
products of two such components. Bounding each remaining component
by \(\delta\) therefore gives, for \(q=3,4\),
\[
 \begin{aligned}
 \sum_i\E\prod_{a=1}^{q}|u_{\alpha_a}^{\top}\zeta_i|
 &\le\delta^{q-2}\sum_i\E
     |(u_{\alpha_1}^{\top}\zeta_i)(u_{\alpha_2}^{\top}\zeta_i)|\\
 &\le\delta^{q-2}.
 \end{aligned}
\]
The same reasoning applies to the fourth-order terms involving two
independent copies: we retain the two factors from the repeated copy
and bound the other two by \(\delta\). Thus the summands contribute a
factor \(\delta\) to the third-order terms and \(\delta^2\) to the
fourth-order terms.

We now use these moment bounds to estimate the third-order term and
the fourth-order remainder in~\eqref{eq:path-summary}. Expanding the
directional derivatives expresses each contribution as a sum of
products of logistic weights and components of the summands along
face normals. The preceding inequalities control the summand
components. To apply them, however, we must separate these components
from the logistic weights: in the remainder, the weights are
evaluated at shifted points that depend on the same copies of the
summands. The local envelope from
Subsection~\ref{sec:outline-smoothing} allows us to bound these weights
by weights evaluated at the common interpolation point, which is
independent of those copies.

The shifted evaluation points differ from the common interpolation
point by at most \(4\delta\) along each face normal. Since
\(h\ge4\delta\), this comparison costs only an absolute constant.
Each derivative contributes a factor \(h^{-1}\), while rectangle
anti-concentration bounds the expected sum of the remaining weights
by \(Ch\sqrt L\).
Consequently, the third- and fourth-order contributions are bounded,
respectively, by
\[
 C\underbrace{\delta}_{\text{moments}}\,
  \underbrace{h^{-3}}_{\text{derivatives}}\,
  \underbrace{h\sqrt L}_{\text{anti-concentration}}
  =\frac{C\delta\sqrt L}{h^2},
 \qquad
 C\delta^2h^{-4}\,h\sqrt L=\frac{C\delta^2\sqrt L}{h^3}.
\]
Including the coefficient \(t/3\) of the third-order term
in~\eqref{eq:path-summary} gives the bounds~\eqref{an:early}
and~\eqref{eq:remainder-early}, proved in
Appendix~\ref{appx-an:sec:early}:
\[
 |P_a'(t)|\le C\sqrt L
   \left(\frac{t\delta}{h^2}+\frac{\delta^2}{h^3}\right),
 \qquad
 |R_{4,a,h}(t)|\le\frac{C\delta^2\sqrt L}{h^3}.
\]
\subsection{The late bounds}
\label{sec:outline-late}

We now explain the late bounds~\eqref{an:late} and~\eqref{an:remainder},
which apply when \(t\ge d_0\delta^2L\). Their purpose is to improve
the dependence on the logistic smoothing width \(h\): the early bounds
contain factors \(h^{-2}\) and \(h^{-3}\) for the third- and fourth-order
contributions, respectively, whereas the late bounds are independent
of \(h\) for the third-order contribution and contain only \(h^{-1}\)
for the fourth-order remainder. To obtain this improvement, two
Gaussian integrations by parts rewrite the third and fourth
derivatives as first and second derivatives multiplied by explicit
quadratic functions of Gaussian variables. The resulting bounds
involve inverse powers of the Gaussian variance, which explains why
we use them away from \(t=0\). The remaining task is to control these
Gaussian factors without introducing dependence on the covariance
rank or losing the gain from anti-concentration.

To use the Gaussian noise for both integration by parts and
anti-concentration, we split it into two independent components.
The bounded component of the path~\eqref{eq:path} is
\[
 V_t\coloneqq\gamma\sum_{i=1}^n\mathbb B_{i,t}\zeta_i^0.
\]
Let \(G_0,G_1\sim N(0,I_r)\) be independent of one another, the
summands, and the masks. We can represent the interpolation point,
without changing its distribution, as
\[
 U_t=V_{0,t}+\sqrt{t/2}\,G_1,
 \qquad V_{0,t}\coloneqq V_t+\sqrt{t/2}\,G_0.
\]
We integrate by parts with respect to \(G_1\), conditioning on
\(V_{0,t}\). Keeping the other Gaussian component inside \(V_{0,t}\)
allows this remaining vector to retain the anti-concentration
property needed below, even as the bounded component disappears
near \(t=1\).

For every unit face normal \(u_\alpha\), the two independent
components of \(V_{0,t}\) satisfy
\[
 \Var(u_\alpha^\top V_t)=1-t,
 \qquad
 \Var\!\left(u_\alpha^\top\sqrt{t/2}\,G_0\right)=t/2.
\]
Since
\[
 \max\{1-t,t/2\}\ge\frac13,
\]
at least one component has variance bounded below by \(1/3\) along
every face normal, uniformly over interpolation time. The bounded
component \(V_t\) is a sum of independent centered vectors whose
components along each face normal are at most \(2\delta\) in absolute
value, because \(\gamma\le2\). Thus either the bounded component or
the Gaussian component satisfies the assumptions of
Lemma~\ref{lem:concentration} with a variance bounded away from zero.
Conditioning on the other component,
Lemma~\ref{lem:concentration} bounds the probability that \(V_{0,t}\)
lies in a boundary layer formed by shifting all face thresholds
outward and inward by \(s>0\) by
\[
 C(s+\delta)\sqrt L,
\]
uniformly in the thresholds and in \(t\). We call a vector with this
property a \emph{concentrated base}; equation~\eqref{appx-an:concentration}
records its precise definition.

We first use integration by parts to reduce the derivative order.
Condition on the bounded summands, the masks, \(G_0\), and the
independent copies, so that only \(G_1\) remains random. Consider
a derivative of order \(q\in\{3,4\}\), with repeated direction
\(\zeta_i\), remaining directions \(z_3,\ldots,z_q\), and evaluation
point \(U_t+v_i\). The shift and all these directions are fixed under
the conditioning. Since \(U_t=V_{0,t}+\sqrt{t/2}\,G_1\), two Gaussian
integrations by parts, using Stein's identity~\cite{Stein1981}, give
\[
 \begin{aligned}
 &\E_{G_1}D^qF(U_t+v_i)
       [\zeta_i,\zeta_i,z_3,\ldots,z_q]\\
 &\qquad=\frac2t\,\E_{G_1}\!\left[
   \bigl((G_1^\top\zeta_i)^2-\norm{\zeta_i}^2\bigr)
   D^{q-2}F(U_t+v_i)[z_3,\ldots,z_q]
   \right],
 \end{aligned}
\]
where \(\E_{G_1}\) denotes expectation over \(G_1\) alone.
The same identity holds when the repeated direction is
\(\widetilde\zeta_i\). Thus two derivatives, each contributing a
factor \(h^{-1}\) in the direct estimate, are replaced by the factor
\(2/t\) and an explicit quadratic Gaussian factor. Only first
derivatives remain when \(q=3\), and only second derivatives when
\(q=4\).

The remaining derivatives contain weights of the form
\(F(x)\prod_{a=1}^k\pi_{\alpha_a}(x)\), with \(k=1\) or \(2\).
We need to express their values at the shifted points \(U_t+v_i\)
in terms of weights at the common point \(U_t\). To see how, recall
that shifting \(x\) to \(x+v\) multiplies the exponential associated
with face \(\beta\) by \(e^{u_\beta^\top v/h}\). Define
\[
 d_\beta(v)\coloneqq e^{u_\beta^\top v/h}-1.
\]
Comparing the numerators and denominators of the logistic weights
gives the exact identity
\[
 \begin{aligned}
 F(x+v)\prod_{a=1}^k\pi_{\alpha_a}(x+v)
 ={}&F(x)\prod_{a=1}^k\pi_{\alpha_a}(x)\\
 &\times\frac{
   \exp\!\left(\sum_{a=1}^k u_{\alpha_a}^\top v/h\right)}{
   \left[1+\sum_\beta\pi_\beta(x)d_\beta(v)\right]^{k+1}}.
 \end{aligned}
\]
Since \(|u_\beta^\top v|\le4\delta\) and
\(\sum_\beta\pi_\beta(x)\le1\),
\[
 \left|\sum_\beta\pi_\beta(x)d_\beta(v)\right|
 \le e^{4\delta/h}-1<1
\]
when \(h\) is a sufficiently large multiple of \(\delta\). We can
therefore expand the denominator factor in an absolutely convergent
power series. Each term is a product of logistic weights evaluated
at \(x\), multiplied by a coefficient depending on the shift \(v\).
Taking \(x=U_t\) and \(v=v_i\) gives common weights for all
observations, with the dependence on the independent copies confined
to the coefficients. Because the shifts are fixed under the
conditioning, these coefficients do not depend on \(G_1\). This exact
expansion allows the signed Gaussian factors to remain inside the
expectation while we sum the coefficient matrices.
Equation~\eqref{eq:weight-translation} and the proof of
Lemma~\ref{appx-an:lem:restoration} give the details.

After expanding the shifted weights, we collect the summands into
coefficient matrices. A typical matrix has the form
\[
 \Xi\coloneqq\sum_i\E[c_i\zeta_i\zeta_i^\top\mid V_{0,t}],
\]
where the scalar coefficients \(c_i\) contain the remaining summand
components and the factors arising from the shifts. If \(|c_i|\le K\),
the conditional covariance identity gives
\[
 |u^\top\Xi u|
 \le K\sum_i\E[(u^\top\zeta_i)^2\mid V_{0,t}]
 =K
 \qquad\text{for every unit vector }u.
\]
Thus \(\norm\Xi_\op\le K\), where the operator norm of a symmetric
matrix is its largest absolute eigenvalue. The same argument applies
to \(\widetilde\zeta_i\). Operator--nuclear norm
duality~\cite[Chapter~IV]{Bhatia1997} gives, for any symmetric matrix \(M\),
\[
 |\tr(\Xi M)|\le\norm\Xi_\op\norm M_*,
\]
where \(\norm M_*\), the nuclear norm, is the sum of its absolute
eigenvalues. This inequality separates the contribution of the
summands from the Gaussian matrix average, while keeping the
coefficient bound independent of the rank \(r\).

We therefore need to control Gaussian matrix averages weighted by
the common logistic factors. For face indices
\(\alpha_1,\ldots,\alpha_k\), define
\[
 M_{\alpha_1,\ldots,\alpha_k,t}\coloneqq
 \E_{G_1}\!\left[
 (G_1G_1^\top-I_r)F(U_t)\prod_{a=1}^k\pi_{\alpha_a}(U_t)
 \right].
\]
This matrix depends on the base \(V_{0,t}\), which is held fixed in
the inner expectation. Lemma~\ref{appx-an:lem:matrix} proves that
\[
 \E_{V_{0,t}}\sum_{\alpha_1,\ldots,\alpha_k}
 \norm{M_{\alpha_1,\ldots,\alpha_k,t}}_*
 \le C^{k+1}hL^{3/2},\qquad k\ge1,
\]
provided \(\sqrt{t/2}\ge\delta\sqrt L\). The late-time condition
\(t\ge d_0\delta^2L\), with \(d_0\ge2\), ensures this requirement.
Although the original first and second derivatives contain only
one or two face probabilities, the expansion of the shifted weights
produces longer products; this is why the estimate is needed for
every \(k\). Taking \(h\) sufficiently large relative to \(\delta\)
makes that expansion summable even after the factors \(C^{k+1}\)
are included.

The factor \(h\) in this matrix estimate is essential to the improved
dependence on the smoothing width. To understand its origin, fix the
base and a face tuple, multiply the standard Gaussian density by the
corresponding logistic weight, and normalize the result to integrate
to one. The weight is log-concave, so the covariance matrix of this
reweighted Gaussian distribution is at most \(I_r\), by
Pr\'ekopa's marginal log-concavity theorem~\cite[Theorem~6]{Prekopa1973},
as shown in Lemma~\ref{appx-an:lem:matrix}. This fact bounds
the nuclear norm by two scalar contributions: one involving the
squared length of the new mean, and one involving a signed trace.
Lemma~\ref{appx-an:lem:matrix} bounds the mean contribution using
anti-concentration and Gaussian projections onto the face normals.
Summed over the face tuples and averaged over the base, its bound
is \(C(k+1)hL^{3/2}\).

For the trace contribution, we first sum over the face tuples while
retaining its sign. Since the face probabilities sum to
\(\theta=1-F\), the resulting expression is
\[
 \E\!\left[(\norm{G_1}^2-r)F(U_t)\theta(U_t)^k\right].
\]
The proof expresses the logistic weights as averages of indicators
of regions defined by shifted face thresholds. Differentiating a
common threshold shift then turns this expectation into Gaussian
boundary integrals, with a factor proportional to \(h/\sqrt{t/2}\).
Lemma~\ref{appx-an:lem:surface} controls these boundary integrals
after averaging over the concentrated base \(V_{0,t}\). Their bound
has order
\[
 \sqrt{t/2}\,L^{3/2}+\delta L^2
 \le2\sqrt{t/2}\,L^{3/2},
 \qquad\sqrt{t/2}\ge\delta\sqrt L.
\]
The factor \(\sqrt{t/2}\) therefore cancels the denominator from
the threshold differentiation, leaving the required factor
\(hL^{3/2}\). Averaging over the base is the step that supplies this
gain through anti-concentration.
Appendix~\ref{appx-an:sec:surface} gives the geometric argument.

We can now combine the estimates. For a derivative of order
\(q\in\{3,4\}\), integration by parts contributes \(2/t\). The
remaining \(q-2\) summand components contribute \(\delta^{q-2}\),
and the remaining derivatives contribute \(h^{-(q-2)}\). The
Gaussian matrix estimate supplies \(hL^{3/2}\). After summing the
convergent expansion, the resulting bound is therefore of order
\[
 \underbrace{\frac2t}_{\text{integration by parts}}\,
 \underbrace{\delta^{q-2}}_{\text{summand components}}\,
 \underbrace{h^{-(q-2)}}_{\text{remaining derivatives}}\,
 \underbrace{hL^{3/2}}_{\text{Gaussian matrix estimate}}
 =\frac{2\delta^{q-2}h^{3-q}L^{3/2}}{t}.
\]
For the third-order contribution, \(q=3\), and the coefficient
\(t/3\) in~\eqref{eq:path-summary} cancels the factor \(1/t\),
giving \(C\delta L^{3/2}\). For the fourth-order remainder, \(q=4\),
and integrating the uniformly bounded scalar kernels gives
\(C\delta^2L^{3/2}/(th)\). Consequently,
\[
 |R_{4,a,h}(t)|\le\frac{C\delta^2L^{3/2}}{th},
 \qquad
 |P_a'(t)|\le CL^{3/2}\left(\delta+\frac{\delta^2}{th}\right).
\]
These are the late bounds~\eqref{an:remainder} and~\eqref{an:late}.


% Source: rate_explanation
\subsection{\texorpdfstring{Integrating the comparison and choosing the scale}{Integrating the comparison and choosing the scale}}\label{sec:rate}
We now combine the early and late bounds to control the change in the
smoothed expectation along the entire interpolation path. We then
choose the logistic smoothing width \(h\) to balance this comparison
error against the error from replacing rectangle indicators by their
logistic approximations. Recall that
\[
 M_h=\sup_a|P_a(1)-P_a(0)|.
\]
Thus integration along the path gives
\[
 M_h\le\sup_a\int_0^1|P_a'(t)|\dd t.
\]

The early and late estimates for the fourth-order remainder have
respective orders
\[
 \frac{\delta^2\sqrt L}{h^3}
 \qquad\text{and}\qquad
 \frac{\delta^2L^{3/2}}{th}.
\]
These expressions agree at the interpolation time
\[
 \tau_h\coloneqq Lh^2.
\]
We therefore use the early estimate before \(\tau_h\) and the late
estimate afterwards. This choice is permitted when
\(h\ge\mathsf A\delta\), with \(\mathsf A\) sufficiently large:
taking \(\mathsf A^2\ge d_0\) ensures that
\(\tau_h\ge d_0\delta^2L\), so the late estimate is available
throughout the second time interval. If \(\tau_h\ge1\), we use the
early estimate along the whole path.

Integrating these estimates yields the following comparison bound.
It is uniform over all finite threshold vectors \(a\in\R^m\),
with the same \(m=2p\) face normals, including when the associated
rectangle is empty. This scope also covers the threshold shifts
used in the smoothing argument.

\begin{proposition}[Smoothed comparison]\label{prop:moving-cutoff}
\label{zs:smooth-proposition}
There are absolute constants \(\mathsf A,C\) such that, under
\eqref{eq:setup}, for every \(h\ge\mathsf A\delta\),
\begin{equation}
 M_h\le CL^{3/2}\left[
 \delta+\frac{\delta^2}{h}\log_+\frac1{Lh^2}\right].
 \tag{S*}\label{eq:smooth-moving}
\end{equation}
If, in addition, the aggregate third-moment tensor vanishes,
as in~\eqref{eq:zero-third}, then
\begin{equation}
 M_h\le\frac{C\delta^2L^{3/2}}h
 \left[1+\log_+\frac1{Lh^2}\right].
 \label{zs:smooth}
\end{equation}
Both estimates hold uniformly for all \(a\in\R^m\).
\end{proposition}

The two terms inside the brackets in~\eqref{zs:smooth} come from the
two time intervals. When \(\tau_h<1\), integrating the fourth-order
estimates gives
\[
 \begin{aligned}
 \int_0^1|R_{4,a,h}(t)|\dd t
 &\le\frac{C\delta^2\sqrt L}{h^3}\int_0^{\tau_h}\dd t
    +\frac{C\delta^2L^{3/2}}h\int_{\tau_h}^1\frac{\dd t}{t}\\
 &=\frac{C\delta^2L^{3/2}}h
    \left[1+\log\frac1{Lh^2}\right].
 \end{aligned}
\]
The additive one records the early contribution, while the logarithm
comes from integrating \(1/t\) in the late interval. When
\(\tau_h\ge1\), integration of the early estimate alone gives the
same upper bound with the logarithmic term set to zero.

The third-order contribution, when present, has integrated magnitude
at most \(C\delta L^{3/2}\). Indeed, before \(\tau_h\), its early
bound is \(Ct\delta\sqrt L/h^2\le C\delta L^{3/2}\); afterwards,
its late bound is \(C\delta L^{3/2}\). Since \(h\ge\delta\), the
early fourth-order contribution is also at most
\(C\delta L^{3/2}\) and can be absorbed into this term, giving
\eqref{eq:smooth-moving}. When the aggregate third moments vanish,
the third-order term in~\eqref{eq:path-summary} is identically zero.
We then retain the early fourth-order contribution explicitly, which
gives the additive one in~\eqref{zs:smooth}.

It remains to choose \(h\). A larger logistic smoothing width reduces
the fourth-order comparison error but increases the error from
approximating the rectangle indicator. Combining
\eqref{eq:smooth-moving} with the smoothing-error bound in
Lemma~\ref{lem:endpoint} gives
\[
 \rho\le C\left[
 \delta L^{3/2}
 +\frac{\delta^2L^{3/2}}h\log_+\frac1{Lh^2}
 +h\sqrt L\,\kappa\right].
\]
The first term is independent of \(h\); the choice of width balances
the last two. Recall that
\[
 \kappa=\min\{L,c^{-1}\},
 \qquad E_\kappa=1+\log_+\frac{\sqrt\kappa}{\delta L},
\]
where \(c\) is the correlation gap and \(\kappa=L\) when \(c=0\).
A suitable width is
\begin{equation}
 h\coloneqq\mathsf A\delta\sqrt{\frac{LE_\kappa}{\kappa}}.
 \label{eq:optimal-scale}
\end{equation}
The factor \(E_\kappa\) accounts for the logarithm produced by time
integration.

This width satisfies \(h\ge\mathsf A\delta\), because
\(\kappa\le L\) and \(E_\kappa\ge1\). Moreover, taking
\(\mathsf A\ge1\) gives
\[
 1+\log_+\frac1{Lh^2}
 =1+2\log_+\frac{\sqrt\kappa}
     {\mathsf A\delta L\sqrt{E_\kappa}}
 \le2E_\kappa.
\]
Consequently, the smoothing error and the integrated fourth-order
error are both bounded by a constant times
\[
 \delta L\sqrt{\kappa E_\kappa}.
\]
Adding the third-order contribution proves the bound in
Theorem~\ref{thm:optimized},
\[
 \rho\le C\min\!\left\{1,\,
 \delta L^{3/2}+\delta L\sqrt{\kappa E_\kappa}\right\}.
\]
Under aggregate third-moment cancellation, that contribution
disappears, and the same choice of \(h\), now applied
to~\eqref{zs:smooth}, proves Theorem~\ref{thm:zero-skew}:
\[
 \rho\le C\min\!\left\{1,\,
 \delta L\sqrt{\kappa E_\kappa}\right\}.
\]
The minimum with one uses the elementary bound \(\rho\le1\).
Finally, \(\kappa\le L\) and \(E_\kappa\le D_L\), where
\(D_L=1+\log_+\{1/(\delta\sqrt L)\}\), yield
Theorem~\ref{thm:unrestricted}. Appendix~\ref{sec:main-proof}
gives the complete proofs.


\subsection{\texorpdfstring{Sources of the argument and the power \(n^{-1/2}\)}{Sources of the argument and the power n to the -1/2}}
\label{sec:why-rate}\label{sec:sources}

The argument combines several existing constructions. We record
their sources and explain how the dependence of the fourth-order
remainder on the logistic smoothing width \(h\) determines the power
\(n^{-1/2}\). The logistic approximation uses the baseline probability
in a multinomial logit model; its formula appears
in~\cite[p.~1001]{McFadden1989} and~\cite[equation~(5)]{GaoZhang}.
The derivative formulas and the representation through randomly
shifted thresholds correspond to~\cite[Lemmas~13 and~14]{GaoZhang}.

The covariance-preserving Bernoulli--Gaussian interpolation follows
\cite[Section~4]{AL}. Its amplitude and retention probability are
calibrated both to preserve covariance and to make the third-order
coefficient vanish at \(t=0\). The bounded-sum anti-concentration
estimate in Lemma~\ref{lem:concentration} is also the estimate
of~\cite[Lemma~2.1]{AL}. In the exact fourth-order identity of
Lemma~\ref{appx-an:lem:identity}, the scalar kernel \(a_t(s)\)
agrees with the kernel in~\cite[Lemma~4.1]{AL}. Evaluating all
third derivatives at the common interpolation point \(U_t\) gives
the same structure as the derivative identity
in~\cite[Lemma~22]{GaoZhang}. This common evaluation point allows
the third-moment tensors to be summed before their magnitude is
estimated, so aggregate third-moment cancellation remains available.

The transformation to the \(\R^r\) coordinates of
Subsection~\ref{sec:outline-smoothing} follows~\cite[Section~2.3]{GaoZhang}.
The transformed summands have aggregate covariance \(I_r\),
and their components along the transformed face normals retain the
original bound \(\delta\). The late argument uses the covariance
summation and two Gaussian integrations by parts
of~\cite[Lemma~23]{GaoZhang}. The exact translation series in
Subsection~\ref{sec:outline-late} expresses weights at the shifted
evaluation points as a convergent series of weights at \(U_t\),
without introducing an additional Taylor remainder.

The matrix and surface estimates use classical log-concavity and
convex geometry. Gaussian covariance contraction follows from
Pr\'ekopa's theorem~\cite[Theorem~6]{Prekopa1973}; for smooth positive
weights it also follows from the Brascamp--Lieb variance
inequality~\cite{BrascampLieb1976}. The facet--ridge identity is a
Gaussian specialization of the Gauss--Green theorem
\cite[Chapter~5]{EvansGariepy2015}. The outward facet tubes follow
Nazarov~\cite{Nazarov2003}, with a detailed proof
in~\cite[Lemmas~1--2]{CCK17N}. We derive the variable tube widths,
paired-ridge estimates, and averaging needed here in
Subsection~\ref{appx-an:sec:surface}. For related convex-cell and
polyhedral applications, see~\cite[Lemma~12 and Proposition~4]{GaoZhang}.

The improvement in the sample-size dependence comes from the late
fourth-order estimate. Expressed in the notation of this paper, the
first-stage remainder bound in~\cite[Proposition~2]{GaoZhang} has order
\[
 \frac{\delta^2L}{th^2},
\]
whereas the late remainder bound in
Subsection~\ref{sec:outline-late} has order
\[
 \frac{\delta^2L^{3/2}}{th}.
\]
Integration in \(t\) produces a logarithmic factor in both cases,
but leaves different powers of \(h\). Suppressing logarithmic
factors, balancing the integrated remainder against the smoothing
error therefore gives
\[
 h\asymp\frac{\delta^2}{h^2}
 \quad\Longrightarrow\quad h\asymp\delta^{2/3}
\]
for the first estimate, and
\[
 h\asymp\frac{\delta^2}{h}
 \quad\Longrightarrow\quad h\asymp\delta
\]
for the second. With \(\delta=B/\sqrt n\) and fixed \(B\), these
balances give the powers \(n^{-1/3}\) and \(n^{-1/2}\), respectively.
The integrated third-order contribution in this paper is also of
order \(\delta\), up to logarithmic factors, and therefore does
not change the latter power.

The ratio of the two displayed remainder scales is \(h\sqrt L\).
The source of this factor is anti-concentration: the logistic
derivative weights have expected total mass at most \(Ch\sqrt L\),
because their mixture representation concentrates them near randomly
shifted rectangle boundaries. This bound enters the early estimates
directly. For the late estimates, it must survive the two Gaussian
integrations by parts. The matrix estimate retains this small mass
when controlling the contribution of the reweighted Gaussian means.
For the remaining signed trace, differentiation of the common
threshold shift introduces \(h/\sqrt{t/2}\), while the averaged
surface estimate supplies the compensating factor \(\sqrt{t/2}\).
Their product leaves the required factor \(h\). Averaging over the
concentrated base \(V_{0,t}\) is the step that uses anti-concentration
to obtain this gain.

The bound in~\cite{AL} also attains the power \(n^{-1/2}\), but its
Gaussian derivative estimates involve negative powers of the
smallest eigenvalue of \(\Sigma\). The estimates used here depend
instead on aggregate covariance and bounded components along the
face normals. They therefore retain the \(n^{-1/2}\) power without
a positive lower bound on that eigenvalue, including when
\(\Sigma\) is singular.


% Source: sharpness_and_discussion
\section{Sharpness and remaining questions}\label{opt:section}
This section compares the upper bounds with examples that limit how
accurate a Gaussian approximation can be. Two questions are distinct:
what is the smallest possible uniform error bound, and how quickly may
the dimension grow while that error still tends to zero? At a fixed
positive correlation gap, the results below settle the second question
both with and without cancellation of third moments. The sharp error
rate under cancellation remains open. Allowing the correlation gap to
shrink introduces a further obstruction, illustrated by a centrally
symmetric example.

Even in one dimension, the dependence on sample size cannot improve
beyond \(n^{-1/2}\) under the present assumptions. A normalized sum of
independent Rademacher signs has an atom of order \(n^{-1/2}\) at one
of its lattice points. An interval shrinking to that point retains the
atom, whereas its Gaussian probability tends to zero. Thus the
Kolmogorov distance over intervals is at least of order \(n^{-1/2}\).
This obstruction persists under central symmetry and hence under the
third-moment cancellation condition~\eqref{eq:zero-third}.

Growing dimension creates another obstruction through the tails of
coordinate maxima. Chernozhukov, Chetverikov, and
Koike~\cite[Proposition~2.1]{CCK} used independent two-point coordinates
to show that the factor \((\log p)^{3/2}\) is necessary in general.
The following theorem retains that construction and removes its
additional growth restrictions on the envelope and dimension. The
condition \(r_n\to0\) places the comparison in the regime where a
vanishing approximation error is possible.
\begin{samepage}
\begin{theorem}[Lower scale]\label{opt:skew-lower}
Suppose \(p_n\to\infty\), \(B_n\ge2\), and
\(r_n\coloneqq B_n(\log p_n)^{3/2}/\sqrt n\to0\). There are i.i.d.
centered observations with independent coordinates, covariance
\(I_{p_n}\), and envelope \(B_n\), for which
\begin{equation}
 \liminf_n\frac{\rh(W_n,Z_n)}{r_n}\ge\frac{\sqrt2}{4e}.
 \label{opt:lower-constant}
\end{equation}
Here \(W_n\coloneqq n^{-1/2}\sum_iX_{n,i}\),
\(Z_n\sim N(0,I_{p_n})\), and the envelope is attained.
\end{theorem}
\end{samepage}
Each scalar coordinate of the observations has the centered,
unit-variance two-point law
\[
 \Pp(X=B_n)=\frac1{B_n^2+1},\qquad
 \Pp(X=-B_n^{-1})=\frac{B_n^2}{B_n^2+1}.
\]
Its third moment is \(B_n-B_n^{-1}\), which is of order \(B_n\)
for \(B_n\ge2\). Write \(\Phi\) for the standard normal distribution
function. Near the threshold \(x_n\) defined by
\(p_n\{1-\Phi(x_n)\}=1\), we have
\(x_n\asymp\sqrt{\log p_n}\); this is where a Gaussian maximum
has a nontrivial distribution. The third-moment correction changes
the relative tail probability by order \(r_n\). Since each tail is
of order \(1/p_n\), multiplying the \(p_n\) coordinate distribution
functions produces a discrepancy of order \(r_n\) between
the distribution functions of the two maxima. The Supplement~\cite[Section~S.4.1, Theorem~S.3 and
Lemma~S.2]{Supplement} proves the tail estimate uniformly in \(B_n\),
including the lattice corrections needed for this conclusion.

The covariance in this example is the identity, so its correlation
gap is one. It therefore supplies a lower bound for every class with
a prescribed gap \(c_0\in(0,1]\). To state the uniform comparison,
let \(M_n(c_0;B,p)\) be the supremum of \(\rh(W_n,Z_n)\) over
independent centered observations \(X_1,\ldots,X_n\in\R^p\) with
\(\norm{X_i}_\infty\le B\),
\(n^{-1}\sum_i\E X_{ij}^2=1\) for every coordinate, and correlation
gap at least \(c_0\). Here \(W_n=n^{-1/2}\sum_iX_i\), and \(Z_n\)
is centered Gaussian with the same covariance. The observations in
this supremum need not be identically distributed. Let
\(M_n^0(c_0;B,p)\) denote the same supremum restricted to observations
satisfying \(\sum_i\E X_i^{\otimes3}=0\).

The upper and lower rates match when the dimension is large enough
to absorb the additional logarithm in the upper bound. More precisely,
fix \(\Lambda>0\) and suppose
\[
 1+\log_+(\sqrt n/B_n)\le\Lambda\log(ep_n).
\]
If also \(p_n\to\infty\), \(B_n\ge2\), and \(r_n\to0\),
Theorems~\ref{thm:optimized} and~\ref{opt:skew-lower} give
\begin{equation}
 M_n(c_0;B_n,p_n)\asymp_{c_0,\Lambda}
 \frac{B_n(\log p_n)^{3/2}}{\sqrt n}.
 \label{opt:matching-minimax}
\end{equation}
The symbol \(\asymp_{c_0,\Lambda}\) means that the ratio is bounded
above and below by positive constants depending only on \(c_0\) and
\(\Lambda\). In particular, this conclusion holds when
\(p_n\asymp n^a\) for a fixed \(a>0\), with \(B_n\) allowed to grow
subject to \(r_n\to0\).

Uniform convergence requires no additional condition relating the
two logarithms. The next theorem gives its exact criterion and shows
how third-moment cancellation enlarges the admissible dimension range.
The general consistency argument is proved in the
Supplement~\cite[Section~S.4.2, Theorem~S.4]{Supplement}; the
zero-third-moment case is proved there in Section~S.3.
\begin{theorem}[Uniform consistency at a fixed gap]
\label{opt:consistency}
For fixed \(c_0\in(0,1]\), arbitrary dimension sequences \(p_n\ge1\),
and envelopes \(B_n\ge2\),
\begin{equation}
 M_n(c_0;B_n,p_n)\longrightarrow0
 \quad\Longleftrightarrow\quad
 \frac{B_n\{\log(ep_n)\}^{3/2}}{\sqrt n}\longrightarrow0.
 \label{opt:consistency-iff}
\end{equation}
For each fixed \(B\ge1\),
\[
 M_n^0(c_0;B,p_n)\longrightarrow0
 \quad\Longleftrightarrow\quad \log p_n=o(\sqrt n).
\]
\end{theorem}
For a fixed envelope \(B\ge2\), the first criterion is
\(\log p_n=o(n^{1/3})\). Cancellation improves it to
\(\log p_n=o(\sqrt n)\), but this larger range does not by itself
identify the sharp error rate within the range.

Independent Rademacher coordinates explain the second threshold.
If \(\log p_n\sim a\sqrt n\) for a fixed \(a>0\), their normalized
sums satisfy
\[
 \Pp\!\left(\max_{j\le p_n}W_{n,j}\le x_n\right)
 -\Pp\!\left(\max_{j\le p_n}Z_{n,j}\le x_n\right)
 \longrightarrow e^{-\exp(-a^2/3)}-e^{-1}>0,
\]
where \(x_n\) again satisfies \(p_n\{1-\Phi(x_n)\}=1\).
The fourth-order tail correction therefore prevents convergence at
this dimension scale, even though all third moments vanish.
The calculation and the necessity argument for arbitrary dimension
sequences are in the Supplement~\cite[Section~S.3, Proposition~S.2
and the following proof]{Supplement}.

Below this threshold, a gap between the known upper and lower rates
remains. In the range \(n^{1/3}\ll L\ll\sqrt n\), the fourth-order
tail calculations give a lower-bound scale \(L^2/n\), whereas
Theorem~\ref{thm:zero-skew}, for fixed \(B\) and \(c_0\), gives
\[
 \rh(W_n,Z_n)\le C_{B,c_0}\frac L{\sqrt n}
       \sqrt{1+\log_+\frac{\sqrt n}{L}}.
\]
The Supplement~\cite[Section~S.3.1]{Supplement} discusses this
difference. Any improvement valid also in fixed dimension must retain
the lattice obstruction of order \(n^{-1/2}\) described above;
\(L^2/n\) alone cannot be a uniform upper bound under the present
assumptions.

When the correlation gap is allowed to shrink, central symmetry alone
does not improve the lower scale \(\delta L^{3/2}\). The next theorem
gives a singular example with unit marginal variances and a positive,
but vanishing, correlation gap.

\begin{theorem}[Centrally symmetric lower bound]\label{opt:symmetric-lower}
Let \(m_n\ge2\) be integers such that \(m_n\to\infty\) and
\((\log m_n)^3/n\to0\). Put
\[
 p_n\coloneqq m_n^2,\qquad \delta_n\coloneqq\frac2{\sqrt n},
 \qquad L_n\coloneqq\log(ep_n).
\]
For all sufficiently large \(n\), there are i.i.d.\ centered vectors
\(X_{n,i}\in\R^{p_n}\), \(1\le i\le n\), with centrally symmetric
law, unit marginal variances, and \(\norm{X_{n,i}}_\infty\le2\).
Their covariance \(\Sigma_n\) has rank \(2m_n-1\) and correlation
gap \(c_n=(\log m_n)/(2n)\). The normalized sum
\(W_n\coloneqq n^{-1/2}\sum_{i=1}^nX_{n,i}\) and its Gaussian
counterpart \(Z_n\sim N(0,\Sigma_n)\) satisfy
\[
 \liminf_{n\to\infty}
 \frac{\rh(W_n,Z_n)}{\delta_nL_n^{3/2}}\ge\frac1{8e}.
\]
\end{theorem}

Appendix~\ref{app:symmetric-lower} proves the theorem by modifying
the common Bernoulli activation example in~\cite[Proposition~E.1]{FKLZ}.
A single lower orthant attains the stated asymptotic constant.
Since \(c_n\to0\), the result leaves the sharp rate under a fixed
positive gap open.

For unrestricted covariance, these examples leave at most the factor
\(\sqrt{D_L}\) between the general upper bound and the attained lower
scale \(\delta L^{3/2}\), even under central symmetry. Recall that
\(D_L=1+\log_+(1/(\delta\sqrt L))\). Whether that factor is necessary,
and how the optimal rate depends on a shrinking correlation gap,
remain open. The note~\cite{Obstruction} disproves particular
pointwise improvements of the fourth-order remainder estimate for a
two-scale smoother, even after the full covariance contraction.
Those counterexamples rule out the proposed pointwise estimates;
they do not rule out an improved CLT rate. The same note formulates
an integrated remainder estimate that would support the desired
improvement, and~\cite{IntegratedRemainder} proves that estimate for
rank-one covariance. The general integrated estimate remains open.

Some singular covariance structures do admit a simpler reduction.
Exactly equal or opposite coordinates can be removed without changing
the Kolmogorov distance over rectangles: their constraints combine
into intervals for the retained coordinates. In rank one, all unit-variance
coordinates are equal up to sign, so this reduction leaves a scalar
comparison. Its worst-case error has order
\(\min\{1,B/\sqrt n\}\), with no dimension logarithm.
The Supplement proves the coordinate reduction~\cite[Section~S.5,
Proposition~S.6]{Supplement} and the matching rank-one bounds
\cite[Section~S.4.3, Proposition~S.4]{Supplement}. The upper-bound
proof for arbitrary rank does not require either reduction.

For sub-exponential observations, the remaining comparison also
depends on the covariance assumptions. The companion
paper~\cite{Local} gives a sharper logarithmic dependence when every
principal covariance submatrix of order at most four has its smallest
eigenvalue bounded below by a positive constant.
Corollary~\ref{cor:subexponential} retains the power \(n^{-1/2}\), up
to logarithms, without that local nondegeneracy assumption.
Appendix~\ref{app:subexponential} compares the assumptions and rates
in detail.

% Source: references
\begin{thebibliography}{99}
\setlength{\itemsep}{3pt}
\phantomsection
\addcontentsline{toc}{section}{References}
\bibitem{AL}
P.~M. Aronow and P. Lopatto (2026).
\emph{Optimal central limit theorem for bounded random variables in high dimensions}.
Manuscript, arXiv:2609.07067v1.
\url{https://arxiv.org/abs/2609.07067v1}.

\bibitem{CCK2013}
V. Chernozhukov, D. Chetverikov, and K. Kato (2013).
\emph{Gaussian approximations and multiplier bootstrap for maxima of
sums of high-dimensional random vectors}.
The Annals of Statistics \textbf{41}(6), 2786--2819.
\href{https://doi.org/10.1214/13-AOS1161}{doi:10.1214/13-AOS1161}.

\bibitem{CCK2017}
V. Chernozhukov, D. Chetverikov, and K. Kato (2017).
\emph{Central limit theorems and bootstrap in high dimensions}.
The Annals of Probability \textbf{45}(4), 2309--2352.
\href{https://doi.org/10.1214/16-AOP1113}{doi:10.1214/16-AOP1113}.

\bibitem{CCKK2022}
V. Chernozhukov, D. Chetverikov, K. Kato, and Y. Koike (2022).
\emph{Improved central limit theorem and bootstrap approximations in
high dimensions}. The Annals of Statistics \textbf{50}(5), 2562--2586.
\href{https://doi.org/10.1214/22-AOS2193}{doi:10.1214/22-AOS2193}.

\bibitem{CCKK2023}
V. Chernozhukov, D. Chetverikov, K. Kato, and Y. Koike (2023).
\emph{High-dimensional data bootstrap}.
Annual Review of Statistics and Its Application \textbf{10}, 427--449.
Working version: arXiv:2205.09691.
\url{https://arxiv.org/abs/2205.09691}.

\bibitem{Local}
V. Chernozhukov, D. Chetverikov, K. Kato, and Y. Koike,
with AI collaborators (2026).
\emph{High-dimensional Gaussian approximation over rectangles under
local nondegeneracy}. Companion working manuscript, 21 October 2026.

\bibitem{Supplement}
V.~Chernozhukov, D.~Chetverikov, K.~Kato, and Y.~Koike,
with AI collaborators (2026).
\emph{Supplementary material: High-dimensional Gaussian approximations
over rectangles without covariance restrictions}.
Working manuscript, 20 October 2026, revision \texttt{bbf1ab99c733}.
\href{https://github.com/VC2015/HDCLTs/blob/bbf1ab99c733a398c181ca8c9f5a243a41f258c6/research-rectangle-clt-singular/Rectangle_CLT_Supplement.pdf}{Supplement}.

\bibitem{CCK}
V. Chernozhukov, D. Chetverikov, and Y. Koike (2023).
\emph{Nearly optimal central limit theorem and bootstrap approximations
in high dimensions}.
The Annals of Applied Probability \textbf{33}(3), 2374--2425.
\href{https://doi.org/10.1214/22-AAP1870}{doi:10.1214/22-AAP1870}.
Working version: arXiv:2012.09513v2, May 12, 2021.
\url{https://arxiv.org/abs/2012.09513v2}.

\bibitem{FKLZ}
X. Fang, Y. Koike, S.-H. Liu, and Y.-K. Zhao (2026).
\emph{High-dimensional central limit theorems by Stein's method in
the degenerate case}. The Annals of Applied Probability, to appear.
arXiv:2305.17365.
\href{https://research.cuhk.edu.hk/en/publications/high-dimensional-central-limit-theorems-by-steins-method-in-the-d/}{Accepted manuscript (2026)}.

\bibitem{GaoZhang}
Z. Gao and W. Zhang (2026).
\emph{Cubic-root Gaussian approximation under unrestricted covariance}.
arXiv:2608.30221v1, 31 August 2026.
\url{https://arxiv.org/abs/2608.30221v1}.

\bibitem{Kane}
D.~M. Kane (2014).
\emph{The average sensitivity of an intersection of half spaces}.
Research in the Mathematical Sciences \textbf{1}, Article~13.
\href{https://doi.org/10.1186/s40687-014-0013-6}{doi:10.1186/s40687-014-0013-6}.
Working version: arXiv:1309.2987v3.
\url{https://arxiv.org/abs/1309.2987v3}.

\bibitem{McFadden1989}
D. McFadden (1989).
\emph{A method of simulated moments for estimation of discrete
response models without numerical integration}.
Econometrica \textbf{57}(5), 995--1026.
\href{https://doi.org/10.2307/1913621}{doi:10.2307/1913621}.

\bibitem{OST}
R. O'Donnell, R.~A. Servedio, and L.-Y. Tan (2019).
Fooling polytopes.
In \emph{Proceedings of the 51st Annual ACM SIGACT Symposium on
Theory of Computing}, 614--625.
\href{https://doi.org/10.1145/3313276.3316321}{doi:10.1145/3313276.3316321}.

\bibitem{Prekopa1973}
A.~Pr\'ekopa (1973).
\emph{On logarithmic concave measures and functions}.
Acta Scientiarum Mathematicarum (Szeged) \textbf{34}, 335--343.
\url{https://rutcor.rutgers.edu/Prekopa/pdf/SCIENT2.pdf}.

\bibitem{IntegratedRemainder}
\emph{Toward an integrated curvature estimate: a proved rank-one case
and an entropy identity for the general problem}.
Companion research note, 24 October 2026.
\href{https://github.com/VC2015/HDCLTs/blob/main/research-rectangle-clt-singular/Rectangle_CLT_Integrated_Remainder.pdf}{\nolinkurl{Rectangle_CLT_Integrated_Remainder.pdf}}.

\bibitem{Obstruction}
\emph{A curvature obstruction for two-scale smoothing}.
Companion research note, 23 October 2026.
\href{https://github.com/VC2015/HDCLTs/blob/main/research-rectangle-clt-singular/Rectangle_CLT_Curvature_Obstruction.pdf}{\nolinkurl{Rectangle_CLT_Curvature_Obstruction.pdf}}.

\bibitem{Dropout2014}
N. Srivastava, G. Hinton, A. Krizhevsky, I. Sutskever, and
R. Salakhutdinov (2014).
\emph{Dropout: A simple way to prevent neural networks from overfitting}.
Journal of Machine Learning Research \textbf{15}(56), 1929--1958.
\url{https://www.jmlr.org/papers/v15/srivastava14a.html}.

\bibitem{Bhatia1997}
R. Bhatia (1997).
\emph{Matrix Analysis}.
Graduate Texts in Mathematics \textbf{169}. Springer, New York.
\href{https://doi.org/10.1007/978-1-4612-0653-8}{doi:10.1007/978-1-4612-0653-8}.

\bibitem{Comtet1974}
L. Comtet (1974).
\emph{Advanced Combinatorics: The Art of Finite and Infinite Expansions}.
Reidel, Dordrecht.
\href{https://doi.org/10.1007/978-94-010-2196-8}{doi:10.1007/978-94-010-2196-8}.

\bibitem{CoverThomas2006}
T. M. Cover and J. A. Thomas (2006).
\emph{Elements of Information Theory}, 2nd edition.
Wiley, Hoboken, NJ.

\bibitem{EvansGariepy2015}
L. C. Evans and R. F. Gariepy (2015).
\emph{Measure Theory and Fine Properties of Functions}, revised edition.
CRC Press, Boca Raton, FL.

\bibitem{Folland1999}
G. B. Folland (1999).
\emph{Real Analysis: Modern Techniques and Their Applications}, 2nd edition.
Wiley, New York.

\bibitem{RockafellarWets1998}
R. T. Rockafellar and R. J.-B. Wets (1998).
\emph{Variational Analysis}.
Grundlehren der mathematischen Wissenschaften \textbf{317}.
Springer, Berlin.
\href{https://doi.org/10.1007/978-3-642-02431-3}{doi:10.1007/978-3-642-02431-3}.

\bibitem{Stein1981}
C. M. Stein (1981).
\emph{Estimation of the mean of a multivariate normal distribution}.
The Annals of Statistics \textbf{9}(6), 1135--1151.
\href{https://doi.org/10.1214/aos/1176345632}{doi:10.1214/aos/1176345632}.

\bibitem{Ziegler1995}
G. M. Ziegler (1995).
\emph{Lectures on Polytopes}.
Graduate Texts in Mathematics \textbf{152}. Springer, New York.
\href{https://doi.org/10.1007/978-1-4613-8431-1}{doi:10.1007/978-1-4613-8431-1}.

\bibitem{BLM2013}
S. Boucheron, G. Lugosi, and P. Massart (2013).
\emph{Concentration Inequalities: A Nonasymptotic Theory of Independence}.
Oxford University Press, Oxford.
\href{https://doi.org/10.1093/acprof:oso/9780199535255.001.0001}{\nolinkurl{doi:10.1093/acprof:oso/9780199535255.001.0001}}.

\bibitem{Berry1941}
A. C. Berry (1941).
\emph{The accuracy of the Gaussian approximation to the sum of independent variates}.
Transactions of the American Mathematical Society \textbf{49}(1), 122--136.
\href{https://doi.org/10.1090/S0002-9947-1941-0003498-3}{doi:10.1090/S0002-9947-1941-0003498-3}.

\bibitem{BrascampLieb1976}
H. J. Brascamp and E. H. Lieb (1976).
\emph{On extensions of the Brunn--Minkowski and Pr\'ekopa--Leindler
theorems, including inequalities for log concave functions, and with
an application to the diffusion equation}.
Journal of Functional Analysis \textbf{22}(4), 366--389.
\href{https://doi.org/10.1016/0022-1236(76)90004-5}{doi:10.1016/0022-1236(76)90004-5}.

\bibitem{Rockafellar1970}
R. T. Rockafellar (1970).
\emph{Convex Analysis}.
Princeton Mathematical Series \textbf{28}.
Princeton University Press, Princeton, NJ.

\bibitem{Nazarov2003}
F. Nazarov (2003).
\emph{On the maximal perimeter of a convex set in \(\R^n\)
with respect to a Gaussian measure}.
In V. D. Milman and G. Schechtman (eds.),
\emph{Geometric Aspects of Functional Analysis (2001--2002)},
Lecture Notes in Mathematics \textbf{1807}, 169--187.
Springer, Berlin.

\bibitem{CCK17N}
V. Chernozhukov, D. Chetverikov, and K. Kato (2017).
\emph{Detailed proof of Nazarov's inequality}.
arXiv:1711.10696.
\url{https://arxiv.org/abs/1711.10696}.

\end{thebibliography}

\clearpage
\appendix
\numberwithin{equation}{section}
\numberwithin{theorem}{section}
\numberwithin{lemma}{section}
\numberwithin{proposition}{section}
\numberwithin{corollary}{section}
\numberwithin{remark}{section}

\section*{How to read the appendices}
\addcontentsline{toc}{section}{How to read the appendices}
The appendices present the proof from its conclusion back to its
ingredients. Appendix~\ref{sec:main-proof} first assembles the main
bounds: it integrates the interpolation estimates proved in
Appendix~\ref{appx-sec:support} and adds the smoothing error proved in
Appendix~\ref{appx-sec:smoothing}. Both of those appendices use the
anti-concentration inequality proved independently in
Appendix~\ref{appx-app:concentration}. Thus a reader who wants to see
how the rate emerges can start with Appendix~\ref{sec:main-proof},
treating these inputs as given; a reader who prefers to establish the
inputs first can read Appendices~\ref{appx-app:concentration},
\ref{appx-sec:smoothing}, \ref{appx-sec:support}, and
\ref{sec:main-proof}, in that order. The final subsection of
Appendix~\ref{sec:main-proof} then transfers the rectangle bounds
to polyhedra by a linear transformation.

Within Appendix~\ref{appx-sec:support}, the covariance and concentration
facts and the interpolation identity come first, followed by the early and late
bounds. The late-bound proof uses a matrix estimate established in
Subsection~\ref{appx-an:sec:matrix}; that estimate in turn uses the
Gaussian surface estimate in Subsection~\ref{appx-an:sec:surface}.
These two supporting arguments are placed last so that their purpose
is visible before their proofs. Appendix~\ref{app:subexponential}
then extends the bounded-summand results by truncation and transfers
the resulting bounds to polyhedra, whereas
Appendix~\ref{app:symmetric-lower} proves
Theorem~\ref{opt:symmetric-lower} independently of the interpolation argument.

\section*{Notation in the appendices}
\addcontentsline{toc}{section}{Notation in the appendices}
The following symbols, most of which Section~\ref{sec:smoothing}
introduces, are used throughout
Appendices~\ref{sec:main-proof}, \ref{appx-sec:smoothing},
and~\ref{appx-sec:support}.
Appendices~\ref{appx-app:concentration}, \ref{app:subexponential},
and~\ref{app:symmetric-lower} use additional local symbols, introduced
where they occur.

\medskip
\begingroup
\small
\renewcommand{\arraystretch}{1.15}
\begin{longtable}{@{}p{0.27\textwidth}p{0.68\textwidth}@{}}
\textbf{Symbol} & \textbf{Meaning}\\
\endhead
\(\delta,\ L,\ c,\ \kappa,\ E_\kappa\) & summand envelope,
\(\log(ep)\), correlation gap, and the parameters
of~\eqref{eq:parameters}\\
\(F_{a,h},\ M_h\) & logistic approximation of \(A_a\) at scale
\(h\), and \(\sup_a|\E F_{a,h}(S)-\E F_{a,h}(G)|\),
see~\eqref{eq:selector} and~\eqref{eq:smoothed-discrepancy}\\
\(P_a(t),\ \tau_h\) & \(\E F_{a,h}(U_t)\) at time \(t\in[0,1]\) on
the path~\eqref{eq:path}, and the time \(Lh^2\) at which the early and
late bounds meet\\
\(S,\ G,\ \rho\) & sum \(\sum_i\zeta_i^0\), standard Gaussian target
in \(\R^r\), and \(\sup_a|\Pp(S\in A_a)-\Pp(G\in A_a)|=\rh(W,Z)\)\\
\(u_\alpha,\ a_\alpha,\ m\) & unit face normal, its threshold, and
the fixed number \(m=2p\) of defining inequalities\\
\(A_a,\ A_a^s\) & the set defined by \(u_\alpha\cdot x\le a_\alpha\)
and the set with every threshold increased by \(s\)\\
\(F,\ \Psi,\ w_\alpha,\ T,\ N\) & logistic approximation \(F_{a,h}\), its potential
\(h\log(1+T)\), the face weights, their sum
\(T=\sum_\alpha w_\alpha\), and the denominator \(N=1+T\) of \(F\)\\
\(\pi_\alpha,\ \theta,\ \bar u,\ \chi\) & normalized weights, total
mass \(1-F\), weighted face normal \(\nabla\Psi\), and curvature
\(h\tr D^2\Psi\)\\
\(\omega_0,\ \omega_\alpha\) & independent unit exponentials in the
mixture representation of \(F\)\\
\(\Gamma,\ r,\ \norm\cdot_U\) & factor \(\Sigma=\Gamma\Gamma^\top\),
rank of \(\Sigma\), and face norm \(\max_{\alpha\le m}|u_\alpha\cdot x|\)\\
\(\zeta_i^0;\ \zeta_i,\ \widetilde\zeta_i\) & summands in support coordinates; two
independent copies of them\\
\(\gamma,\ \varpi,\ \mathbb B_{i,t}\) & amplitude \(1+t\), selection
probability \((1-t)/(1+t)^2\), and the selection variables\\
\(Y_i,\ V_t,\ U_t,\ H_i\) & selected summands
\(\gamma\mathbb B_{i,t}\zeta_i^0\), their sum, the path
\(V_t+\sqrt t\,G\), and the path with summand \(i\) deleted; all in \(\R^r\)\\
\(V_0,\ \sigma,\ G,\ U\) & concentrated base, Gaussian scale, standard
Gaussian vector on \(\R^r\), and \(U=V_0+\sigma G\); the coordinate
standard deviations \(\sigma_j\) of Corollary~\ref{cor:subexponential}
carry a subscript\\
\(\Phi_r,\ \varphi_r,\ H_2,\ \surf\) & standard Gaussian measure and
density on \(\R^r\), \(H_2(g)=\norm g^2-r\), and the Gaussian surface
integral\\
\(\mathsf A,\ \mathsf A_b,\ d_0\) & large absolute constants in
\(h\ge\mathsf A\delta\), \(h\ge\mathsf A_b\delta\), and
\(t\ge d_0\delta^2L\)\\
\(C,\ C_q;\ c_1,\ c_2,\ c_*\) & absolute constants, large and small;
they may change between occurrences, and a subscript records any
dependence. The letter \(c\) alone is the correlation gap.\\
\end{longtable}
\endgroup

% Source: complete_integration_proofs
\section{\texorpdfstring{Proofs of Proposition~\ref{prop:moving-cutoff} and the theorems}{Proofs of Proposition 2 and the theorems}}\label{sec:main-proof}
We obtain the probability bounds in two steps: integrate the derivative
estimates along the interpolation path, then choose the smoothing width
to balance the comparison and endpoint errors. The inputs are
Proposition~\ref{an:prop:derivative} and Lemma~\ref{lem:endpoint}, proved in
Appendices~\ref{appx-sec:support} and~\ref{appx-sec:smoothing}.
After proving the theorems, Subsection~\ref{app:polyhedral-transfer}
deduces the polyhedral corollary and its refinements.
Here \(\rho=\rh(W,Z)\) by the reduction in
Subsection~\ref{sec:outline-smoothing}. Recall that \(P_a(t)=\E F_{a,h}(U_t)\) along the
path~\eqref{eq:path}, so that \(P_a(0)=\E F_{a,h}(S)\) and
\(P_a(1)=\E F_{a,h}(G)\), and that
\(M_h=\sup_a|P_a(1)-P_a(0)|\) by~\eqref{eq:smoothed-discrepancy}.

\begin{proof}[Proof of~\eqref{eq:smooth-moving} in
Proposition~\ref{prop:moving-cutoff}]
The early estimate is integrable at time zero, whereas the late estimate
has a smaller fourth-order term once enough Gaussian variance is present.
Proposition~\ref{an:prop:derivative}, equations~\eqref{an:early}
and~\eqref{an:late}, gives
\[
 \begin{aligned}
 |P_a'(t)|&\le C\sqrt L\left(\frac{t\delta}{h^2}
                    +\frac{\delta^2}{h^3}\right),&&0<t<1,\\
 |P_a'(t)|&\le CL^{3/2}\left(\delta+\frac{\delta^2}{th}\right),
                    &&d_0\delta^2L\le t<1.
 \end{aligned}
\]
Increase \(\Ac\) so that these estimates apply for \(h\ge\Ac\delta\)
and \(\Ac^2\ge d_0\). Put \(\tau_h=Lh^2\), obtained by equating
\(\sqrt L/h^3\) and \(L^{3/2}/(th)\), the two fourth-order factors. Since \(h\ge\Ac\delta\), we have
\(\tau_h\ge d_0\delta^2L\). If \(\tau_h<1\), the later estimate
is therefore available throughout \([\tau_h,1]\). Integrating the
early estimate gives
\begin{align*}
 |P_a(\tau_h)-P_a(0)|
 &\le C\sqrt L\left(
 \frac{\delta\tau_h^2}{h^2}
 +\frac{\delta^2\tau_h}{h^3}\right)\\
 &=C\delta L^{3/2}
 \left(Lh^2+\frac\delta h\right)
 \le C\delta L^{3/2}.
\end{align*}
The later estimate gives
\[
 |P_a(1)-P_a(\tau_h)|
 \le CL^{3/2}\left[
 \delta+\frac{\delta^2}{h}\log\frac1{\tau_h}\right].
\]
The constant cubic contribution integrates to \(C\delta L^{3/2}\);
the fourth-order factor \(1/t\) gives the logarithm. Combining
the two time ranges and taking the supremum over thresholds proves
\eqref{eq:smooth-moving} when \(\tau_h<1\).

If \(Lh^2\ge1\), integrate the early estimate over the whole path:
\[
 |P_a(1)-P_a(0)|
 \le C\sqrt L\left(\frac\delta{h^2}
 +\frac{\delta^2}{h^3}\right)
 =C\delta L^{3/2}\frac{1+\delta/h}{Lh^2}
 \le C\delta L^{3/2}.
\]
The logarithmic term in \eqref{eq:smooth-moving} is zero in this case.
Continuity at the interpolation endpoints permits integration up to
\(0\) and \(1\). All bounds are uniform in the thresholds.
\end{proof}

\begin{proof}[Proof of Theorem~\ref{thm:optimized}]
The smoothed comparison just proved in
Proposition~\ref{prop:moving-cutoff}, equation~\eqref{eq:smooth-moving},
and the smoothing inequality of Lemma~\ref{lem:endpoint},
equation~\eqref{eq:endpoint-min}, give, respectively,
\[
 M_h\le CL^{3/2}\left[\delta+\frac{\delta^2}{h}
          \log_+\frac1{Lh^2}\right],
 \qquad \rho\le M_h+Ch\sqrt L\,\kappa.
\]
Substituting the first bound into the second yields, for every
\(h\ge\Ac\delta\),
\begin{equation}
 \rho\le C\left[
 \delta L^{3/2}
 +\frac{\delta^2}{h}L^{3/2}\log_+\frac1{Lh^2}
 +h\sqrt L\,\kappa\right].
 \label{eq:variational}
\end{equation}
The first term is independent of \(h\); the other two decrease and
increase with the width, respectively. The logarithm prevents an exact
power-law balance, so we use \(E_\kappa\) to bound it at the chosen width.
The admissible choice from~\eqref{eq:optimal-scale} is
\[
 h=\Ac\delta\sqrt{LE_\kappa/\kappa},\qquad
 E_\kappa=1+\log_+\frac{\sqrt\kappa}{\delta L},\qquad
 \kappa=\min\{L,c^{-1}\},
\]
with \(\kappa=L\) when \(c=0\). Because \(\kappa\le L\) and \(E_\kappa\ge1\), this scale is at least
\(\Ac\delta\). Moreover,
\[
 \log_+\frac1{Lh^2}
 =2\log_+\frac{\sqrt\kappa}{\Ac\delta L\sqrt{E_\kappa}}
 \le2\log_+\frac{\sqrt\kappa}{\delta L}
 \le2E_\kappa.
\]
The endpoint term in \eqref{eq:variational} equals
\(\Ac\delta L\sqrt{\kappa E_\kappa}\), while its logarithmic term is at
most \((2/\Ac)\delta L\sqrt{\kappa E_\kappa}\).
This proves \eqref{eq:optimized} after taking the minimum with one.
The version with \(D_L\) follows from \(E_\kappa\le D_L\).
\end{proof}

\begin{proof}[Proof of Theorem~\ref{thm:unrestricted}]
Since \(\kappa\le L\), \(E_\kappa\le D_L\), and \(D_L\ge1\),
\[
 L^{3/2}+L\sqrt{\kappa E_\kappa}
 \le 2L^{3/2}\sqrt{D_L}.
\]
Apply Theorem~\ref{thm:optimized}.
\end{proof}

% Source: zero_skewness_proof
\subsection{\texorpdfstring{Third-moment cancellation: the second part of Proposition~\ref{prop:moving-cutoff} and Theorem~\ref{thm:zero-skew}}{Third-moment cancellation: the second part of Proposition 2 and Theorem 3}}\label{zs:section}
When the aggregate third moments vanish, the cubic contribution
disappears before we take absolute values. We can therefore repeat
the integration using only the fourth-order remainder; the endpoint
error is unchanged.

\begin{proof}[Proof of~\eqref{zs:smooth} in
Proposition~\ref{prop:moving-cutoff}]
The coordinate reduction preserves~\eqref{eq:zero-third}, so
\(\sum_i\E(\zeta_i^0)^{\otimes3}=0\). This cancels the cubic
term in the interpolation identity, leaving
\(P_a'(t)=R_{4,a,h}(t)\), as stated in equation~\eqref{zs:identity}.
The fourth-order remainder estimates of
Proposition~\ref{an:prop:derivative}, equations
\eqref{eq:remainder-early} and~\eqref{an:remainder}, therefore give
\begin{align}
 |P_a'(t)|
 &\le \frac{C\delta^2\sqrt L}{h^3},
 &&0<t<1,\label{zs:early}\\
 |P_a'(t)|
 &\le \frac{C\delta^2L^{3/2}}{t h},
 &&d_0\delta^2L\le t<1.\label{zs:late}
\end{align}
Increase \(\Ac\) so that \(\Ac^2\ge d_0\), and put \(\tau_h=Lh^2\).
If \(\tau_h<1\), then \(\tau_h\ge d_0\delta^2L\). Integrating
\eqref{zs:early} up to \(\tau_h\), and \eqref{zs:late} after it,
yields
\begin{align*}
 |P_a(1)-P_a(0)|
 &\le \frac{C\delta^2\sqrt L}{h^3}\tau_h
       +\frac{C\delta^2L^{3/2}}h\log\frac1{\tau_h}\\
 &=\frac{C\delta^2L^{3/2}}h
       \left[1+\log\frac1{Lh^2}\right].
\end{align*}
If \(Lh^2\ge1\), integrate \eqref{zs:early} on the whole path and
use \(\sqrt L/h^3\le L^{3/2}/h\). Endpoint continuity, already
established after Proposition~\ref{an:prop:derivative}, justifies
integration to \(0\) and \(1\). Taking the supremum over thresholds
proves \eqref{zs:smooth}.
\end{proof}

\begin{proof}[Proof of Theorem~\ref{thm:zero-skew}]
The second part of Proposition~\ref{prop:moving-cutoff} gives
\[
 M_h\le\frac{C\delta^2L^{3/2}}h
       \left[1+\log_+\frac1{Lh^2}\right].
\]
Insert this into the smoothing inequality
\(\rho\le M_h+Ch\sqrt L\,\kappa\) from
Lemma~\ref{lem:endpoint}, equation~\eqref{eq:endpoint-min}.
For every \(h\ge\Ac\delta\), this yields
\begin{equation}
 \rho\le C\left\{
 \frac{\delta^2L^{3/2}}h
 \left[1+\log_+\frac1{Lh^2}\right]+h\sqrt L\,\kappa\right\}.
 \label{zs:variational}
\end{equation}
To balance the comparison error against the endpoint error, choose
\(h=\Ac\delta\sqrt{LE_\kappa/\kappa}\), the width in
equation~\eqref{eq:optimal-scale}, again; it is admissible because \(\kappa\le L\) and \(E_\kappa\ge1\).
As before,
\[
 \log_+\frac1{Lh^2}
 =2\log_+\frac{\sqrt\kappa}{\Ac\delta L\sqrt{E_\kappa}}
 \le2E_\kappa.
\]
The endpoint term in \eqref{zs:variational} equals
\(\Ac\delta L\sqrt{\kappa E_\kappa}\); its comparison term is at most
\[
 \frac{\delta L\sqrt\kappa}{\Ac\sqrt{E_\kappa}}(1+2E_\kappa)
 \le\frac3{\Ac}\delta L\sqrt{\kappa E_\kappa}.
\]
Taking the minimum with one proves \eqref{eq:zero-skew}.

\end{proof}

\subsection{\texorpdfstring{Polyhedral transfer: proof of Corollary~\ref{cor:polyhedra}}{Polyhedral transfer: proof of Corollary 2}}
\label{app:polyhedral-transfer}

Membership in a polyhedron becomes a coordinatewise inequality after
projection onto its face normals. Following the reduction
in~\cite[Section~3.1]{CCK2017}, we standardize these projections and
verify the hypotheses of our rectangle theorem. The same calculation
identifies the covariance and third moments needed for the refinements.

\begin{proof}[Proof of Corollary~\ref{cor:polyhedra}]
The empty set and \(\R^p\) have zero approximation error. For any
other \(P\in\mathcal P_m(K;\Sigma)\), fix a representation satisfying
the stated conditions. If \(q<m\), repeat one of its inequalities
until there are exactly \(m\). This lets us apply the rectangle
bound in a fixed dimension; repeated inequalities leave \(P\)
unchanged and preserve the conditions on its normals.

Let \(T\) be the \(m\times p\) matrix whose \(\alpha\)th row is
\(\nu_\alpha^\top/\sigma_\alpha\), and define
\[
 \eta_i\coloneqq T\xi_i,\qquad
 \widetilde W\coloneqq TW=\sum_i\eta_i,\qquad
 \widetilde Z\coloneqq TZ.
\]
The transformed summands remain independent and centered. Moreover,
\[
 |\eta_{i\alpha}|
 \le\frac{\norm{\nu_\alpha}_1\norm{\xi_i}_\infty}{\sigma_\alpha}
 \le K\delta=\Delta,
 \qquad
 \Var(\widetilde W_\alpha)
 =\frac{\nu_\alpha^\top\Sigma\nu_\alpha}{\sigma_\alpha^2}=1.
\]
The vector \(\widetilde Z\) is centered Gaussian with covariance
\(T\Sigma T^\top=\Cov(\widetilde W)\). Thus
Theorem~\ref{thm:unrestricted} applies in dimension \(m\), with
summand envelope \(\Delta\). For the rectangle
\[
 B\coloneqq\prod_{\alpha=1}^m
 \left(-\infty,\frac{a_\alpha}{\sigma_\alpha}\right],
\]
we have \(\Pp(W\in P)=\Pp(\widetilde W\in B)\) and
\(\Pp(Z\in P)=\Pp(\widetilde Z\in B)\). Applying the theorem
therefore gives the claimed bound. Its right-hand side is independent
of \(P\), so taking the supremum completes the proof.
\end{proof}

For the gap-sensitive and third-moment refinements, use the original
\(q\) inequalities and let \(T\) have the \(q\) rows
\(\nu_\alpha^\top/\sigma_\alpha\). Repeating an inequality would
create two identical coordinates and hence a zero correlation gap.
The preceding calculations give
\[
 \norm{T\xi_i}_\infty\le\Delta,\qquad
 \Cov(TW)=R,\qquad
 R_{\alpha\beta}
 =\frac{\nu_\alpha^\top\Sigma\nu_\beta}
       {\sigma_\alpha\sigma_\beta}.
\]
The diagonal entries of \(R\) are one. Theorem~\ref{thm:optimized}
therefore applies with dimension \(q\), envelope \(\Delta\), and
the correlation gap of \(R\). To check the additional hypothesis
of Theorem~\ref{thm:zero-skew}, apply the linear map to each factor
of the third-moment tensor:
\[
 \sum_i\E(T\xi_i)^{\otimes3}
 =T^{\otimes3}\left(\sum_i\E\xi_i^{\otimes3}\right).
\]
Thus cancellation of the original aggregate tensor implies
cancellation of the transformed tensor, as asserted.

A representation may also contain zero-variance directions, which
must be removed before standardization. Since \(W\) and \(Z\) are centered,
\(\nu_\alpha^\top\Sigma\nu_\alpha=0\) implies
\(\nu_\alpha^\top W=\nu_\alpha^\top Z=0\) almost surely.
If \(a_\alpha\ge0\), removing that inequality leaves both
probabilities unchanged; if \(a_\alpha<0\), both probabilities
are zero. The corollary applies to the remaining inequalities when
they satisfy its stated conditions.

% Source: smoothing_proofs
\section{\texorpdfstring{The logistic approximation: proof of Lemma~\ref{lem:endpoint}}{The logistic approximation: proof of Lemma 2}}\label{appx-sec:smoothing}

We work throughout in \(\R^r\), with the \(m=2p\) unit face normals,
threshold sets \(A_a\), and logistic approximation \(F_{a,h}\) from
Subsection~\ref{sec:outline-smoothing}. We first explain the scope of
the anti-concentration bound, then derive the mixture and derivative
formulas, and finally bound the smoothing error at \(S\) and \(G\).
The derivative formulas and the averaged envelope will also be used
in Appendix~\ref{appx-sec:support}.

\subsection{Anti-concentration for the face inequalities}

Increasing every threshold by \(s\) gives \(A_a^s\), as defined
in Subsection~\ref{sec:outline-smoothing}. Lemma~\ref{lem:concentration}
controls the layer \(A_a^s\setminus A_a^{-s}\), including when
the original inequalities are inconsistent. The same bound applies
to any subsystem: send the unused thresholds to infinity and use
continuity of probability. A random translation independent of the vector being bounded is
handled by conditioning on that translation. Thus we keep the fixed
list of \(m=2p\) normals while allowing the auxiliary sets used
below and in Appendix~\ref{appx-sec:support}.

\subsection{\texorpdfstring{Mixture and derivative formulas}{Mixture and derivative formulas}}\label{appx-sec:mixture}

Fix \(a\in\R^m\) and \(h>0\), and recall
\[
 w_\alpha(x)=e^{(u_\alpha\cdot x-a_\alpha)/h},\qquad
 T=\sum_{\alpha=1}^m w_\alpha,\qquad
 F=F_{a,h}=(1+T)^{-1},\qquad \pi_\alpha=Fw_\alpha.
\]
The face mass \(\theta=1-F\) measures the complement of the baseline
probability. For the curvature argument we also recall the potential
\(\Psi=h\log(1+T)\) and write
\[
 \pi_0\coloneqq F,\qquad
 \bar u\coloneqq\sum_{\alpha=1}^m\pi_\alpha u_\alpha,
 \qquad \chi\coloneqq h\,\tr D^2\Psi.
\]
Here \(\pi_0\) is the probability of the baseline category,
\(\bar u=\nabla\Psi\) is the weighted face normal, and
\(\chi\) is the curvature used in the gap-sensitive endpoint
estimate. The formulas refer to the specified thresholds even when
\(A_a\) is empty. Their denominators remain positive, and limits
with unused thresholds tending to infinity give the corresponding
formulas for subsystems.

The mixture representation~\eqref{eq:mixture} is
\[
 F_{a,h}(x)=\E_\omega\ind_{A_{a,\omega}}(x),\qquad
 A_{a,\omega}=\{x:u_\alpha\cdot x\le
 a_\alpha+h\log(\omega_\alpha/\omega_0),\ \alpha\le m\},
\]
where \(\omega_0,\ldots,\omega_m\) are independent mean-one
exponentials, independent of every random vector at which \(F\) is
evaluated. Conditional on \(\omega_0\), membership requires
\(\omega_\alpha\ge\omega_0w_\alpha(x)\) for every face.
Its probability is therefore \(e^{-\omega_0T(x)}\), whose
expectation is \((1+T(x))^{-1}\). Equivalently,
\[
 F_{a,h}(x)=\Pp\!\left\{-h\log\omega_0\ge
 \max_{\alpha\le m}
 (u_\alpha\cdot x-a_\alpha-h\log\omega_\alpha)\right\}.
\]
The variables \(-\log\omega_\beta\) are independent standard
Gumbel variables. The face displacements have logistic marginals and
are dependent through \(\omega_0\). McFadden~\cite[p.~1001]{McFadden1989}
uses this random-utility representation; it is stated for softmax
weights in~\cite[Lemma~14]{GaoZhang}.

We next express derivatives through positive products of face weights;
this will let us separate the weights from random derivative directions.
Apply Fa\`a di Bruno's higher-order chain rule~\cite{Comtet1974}
to \(F=f(T)\), where
\(f^{(k)}(T)=(-1)^k k!(1+T)^{-k-1}\). Each block of a partition
groups the directions entering one derivative of \(T\), while the
number of blocks determines the derivative of \(f\). Thus, for
\(z_1,\ldots,z_q\in\R^r\),
\begin{equation}
 \begin{split}
 D^qF(x)[z_1,\ldots,z_q]
 ={}&h^{-q}\sum_{\mathcal P\in\Pi_q}(-1)^{|\mathcal P|}|\mathcal P|!
 \sum_{\alpha_1,\ldots,\alpha_{|\mathcal P|}}
 F(x)\prod_{b=1}^{|\mathcal P|}\pi_{\alpha_b}(x)\\
 &\hspace{3em}\times
 \prod_{b=1}^{|\mathcal P|}\prod_{j\in J_b}
 (u_{\alpha_b}\cdot z_j),
 \end{split}
 \label{appx-eq:partition-derivative}
\end{equation}
where \(\Pi_q\) is the set of partitions of \(\{1,\ldots,q\}\)
and \(J_1,\ldots,J_{|\mathcal P|}\) are the blocks in a fixed
ordering. Every face index ranges from \(1\) to \(m\).
Ordered partitions can be encoded by a permutation and cuts between
successive elements, so
\(\sum_{\mathcal P\in\Pi_q}|\mathcal P|!\le q!2^{q-1}\).
In particular,
\begin{equation}
 |D^qF(x)[z_1,\ldots,z_q]|
 \le C^q q!h^{-q}\prod_{j=1}^q\norm{z_j}_U.
 \label{appx-eq:all-orders}
\end{equation}
The softmax cumulant formula~\cite[Lemma~13]{GaoZhang} gives the
same identities through the joint cumulants of a categorical variable.

The all-orders bound discards the factor \(F\theta\), whose expectation
is small by anti-concentration. To retain it, take absolute values
in the partition formula before summing the positive weight products.
For \(\norm v_U\le\varrho\), every \(w_\alpha(x+v)\) and the
denominator \(1+T(x+v)\) change by factors between
\(e^{-\varrho/h}\) and \(e^{\varrho/h}\). Comparing each weight
product with its value at \(x\), and using
\(\sum_{\alpha_1,\ldots,\alpha_k}\prod_b\pi_{\alpha_b}
=\theta^k\le\theta\), yields
\begin{equation}
 \sup_{\norm v_U\le\varrho}|D^qF(x+v)[z_1,\ldots,z_q]|
 \le C_qh^{-q}F(x)(1-F(x))\prod_{j=1}^q\norm{z_j}_U,
 \quad 0\le\varrho\le h.
 \label{appx-eq:local-envelope}
\end{equation}
The comparison holds separately for each positive weight product.
This is the form needed when the derivative directions themselves
are random and must be summed using covariance.

To average the envelope, it suffices that one independent component
of the argument has controlled boundary probabilities. Here an
independent component means \(U=V+R\), with \(V\) independent of \(R\);
conditioning on \(R\) only changes the thresholds.
\begin{lemma}[Averaging the derivative envelope]\label{appx-lem:envelope}
Suppose \(U\in\R^r\) has an independent component \(V\) that is
either a sum of independent centered vectors with summand face norm
at most \(2\delta\), or a centered Gaussian vector. Suppose also
that \(\Var(u_\alpha\cdot V)=v\ge1/2\) for every face.
Uniformly over thresholds, for \(h\ge4\delta\),
\begin{equation}
 \E F(U)(1-F(U))\le Ch\sqrt L.
 \label{appx-eq:small-envelope}
\end{equation}
Consequently, for fixed \(q\), fixed directions
\(z_1,\ldots,z_q\in\R^r\), and \(h\ge\max\{4\delta,\varrho\}\),
\begin{equation}
 \E\sup_{\norm v_U\le\varrho}|D^qF(U+v)[z_1,\ldots,z_q]|
 \le C_qh^{1-q}\sqrt L\prod_{j=1}^q\norm{z_j}_U.
 \label{appx-eq:averaged-derivatives}
\end{equation}
\end{lemma}

\begin{proof}
We compare \(F(1-F)\) with the change in the smoother across a layer
of width \(2h\). Increase every threshold by \(s\), and write
\(F_s=(1+e^{-s/h}T)^{-1}\). Since
\(\partial_sF_s=F_s(1-F_s)/h\) and
\(F_s(1-F_s)\ge e^{-1}F(1-F)\) for \(|s|\le h\),
\[
 F(1-F)\le\frac e2(F_h-F_{-h}).
\]
The mixture representation writes the difference as the average of
\(\ind_{A_{a,\omega}^h}-\ind_{A_{a,\omega}^{-h}}\).
Condition on the exponentials and on \(R\). Apply
Lemma~\ref{lem:concentration} to \(V\), whose face variance is
\(v\ge1/2\). In the bounded-sum case the summand bound is at most
\(2\delta\); in the Gaussian case that term is absent. The expected
layer indicator is at most \(C(h+\delta)\sqrt L\). Since \(h\ge4\delta\), this proves
\eqref{appx-eq:small-envelope}. Combine it with
\eqref{appx-eq:local-envelope} to obtain
\eqref{appx-eq:averaged-derivatives}.
\end{proof}

\subsection{The error at the endpoints}

The mixture representation gives an endpoint bound for every
covariance matrix. Separation of the face normals gives another
bound, which can be sharper. We choose the smaller bound after fixing
the smoothing scale.

\begin{lemma}[Unrestricted endpoint]\label{appx-lem:endpoint-free}
Under~\eqref{eq:setup}, for \(h\ge\delta\),
\begin{equation}
 \E|F_{a,h}(S)-\ind_{A_a}(S)|
 +\E|F_{a,h}(G)-\ind_{A_a}(G)|\le ChL^{3/2}.
 \label{appx-eq:endpoint-free}
\end{equation}
\end{lemma}

\begin{proof}
Recall that \(R_h\) is the largest absolute displacement of a threshold.
The inclusions \(A_a^{-R_h}\subseteq A_{a,\omega}\subseteq A_a^{R_h}\)
reduce the error to a boundary layer. To bound its expected width,
note that, for \(s\ge0\),
\[
 \Pp\{|\log(\omega_\alpha/\omega_0)|>s\}
 =\frac2{1+e^s}\le2e^{-s}.
\]
A union bound over the \(m=2p\) faces gives
\(\Pp(R_h>hs)\le\min\{1,2me^{-s}\}\). Integrating in \(s\)
yields \(\E R_h\le h\{\log(2m)+1\}\le ChL\); independence
of the ratios is not needed.
For fixed exponentials, the indicators can differ only in
\(A_a^{R_h}\setminus A_a^{-R_h}\). Jensen's inequality moves
the absolute value inside the exponential average.
Lemma~\ref{lem:concentration} then bounds the errors at \(S\)
and \(G\) by \(C(R_h+\delta)\sqrt L\) and \(CR_h\sqrt L\),
respectively. Average over the exponentials and use \(h\ge\delta\)
and \(L\ge1\).
\end{proof}

When distinct face normals are separated, the Hessian of \(\Psi\)
also controls the indicator error. Its trace can be averaged using
the covariance identity, giving the following alternative bound.
\begin{lemma}[Curvature endpoint]\label{appx-lem:endpoint-gap}
Under~\eqref{eq:setup}, if \(c>0\) and \(h\ge\delta\), then
\begin{equation}
 \E|F_{a,h}(S)-\ind_{A_a}(S)|
 +\E|F_{a,h}(G)-\ind_{A_a}(G)|
 \le\frac{Ch\sqrt L}{c}.
 \label{appx-eq:endpoint-gap}
\end{equation}
\end{lemma}

\begin{proof}
We first relate the indicator error to \(\chi=h\tr D^2\Psi\),
then bound the expectation of this curvature. Differentiating the
logistic probabilities and using \(\norm{u_\alpha}=1\) gives
\begin{equation}
 D^2\Psi=\frac1h\left(
 \sum_\alpha\pi_\alpha u_\alpha u_\alpha^\top-\bar u\bar u^\top\right),
 \qquad \norm{\bar u}\le\theta,\qquad
 \chi=\theta-\norm{\bar u}^2.
 \label{appx-eq:hessian}
\end{equation}
Let \(\pi_*\coloneqq\max_\alpha\pi_\alpha\). The separation
\(u_\alpha\cdot u_\beta\le1-c\) for \(\alpha\ne\beta\) implies
\[
 \norm{\bar u}^2
 \le(1-c)\theta^2+c\sum_\alpha\pi_\alpha^2.
\]
Consequently,
\[
 \chi\ge F\theta+c\left(\theta^2-\sum_\alpha\pi_\alpha^2\right)
 \ge c\theta(1-\pi_*),\qquad \chi\ge F\theta,
\]
where the second inequality in the first chain uses \(0<c\le1\).
On \(A_a\), every \(w_\alpha\le1\), so \(\pi_*\le1/2\)
and \(\chi\ge c\theta/2\). Off \(A_a\), some weight exceeds
one, so \(\theta>1/2\) and \(\chi\ge F/2\).
The indicator error is \(\theta\) and \(F\) in the two cases,
respectively. Thus
\begin{equation}
 |F-\ind_{A_a}|\le\frac{2\chi}{c}.
 \label{appx-eq:curvature-error}
\end{equation}
It remains to prove
\begin{equation}
 \E\chi(S)+\E\chi(G)\le Ch\sqrt L.
 \label{appx-eq:gap-contraction}
\end{equation}

To average the Hessian at \(S\), we need to compare it with its value
after deleting one summand. Include the baseline index by setting
\(u_0\coloneqq0\) and \(\pi_0=F\). If \(\norm v_U\le\varrho\), each probability
changes by a factor between \(e^{-2\varrho/h}\) and
\(e^{2\varrho/h}\). In the representation
\[
 D^2\Psi=\frac1{2h}\sum_{\alpha,\beta=0}^m
 \pi_\alpha\pi_\beta
 (u_\alpha-u_\beta)(u_\alpha-u_\beta)^\top,
\]
every summand is positive semidefinite. Termwise comparison gives
\[
 e^{-4\varrho/h}D^2\Psi(x)\preceq D^2\Psi(x+v)
 \preceq e^{4\varrho/h}D^2\Psi(x),
\]
where \(M\preceq N\) means that \(N-M\) is positive semidefinite.

Set \(S_i\coloneqq S-\zeta_i^0\). Centering and independence
allow us to subtract \(\nabla\Psi(S_i)\) before integrating the
Hessian along the segment from \(S_i\) to \(S\).
Compare that Hessian first with \(D^2\Psi(S_i)\), then use
independence to replace \(\zeta_i^0(\zeta_i^0)^\top\) by its
expectation. A second comparison replaces \(D^2\Psi(S_i)\) by
\(D^2\Psi(S)\). Since \(\norm{\zeta_i^0}_U\le\delta\) and
the covariances sum to \(I_r\),
\[
 \begin{split}
 \E[S\cdot\nabla\Psi(S)]
 &=\sum_i\E\int_0^1
 (\zeta_i^0)^\top D^2\Psi(S_i+s\zeta_i^0)\zeta_i^0\dd s\\
 &\ge e^{-8\delta/h}\E\tr D^2\Psi(S).
 \end{split}
\]
Also, \(x\cdot\nabla\Psi(x)=\sum_\alpha\pi_\alpha(u_\alpha\cdot x)
\le\norm x_U\). Bernstein's inequality and the expected-maximum bound
in~\cite[Chapter~2]{BLM2013}, applied to the \(m=2p\) face
components of variance one, therefore bound the left side by
\(\E\norm S_U\le C(\sqrt L+\delta L)\).
If \(\delta\sqrt L\le1\), this is \(C\sqrt L\), while
\(h\ge\delta\) bounds \(e^{8\delta/h}\) absolutely. Hence
\(\E\chi(S)\le Ch\sqrt L\). If \(\delta\sqrt L>1\), the
same conclusion follows from \(\chi\le1\) and \(h\ge\delta\).
For the standard Gaussian vector, integration by parts~\cite{Stein1981}
gives directly
\[
 \E\tr D^2\Psi(G)=\E[G\cdot\nabla\Psi(G)]
 \le\E\norm G_U\le C\sqrt L.
\]
This proves~\eqref{appx-eq:gap-contraction}. Combine it with
\eqref{appx-eq:curvature-error} to finish.
\end{proof}

\begin{proof}[Proof of Lemma~\ref{lem:endpoint}]
Lemma~\ref{appx-lem:endpoint-free} bounds the endpoint error by
\(ChL^{3/2}=Ch\sqrt L\cdot L\). When \(c>0\),
Lemma~\ref{appx-lem:endpoint-gap} bounds it by
\(Ch\sqrt L\cdot c^{-1}\). The smaller bound is
\(Ch\sqrt L\,\kappa\), which proves~\eqref{eq:endpoint-error}.
For each \(a\), the probability difference is bounded by
\(|\E F_{a,h}(S)-\E F_{a,h}(G)|\) plus the two endpoint errors.
Taking the supremum over \(a\in\R^m\) yields
\(\rho\le M_h+Ch\sqrt L\,\kappa\), as in~\eqref{eq:endpoint-min}.
\end{proof}

\section{Comparison along the path: proof of Proposition~\ref{an:prop:derivative}}
\label{appx-sec:support}

The early and late bounds estimate the same third- and fourth-order
terms. Their common ingredient is the covariance identity for the
summands: it controls a pair of directions in each derivative without
introducing the covariance rank. Near the bounded endpoint, we use
the local derivative envelope. Once enough Gaussian variance is
present, we move that pair of derivatives onto the Gaussian density.

We first record the covariance and concentration facts needed below,
then derive the exact interpolation identity. Subsections~\ref{appx-an:sec:early}
and~\ref{appx-an:sec:restoration} then prove the two bounds, using
the matrix estimate of Subsection~\ref{appx-an:sec:matrix} for the
late bound. Subsection~\ref{appx-an:sec:surface} supplies the surface
estimate needed in that matrix argument. The full face family has
\(m=2p\) members; when needed, estimates for auxiliary subsystems
follow by sending unused thresholds to infinity.

\subsection{Covariance and concentration}
\label{appx-an:sec:coords}

We continue with the summands and unit face normals of
Subsection~\ref{sec:outline-smoothing}. Recall the two properties
that drive the derivative estimates:
\begin{equation}
 \sum_i\E\zeta_i^0(\zeta_i^0)^\top=I_r,
 \qquad \norm{\zeta_i^0}_U\le\delta,
 \qquad \norm x_U=\max_{j\le m}|u_j\cdot x|,
 \quad m=2p.
 \label{appx-eq:projected-bound}
\end{equation}
The covariance identity controls pairs of derivative directions;
the face bound controls the remaining directions. Write
\[
 q_j(x)\coloneqq(u_j\cdot x-a_j)/h,\quad
 N(x)\coloneqq1+\sum_{j=1}^m e^{q_j(x)},\quad
 F=F_{a,h}=N^{-1},\quad\pi_j=e^{q_j}/N,\quad\theta=1-F.
\]
These are the function and weights from Appendix~\ref{appx-sec:smoothing}.
We now state a conditional version of the pair estimate, because
the late-bound proof will condition on the original summands and masks.
It uses second moments and does not require conditional centering. If a random array \((Z_i)_i\)
and a \(\sigma\)-field \(\mathcal H\) satisfy
\(\sum_i\E[Z_iZ_i^\top\mid\mathcal H]\preceq I_r\), then, for
fixed or \(\mathcal H\)-measurable unit vectors \(u,v\),
conditional Cauchy--Schwarz gives
\begin{equation}
 \sum_i\E[|(u\cdot Z_i)(v\cdot Z_i)|\mid\mathcal H]
 \le\left\{\sum_i\E[(u\cdot Z_i)^2\mid\mathcal H]
           \sum_i\E[(v\cdot Z_i)^2\mid\mathcal H]\right\}^{1/2}
 \le1.
 \label{appx-an:pair-contraction}
\end{equation}
Bounding the other directions by \(\delta\) will therefore cost
\(\delta\) for a third derivative and \(\delta^2\) for a fourth
derivative. The same second-moment bound will later control coefficient
matrices in operator norm.

The second ingredient is a bound on boundary probabilities. Fix a
positive absolute constant \(C_0\), large enough for the applications
below. With the
threshold sets \(A_a\) and their enlargements already defined in
Subsection~\ref{sec:outline-smoothing}, we call \(V_0\) a
\emph{concentrated base} if
\begin{equation}
 \Pp\{V_0\in A_a^s\setminus A_a^{-s}\}
 \le C_0(s+\delta)\sqrt L,
 \qquad a\in\R^m,\quad s>0.
 \label{appx-an:concentration}
\end{equation}
Empty intersections are allowed. The property passes to subsystems
by sending unused thresholds to infinity, and to independent shifts
by conditioning. Lemma~\ref{lem:concentration} establishes it whenever an independent component is
either a bounded sum with summand face bound \(C\delta\), or a
Gaussian vector, and its equal face variances are bounded below by
a positive absolute constant.

This boundary estimate is the only property of \(V_0\) needed in the
matrix and surface arguments below. Repeating the envelope argument of
Lemma~\ref{appx-lem:envelope} gives, for \(h\ge\delta\),
\begin{equation}
 \sup_z\E F(V_0+z)\theta(V_0+z)\le Ch\sqrt L.
 \label{appx-an:weight-concentration}
\end{equation}
Indeed, shifting all thresholds by \(\pm h\) gives
\(F(1-F)\le(e/2)(F_h-F_{-h})\). The mixture representation writes
the difference as an average of boundary layers, each controlled by
\eqref{appx-an:concentration}. The same argument for one face bounds
\(e^{q_j}/(1+e^{q_j})^2\). We use these facts after conditioning on
an independent Gaussian shift; no conditional concentration of
\(V_0\) itself is required.

\subsection{The path and its derivative}
\label{appx-an:sec:path}

The path~\eqref{eq:path} can be written as
\begin{equation}
 Y_i\coloneqq \gamma\mathbb B_{i,t}\zeta_i^0,\quad
 V_t\coloneqq \sum_iY_i,\quad U_t=V_t+\sqrt t\,G,
 \quad \gamma=1+t,\quad
 \varpi=\frac{1-t}{(1+t)^2},
 \label{appx-an:path}
\end{equation}
where \(\mathbb B_{i,t}\sim\operatorname{Bernoulli}(\varpi)\)
and \(G\sim N(0,I_r)\) are independent of the summands and of one
another. Only the law at each \(t\) is used. Since
\(\varpi\gamma^2=1-t\), the covariance of \(U_t\)
is \(I_r\), and \(P_a(t)=\E F(U_t)\).
The additional calibration \(\varpi\gamma^3=1-t^2\)
makes the cubic coefficient stationary at \(t=0\).

To differentiate the contribution of summand \(i\), first remove it
from the path and put \(H_i\coloneqq U_t-Y_i\). We insert independent
copies at this base: let \(\zeta_i,\widetilde\zeta_i\) have the law of
\(\zeta_i^0\), with all copies mutually independent and independent
of the original summands, masks, and Gaussian variables.
For repeated directions, write
\(D^qF(x)[z^q]=D^qF(x)[z,\ldots,z]\).
To see the cancellations before expanding the remainder, define the
signed density
\[
 \nu_t(s)\coloneqq \frac{4-2(3-t)s}{\gamma},\qquad 0\le s\le1.
\]
Then
\begin{equation}
 P_a'(t)=\frac12\sum_i\E\left[
 D^2F(U_t)[\zeta_i^2]
 -\int_0^1\nu_t(s)
       D^2F(H_i+s\gamma\zeta_i)[\zeta_i^2]\dd s\right].
 \label{appx-an:signed-hessian}
\end{equation}
Its two relevant moments are
\begin{equation}
 \int_0^1\nu_t(s)\dd s=1,\qquad
 \int_0^1s\nu_t(s)\dd s=\frac{2t}{3\gamma}.
 \label{appx-an:signed-moments}
\end{equation}
The mass one cancels the constant Hessian term; the first moment
produces the coefficient \(-t/3\) of the cubic term.
Although \(\nu_t\) has mass one, it can take negative values; the
identity is an algebraic cancellation, not an averaging identity
under a probability law. Its total variation is at most \(5/3\).

To prove~\eqref{appx-an:signed-hessian}, differentiate each Bernoulli
factor of the joint law, then the amplitudes and the Gaussian variance.
In the contribution of the \(i\)th factor or amplitude, hold the other
summands fixed; their dependence on \(t\) is accounted for by their
own contributions. Thus \(H_i\) serves as a fixed base in that term,
not as a quantity whose law is constant along the path. This gives
\begin{align*}
 P_a'(t)
 ={}&\frac12\sum_i\E D^2F(U_t)[\zeta_i^2]\\
 &+\sum_i\left[\varpi'\E\{F(H_i+\gamma\zeta_i)-F(H_i)\}
             +\varpi\E DF(H_i+\gamma\zeta_i)[\zeta_i]\right].
\end{align*}
The first term comes from differentiating the Gaussian variance.
For a fixed \(v\), differentiating under the Gaussian integral and
applying Stein's identity~\cite{Stein1981} gives the heat identity
\[
 \frac{d}{dt}\E_G F(v+\sqrt t\,G)
 =\frac12\E_G\tr D^2F(v+\sqrt t\,G).
\]
Averaging over \(v=V_t\) gives the contribution of the Gaussian
factor at time \(t\); the change in the law of \(V_t\) contributes
the other terms. Since the copies \(\zeta_i\) are independent of
\(U_t\) and \(\sum_i\E\zeta_i\zeta_i^\top=I_r\),
\(\E\tr D^2F(U_t)=\sum_i\E D^2F(U_t)[\zeta_i^2]\).
At a fixed base \(x\), centering and the fundamental theorem of calculus give
\begin{align*}
 \E\{F(x+\gamma\zeta_i)-F(x)\}
 &=\gamma^2\int_0^1(1-s)
       \E D^2F(x+s\gamma\zeta_i)[\zeta_i^2]\dd s,\\
 \E DF(x+\gamma\zeta_i)[\zeta_i]
 &=\gamma\int_0^1\E D^2F(x+s\gamma\zeta_i)[\zeta_i^2]\dd s.
\end{align*}
Since \(\varpi'=-(3-t)/\gamma^3\), their combined coefficient
is \(\varpi'\gamma^2(1-s)+\varpi\gamma=-\nu_t(s)/2\).
This proves~\eqref{appx-an:signed-hessian}; integrating the affine
density gives~\eqref{appx-an:signed-moments}. All differentiations
are justified for fixed \(h>0\) by bounded summands and bounded
derivatives of \(F\), using dominated differentiation under the
integral~\cite[Section~2.3]{Folland1999}.

We now put the cubic term at the full path, where aggregate third
moments can cancel. The fourth-order remainder combines three
corrections: the Hessian expansion along an independent copy, the
difference between full and deleted bases for the quadratic term, and
the restoration of the cubic term to the full path. Exchanging the
names of the independent copies lets all three corrections use the
same fourth-derivative evaluation point. The first kernel agrees with
\cite[Lemma~4.1]{AL}, and the resulting full-path form agrees with
the structure of~\cite[Lemma~22]{GaoZhang}.

\begin{lemma}[Exact fourth-order identity]\label{appx-an:lem:identity}
For \(0<t<1\),
\begin{equation}
 \begin{aligned}
 P_a'(t)
 ={}&-\frac t3\E D^3F(U_t):\sum_i\E\zeta_i^{\otimes3}\\
 &+\int_0^1\sum_i\E\left[
 D^4F(H_i+s\gamma\zeta_i):\mathcal T_{i,t}(s)
 \right]\dd s,
 \end{aligned}
 \label{appx-an:fourth-identity}
\end{equation}
where, in the support coordinates,
\[
 \begin{aligned}
 \mathcal T_{i,t}(s)\coloneqq{}&a_t(s)\zeta_i^{\otimes4}\\
 &+b_t(s)\widetilde\zeta_i^{\otimes2}\otimes\zeta_i^{\otimes2}\\
 &+c_t(s)\widetilde\zeta_i^{\otimes3}\otimes\zeta_i,
 \end{aligned}
\]
and the scalar kernels are
\begin{equation}
 \begin{aligned}
 a_t(s)&\coloneqq\frac\gamma6(1-s)^2\{(3-t)s-2t\},\\
 b_t(s)&\coloneqq\frac{1-t}{2}(1-s),\\
 c_t(s)&\coloneqq\frac{t(1-t)}{3\gamma}.
 \end{aligned}
 \label{appx-an:kernel}
\end{equation}
Their absolute integrals over \(0\le s\le1\) are uniformly bounded
for \(0<t<1\). The colon denotes contraction over corresponding
tensor indices; no explicit symmetrization of \(\mathcal T_{i,t}\)
is needed because \(D^4F\) is symmetric.
\end{lemma}

\begin{proof}
Expand the segment Hessian in~\eqref{appx-an:signed-hessian}
about \(H_i\), through its linear term. The two moments give
\begin{align*}
 P_a'(t)
 ={}&\frac12\sum_i\E\{(D^2F(U_t)-D^2F(H_i))[\zeta_i^2]\}\\
 &-\frac t3\sum_i\E D^3F(H_i)[\zeta_i^3]
       +\sum_i\mathcal A_i,
\end{align*}
where
\begin{equation}
 \mathcal A_i\coloneqq\int_0^1 a_t(s)
 \E D^4F(H_i+s\gamma\zeta_i)
       [\zeta_i,\zeta_i,\zeta_i,\zeta_i]\dd s.
 \label{appx-an:A}
\end{equation}
Indeed, the second-order integral remainder has kernel
\[
 -\frac{\gamma^2}{2}\int_s^1(u-s)\nu_t(u)\dd u
 =\frac\gamma6(1-s)^2\{(3-t)s-2t\}=a_t(s).
\]
In the Hessian difference, the deleted summand \(Y_i\) is centered
and independent of \(H_i,\zeta_i\). A second-order expansion has
zero expected linear term. Since \(\varpi\gamma^2=1-t\), its
contribution is
\begin{equation}
 \mathcal B_i\coloneqq\int_0^1 b_t(s)
 \E D^4F(H_i+s\gamma\widetilde\zeta_i)
       [\zeta_i,\zeta_i,\widetilde\zeta_i,\widetilde\zeta_i]\dd s.
 \label{appx-an:B}
\end{equation}
Finally, add and subtract \(D^3F(U_t)\) in the cubic term.
A first-order integral expansion of that difference, using
\(\varpi\gamma=(1-t)/\gamma\), gives
\begin{equation}
 \mathcal C_i\coloneqq\int_0^1 c_t(s)
 \E D^4F(H_i+s\gamma\widetilde\zeta_i)
       [\zeta_i,\zeta_i,\zeta_i,\widetilde\zeta_i]\dd s.
 \label{appx-an:C}
\end{equation}
The remainder is therefore \(\sum_i(\mathcal A_i+\mathcal B_i+\mathcal C_i)\).
The pair \((\zeta_i,\widetilde\zeta_i)\) has an exchangeable joint
law and is independent of \(H_i\). Interchanging the names of these
copies in \(\mathcal B_i\) and \(\mathcal C_i\) puts every fourth
derivative at \(H_i+s\gamma\zeta_i\) and produces the three tensor
products in \(\mathcal T_{i,t}(s)\). Independence of \(\zeta_i\)
and \(U_t\) gives the displayed cubic contraction. This proves
\eqref{appx-an:fourth-identity}. Finally, \(\gamma\le2\) and the
explicit kernels give
\( |a_t(s)|\le(1-s)^2\),
\(0\le b_t(s)\le(1-s)/2\), and
\(0\le c_t(s)\le1/12\), which prove the uniform integral bound.
\end{proof}

All three corrections are evaluated near the same full path:
\begin{equation}
 H_i+s\gamma\zeta_i=U_t+v_i,\qquad
 v_i\coloneqq s\gamma\zeta_i-Y_i,\qquad \norm{v_i}_U\le4\delta,
 \label{appx-an:translates}
\end{equation}
and the same holds with \(\widetilde\zeta_i\). The bound follows
from \(\gamma\le2\) and the summand face bound. The scalar kernels
in the three corrections have uniformly bounded absolute integrals.
We must preserve the cubic term at its common base until any
aggregate third-moment cancellation has been used.

\subsection{The early bounds}
\label{appx-an:sec:early}

We split the Gaussian part so that one half can provide concentration
and the other can absorb derivatives in the late estimate. For each
fixed time, take independent standard Gaussian vectors \(G_0,G_1\),
independent of the summands and masks, and represent the path in law as
\begin{equation}
 U_t=V_{0,t}+\sigma_t G_1,\qquad
 V_{0,t}=V_t+\sqrt{t/2}\,G_0,\qquad
 \sigma_t^2=t/2.
 \label{appx-an:gaussian-split}
\end{equation}
This leaves its law unchanged. The bounded and Gaussian components
of \(V_{0,t}\) have equal face variances \(1-t\) and
\(t/2\), respectively. For \(t\le2/3\), the bounded component
has face variance at least \(1/3\); for \(t\ge2/3\), the Gaussian
component does. Apply Lemma~\ref{lem:concentration} to that component
and condition on the other. Rescaling a variance bounded below by
\(1/3\) changes only the absolute constant, and the bounded component
has summand face bound \(2\delta\). Thus \(V_{0,t}\) is a concentrated
base for every \(0\le t\le1\).
By~\eqref{appx-an:weight-concentration},
\begin{equation}
 \E F(U_t)\theta(U_t)\le Ch\sqrt L.
 \label{appx-an:path-weight}
\end{equation}

The early estimate bounds all derivatives directly. In the cubic
term and the \(a_t\) term of \(\mathcal T_{i,t}\), retain the pair
\((\zeta_i,\zeta_i)\); in the \(b_t\) and \(c_t\) terms, retain
\((\widetilde\zeta_i,\widetilde\zeta_i)\). The remaining one or two
directions have face norm at most \(\delta\).
We retain the repeated pair for the covariance estimate, bounding
the remaining directions by \(\delta\). For example,~\eqref{appx-an:pair-contraction}
gives, for any face indices \(a,b\),
\[
 \sum_i\E|(u_a\cdot\zeta_i)(u_b\cdot\zeta_i)|\le1.
\]
The same inequality holds for \(\widetilde\zeta_i\). This gives the
moment factor \(\delta^{q-2}\) for derivative order \(q=3,4\),
including the mixed-copy terms.

The shift \(v_i\) may depend on the same copy that supplies the
derivative directions. To separate these moments from the derivatives,
first compare the positive logistic weights at \(U_t+v_i\) with
those at \(U_t\), before using independence. If
\(\norm v_U\le4\delta\) and \(h\ge4\delta\), every weight
\(F\prod_{a=1}^k\pi_{\alpha_a}\), with \(k\le4\), changes by
at most an absolute factor. Apply the derivative formula
\eqref{appx-eq:partition-derivative} at \(U_t+v_i\) and use
this comparison term by term. The copies are independent of \(U_t\),
so the preceding moment bound applies. Finally,
the sum of weights with \(k\) face factors is
\(F\theta^k\le F\theta\). Writing \(\zeta_i\) for the retained
copy in each case, we obtain, for each direction pattern in the identity,
\[
 \begin{aligned}
 \sum_i\E\left|D^qF(U_t+v_i)
       [\zeta_i,\zeta_i,z_{i,3},\ldots,z_{i,q}]\right|
 &\le C\delta^{q-2}h^{-q}\E F(U_t)\theta(U_t)\\
 &\le C\delta^{q-2}h^{1-q}\sqrt L.
 \end{aligned}
\]
The cubic term has the additional coefficient \(t/3\).
The bounded integrals of the fourth-order kernels therefore give
\[
 |R_{4,a,h}(t)|\le C\delta^2\sqrt L/h^3,\qquad
 |P_a'(t)|\le C\sqrt L
       \left(t\delta/h^2+\delta^2/h^3\right).
\]
These are~\eqref{eq:remainder-early} and~\eqref{an:early}.
The moment argument is the one in~\cite[Lemma~4.1]{AL}; the
logistic envelope supplies its boundary factor without an eigenvalue
condition.

\subsection{The late bounds}
\label{appx-an:sec:restoration}

For the late estimate we move the repeated pair of derivatives onto
the Gaussian density. This replaces two factors \(h^{-1}\) by
\(\sigma^{-2}\). To keep the bound independent of rank, we then
sum the coefficient matrices using the second-moment bound before
taking the nuclear norm of the Gaussian average.

The following matrix estimate, proved in
Subsection~\ref{appx-an:sec:matrix}, controls the Gaussian factor
\(GG^\top-I_r\) produced by two integrations by parts. For a concentrated base \(V_0\),
an independent \(G\sim N(0,I_r)\), \(U=V_0+\sigma G\),
\(\delta\sqrt L\le\sigma\le1\), and \(h\ge\delta\), it states
that, for every integer \(k\ge1\),
\begin{equation}
 \E_{V_0}\sum_{\alpha_1,\ldots,\alpha_k}
 \left\|\E_G\left[(GG^\top-I_r)F(U)
              \prod_{a=1}^k\pi_{\alpha_a}(U)\right]\right\|_*
 \le C_1^{k+1}hL^{3/2}.
 \label{appx-an:matrix-bound}
\end{equation}
The sum is over ordered face tuples, \(C_1\ge1\) is absolute,
and \(\norm{\cdot}_*\) denotes the nuclear norm, the sum of singular
values. The next lemma brings derivatives at different nearby points
to the common base where this estimate applies. Its auxiliary vectors \(Z_i,Z_{i,a}\) are unrelated to the
Gaussian target \(Z\).

\begin{lemma}[Paired derivatives at translated bases]
\label{appx-an:lem:restoration}
Let \(\mathcal H\) be a \(\sigma\)-field, let \(V_0\) be an
\(\mathcal H\)-measurable concentrated base, let \(G\sim N(0,I_r)\) be
independent of \(\mathcal H\), let \(\delta\sqrt L\le\sigma\le1\), and put
\(U\coloneqq V_0+\sigma G\). For \(s\in\{1,2\}\), let
\((Z_i,Z_{i,1},\ldots,Z_{i,s},v_i)_i\) be random arrays which, conditional
on \(\mathcal H\), are independent of~\(G\). They may depend on
\(\mathcal H\), and on one another. Suppose that, almost surely,
\begin{equation}
 \sum_i\E[Z_iZ_i^\top\mid\mathcal H]\preceq I_r,
 \qquad \norm{Z_{i,a}}_U\le\delta,
 \qquad \norm{v_i}_U\le b\delta.
 \label{appx-an:conditional-budget}
\end{equation}
There is an absolute constant \(\mathsf A_b\), depending only on \(b\),
such that \(h\ge\mathsf A_b\delta\) implies
\begin{equation}
 \left|\sum_i\E D^{s+2}F(U+v_i)
         [Z_i,Z_i,Z_{i,1},\ldots,Z_{i,s}]\right|
 \le\frac{C_b}{\sigma^2}\,
                 \delta^s h^{1-s}L^{3/2}.
 \label{appx-an:restoration-bound}
\end{equation}
\end{lemma}

\begin{proof}
Condition on \(\mathcal H\) and on the arrays, leaving only \(G\)
random. This conditioning is used for Gaussian integration by parts;
we average again before applying the matrix estimate. Concentration
is required only for the unconditional law of \(V_0\), not for its
conditional law given \(\mathcal H\). Applying Stein's
identity~\cite{Stein1981} twice gives
\begin{align*}
 &\E_G D^{s+2}F(U+v_i)[Z_i,Z_i,Z_{i,1},\ldots,Z_{i,s}]\\
 &\qquad=\sigma^{-2}\E_G\left[
       \tr\{(GG^\top-I_r)Z_iZ_i^\top\}
       D^sF(U+v_i)[Z_{i,1},\ldots,Z_{i,s}]\right].
\end{align*}
Only first and second derivatives remain. With
\(W_j\coloneqq F\pi_j\) and \(W_{jk}\coloneqq F\pi_j\pi_k\), they are
\begin{align*}
 DF(x)[z]&=-h^{-1}\sum_jW_j(x)(u_j\cdot z),\\
 D^2F(x)[z,w]&=h^{-2}\left\{
       2\sum_{j,k}W_{jk}(x)(u_j\cdot z)(u_k\cdot w)
       -\sum_jW_j(x)(u_j\cdot z)(u_j\cdot w)\right\}.
\end{align*}
Thus there are only three forms to estimate. For each, write
\(W_\alpha\coloneqq F\prod_{a=1}^k\pi_{\alpha_a}\), with \(k=1\) or \(2\),
and let \(c_{i,\alpha}\) denote the corresponding product of the
\(s\) remaining face projections. Then \(|c_{i,\alpha}|\le\delta^s\);
the numerical coefficients have absolute value at most two.

We next express every translated weight at the common base. Put
\[
 d_j(v)\coloneqq e^{u_j\cdot v/h}-1,\qquad
 z_v(x)\coloneqq \sum_j\pi_j(x)d_j(v),\qquad
 \epsilon\coloneqq e^{b\delta/h}-1.
\]
Since \(N(x+v)=N(x)(1+z_v(x))\),
\begin{equation}
 W_\alpha(x+v)=W_\alpha(x)
 e^{\sum_a u_{\alpha_a}\cdot v/h}(1+z_v(x))^{-k-1}.
 \label{eq:weight-translation}
\end{equation}
If \(\norm v_U\le b\delta\), then \(|d_j(v)|,|z_v(x)|\le\epsilon\).
For \(\epsilon<1\), expand
\[
 (1+z_v)^{-k-1}
 =\sum_{\ell\ge0}(-1)^\ell\binom{k+\ell}{\ell}
       \sum_{\beta_1,\ldots,\beta_\ell}
       \prod_{q=1}^\ell\pi_{\beta_q}(x)d_{\beta_q}(v).
\]
The expansion is useful because every weight depending on \(G\) is
now evaluated at \(U\), while the translation enters only through
coefficients independent of \(G\). For fixed face tuples \(\alpha,\beta\), all dependence on the
summands is now in the symmetric coefficient matrix
\[
 \Xi_{\alpha,\beta}\coloneqq \sum_i\E\left[
 Z_iZ_i^\top c_{i,\alpha}
 e^{\sum_a u_{\alpha_a}\cdot v_i/h}
 \prod_{q=1}^\ell d_{\beta_q}(v_i)\,\middle|\,\mathcal H\right].
\]
For any vector \(z\), the conditional second-moment bound implies
\[
 |z^\top\Xi_{\alpha,\beta}z|
 \le\delta^s e^{kb\delta/h}\epsilon^\ell
       \sum_i\E[(z\cdot Z_i)^2\mid\mathcal H]
 \le\delta^s e^{kb\delta/h}\epsilon^\ell\norm z^2.
\]
Hence \(\norm{\Xi_{\alpha,\beta}}_\op\) has the same scalar bound.
This is the matrix use of the second-moment bound underlying
\eqref{appx-an:pair-contraction}; replacing it by a trace bound
would introduce the rank.

Apply the operator--nuclear norm inequality
\(|\tr(\Xi M)|\le\norm\Xi_\op\norm M_*\)
from~\cite[Chapter~IV]{Bhatia1997}, sum over face
tuples, and use~\eqref{appx-an:matrix-bound} with \(k+\ell\) factors.
After restoring \(\sigma^{-2}h^{-s}\), each of the three forms
is bounded by \(\sigma^{-2}\delta^sh^{1-s}L^{3/2}\) times
\[
 e^{kb\delta/h}C_1^{k+1}
 \sum_{\ell\ge0}\binom{k+\ell}{\ell}(C_1\epsilon)^\ell
 =e^{kb\delta/h}C_1^{k+1}(1-C_1\epsilon)^{-k-1}.
\]
Choose \(\mathsf A_b\) so that \(h\ge\mathsf A_b\delta\) implies
\(C_1\epsilon\le1/2\). Since \(k\le2\), the last expression is
bounded by \(C_b\), proving~\eqref{appx-an:restoration-bound}.

Fubini's theorem~\cite[Section~2.5]{Folland1999} justifies the
series and expectation interchanges using the convergent scalar
majorant \((1-\epsilon)^{-k-1}\), the bound
\(\sum_i\E\norm{Z_i}^2\le r\), and Gaussian integrability.
This finite preliminary bound is used only for the interchanges;
the operator-norm argument supplies the stated bound.
\end{proof}

\begin{proof}[Proof of the late bounds \eqref{an:late}
and~\eqref{an:remainder}]
Use the decomposition~\eqref{appx-an:gaussian-split} and set
\[
 \mathcal H\coloneqq
 \sigma\bigl((\zeta_i^0,\mathbb B_{i,t})_{i=1}^n,G_0\bigr).
\]
The base \(V_{0,t}\) is concentrated, \(G_1\) is independent
of \(\mathcal H\), and \(\sigma_t^2=t/2\).
Thus \(t\ge d_0\delta^2L\), with \(d_0\ge2\), ensures
\(\sigma_t\ge\delta\sqrt L\). Both arrays of independent copies
remain independent of \(\mathcal H\), so their conditional aggregate
covariances are \(I_r\).

For the cubic term, take \(s=1\), \(Z_i=Z_{i,1}=\zeta_i\), and
\(v_i=0\) in Lemma~\ref{appx-an:lem:restoration}. It gives
\begin{equation}
 \left|\frac t3\sum_i\E D^3F(U_t)[\zeta_i^3]\right|
 \le Ct\sigma_t^{-2}\delta L^{3/2}
 \le C\delta L^{3/2}.
 \label{appx-an:cubic-bound}
\end{equation}
For each fixed point of the remainder integral, use the same lemma
with two remaining directions (its parameter \(s=2\)). The integration
point is held fixed during this application.
Use the common shifts in~\eqref{appx-an:translates}, with \(b=4\). For the \(a_t\) term
of \(\mathcal T_{i,t}\), take \(Z_i=\zeta_i\), leaving directions
\((\zeta_i,\zeta_i)\). For the \(b_t\) and \(c_t\) terms,
take \(Z_i=\widetilde\zeta_i\), leaving directions
\((\zeta_i,\zeta_i)\) and \((\widetilde\zeta_i,\zeta_i)\),
respectively. These directions have face norms at most \(\delta\),
so each summed derivative contribution is bounded by
\(C\delta^2L^{3/2}/(t h)\). Integrating the bounded scalar
kernels proves~\eqref{an:remainder}; combining it with the cubic
bound proves~\eqref{an:late}.
\end{proof}

At the endpoints, \(U_0\) has the law of \(S\), and
\(U_1\) has the law of \(G\). Boundedness and continuity
of \(F\), together with the integrable early derivative bound,
justify integration up to both endpoints. This completes the proof
of Proposition~\ref{an:prop:derivative}, subject to the matrix
estimate proved next.

\subsection{A matrix estimate for the logistic weights}
\label{appx-an:sec:matrix}

We now prove~\eqref{appx-an:matrix-bound}. The key observation is that
tilting a standard Gaussian law by a log-concave weight can only
decrease its covariance, by Pr\'ekopa's
theorem~\cite[Theorem~6]{Prekopa1973}. This reduces the nuclear norm to two scalar
quantities: the squared length of the tilted mean and a signed trace
term. We control the means relative to a single face and sum the
weights before estimating the trace. The latter then becomes a
surface integral, bounded in Subsection~\ref{appx-an:sec:surface}.
For a related application to convex cells,
see~\cite[Lemma~12]{GaoZhang}.

Throughout this subsection, \(V_0\) is a concentrated base,
\(G\sim N(0,I_r)\) is independent of \(V_0\),
\(\delta\sqrt L\le\sigma\le1\), and \(U=V_0+\sigma G\).
We use the face scores \(q_j\) from
Subsection~\ref{appx-an:sec:coords}, so
\(N=1+\sum_j e^{q_j}\), \(F=N^{-1}\),
\(\pi_j=e^{q_j}/N\), and \(\theta=1-F\).

For the mean estimate, comparison with the constant weight one would
introduce the logarithm of the reciprocal Gaussian mass of the weight.
Instead, we compare it with a weight depending on one normal coordinate.
The component along that normal is controlled by a second moment.
The component orthogonal to it is controlled by the logarithm of the
ratio of reference mass to tilted mass. This refinement of the entropy argument
in~\cite[Lemma~12]{GaoZhang} avoids a logarithm of the smoothing scale.

\begin{lemma}[Gaussian mean relative to a single face]
\label{appx-an:lem:mean}
Let \(u\) be a unit vector, and let \(b(g)\coloneqq \beta(u\cdot g)\in[0,1]\) have
positive Gaussian integral \(p_b\). If \(0\le w\le b\),
\(p_w\coloneqq \E w(G)>0\), and \(\mu_w\coloneqq \E[Gw(G)]/p_w\), then
\begin{equation}
 p_w\norm{\mu_w}^2
 \le \E[w(G)(u\cdot G)^2]
        +2p_w\log(p_b/p_w).
 \label{appx-an:mean-entropy}
\end{equation}
\end{lemma}

\begin{proof}
The Gaussian component orthogonal to \(u\) is independent of the
reference weight \(b\). For \(v\perp u\), Jensen's inequality under
the law with density \(w/p_w\), followed by \(w\le b\), therefore gives
\[
 p_we^{v\cdot\mu_w}
 \le\E[w(G)e^{v\cdot G}]
 \le\E[b(G)e^{v\cdot G}]
 =p_be^{\norm v^2/2}.
\]
Take \(v\) to be the projection of \(\mu_w\) onto \(u^\perp\).
Taking logarithms bounds its squared length by
\(2\log(p_b/p_w)\). Cauchy--Schwarz gives
\(p_w|u\cdot\mu_w|^2\le\E[w(G)(u\cdot G)^2]\) for the normal
component. Adding the two bounds proves the claim.
\end{proof}

We can now control the matrix sum without a factor depending on the
rank. We must sum the signed trace terms before taking an absolute value.
The resulting scalar weight has a surface representation, to which
the next subsection's estimate applies.

\begin{lemma}[Matrix estimate for logistic weights]
\label{appx-an:lem:matrix}
If \(h\ge\delta\), then~\eqref{appx-an:matrix-bound} holds for every
integer \(k\ge1\).
\end{lemma}

\begin{proof}
We first record the Gaussian-tilt calculation. Let \(w:\R^r\to[0,1]\)
be log-concave with \(p_w\coloneqq \E_Gw(G)>0\). Define the tilted
probability measure by
\[
 dQ_w(g)\coloneqq\frac{w(g)}{p_w}\,d\Phi_r(g),
\]
and let \(\mu_w,K_w\) be its mean and covariance. Pr\'ekopa's
theorem~\cite[Theorem~6]{Prekopa1973} makes
\[
 z\longmapsto\int_{\R^r}e^{-\norm{g-z}^2/2}w(g)\dd g
\]
log-concave. Differentiating under the integral shows that at zero
the gradient of its logarithm is \(\mu_w\), and the Hessian is
\(\E_{Q_w}GG^\top-\mu_w\mu_w^\top-I_r=K_w-I_r\).
Log-concavity therefore implies \(K_w\preceq I_r\).
This argument includes convex-cell indicators, which are log-concave.
For smooth positive \(w\), the same bound follows from the
Brascamp--Lieb variance inequality~\cite{BrascampLieb1976}, applied
to linear functions: the potential
\(V(g)=\norm g^2/2-\log w(g)\) has \(\nabla^2V\succeq I_r\),
so \(\operatorname{Var}_{Q_w}(a\cdot G)\le\norm a^2\) for every
\(a\in\R^r\).
Writing
\[
 \E_G[(GG^\top-I_r)w(G)]
 =p_w\{\mu_w\mu_w^\top-(I_r-K_w)\}
\]
as \(A-B\), with \(A\coloneqq p_w\mu_w\mu_w^\top\) and
\(B\coloneqq p_w(I_r-K_w)\), we have \(A,B\succeq0\). The triangle
inequality for the nuclear norm gives
\[
 \norm{A-B}_*\le\tr A+\tr B
 =2\tr A-\tr(A-B).
\]
Since \(\tr A=p_w\norm{\mu_w}^2\) and
\(\tr(A-B)=\E_G[H_2(G)w(G)]\), this gives
\begin{equation}
 \left\|\E_G[(GG^\top-I_r)w(G)]\right\|_*
 \le2p_w\norm{\mu_w}^2-\E_G[H_2(G)w(G)],
 \qquad H_2(g)\coloneqq \norm g^2-r.
 \label{appx-an:matrix-trace-reduction}
\end{equation}
Thus the matrix calculation applies to any log-concave weight; only
the two scalar estimates below use the logistic form.

Condition on \(V_0\). For each ordered tuple
\(\alpha=(\alpha_1,\ldots,\alpha_k)\), put
\[
 w_\alpha(g)\coloneqq F(V_0+\sigma g)
                 \prod_{a=1}^k\pi_{\alpha_a}(V_0+\sigma g),
 \quad p_\alpha\coloneqq \E_Gw_\alpha(G),
 \quad \mu_\alpha\coloneqq \frac{\E_G[Gw_\alpha(G)]}{p_\alpha}.
\]
Its logarithm is affine minus \((k+1)\log N(V_0+\sigma g)\), so
\(w_\alpha\) is log-concave. Moreover,
\(\bar w(g)\coloneqq \sum_\alpha w_\alpha(g)
=F(V_0+\sigma g)\theta(V_0+\sigma g)^k\).
Summing~\eqref{appx-an:matrix-trace-reduction} and then averaging over
\(V_0\) reduces the claim to bounds for
\[
 \E_{V_0}\sum_\alpha p_\alpha\norm{\mu_\alpha}^2
 \quad\text{and}\quad
 \E_{V_0,G}[H_2(G)F(U)\theta(U)^k].
\]

To bound the means, use the first face of each tuple as a reference.
For a tuple with \(\alpha_1=j\), set
\[
 b_j(g)\coloneqq \frac{e^{q_j(V_0+\sigma g)}}
                  {\{1+e^{q_j(V_0+\sigma g)}\}^2},
 \qquad p_{b_j}\coloneqq \E_Gb_j(G).
\]
This weight depends only on \(u_j\cdot g\), and
\(w_\alpha\le F\pi_j\le b_j\). Lemma~\ref{appx-an:lem:mean}
and \(|u_j\cdot G|\le\norm G_U\) therefore imply
\begin{equation}
 \sum_\alpha p_\alpha\norm{\mu_\alpha}^2
 \le\E_G[\bar w(G)\norm G_U^2]
       +2\sum_\alpha p_\alpha\log(p_{b_{\alpha_1}}/p_\alpha).
 \label{appx-an:entropy-sum}
\end{equation}
The first term has expectation at most \(ChL^{3/2}\): condition on
\(G\), apply~\eqref{appx-an:weight-concentration} using
\(\bar w\le F\theta\), and then use \(\E\norm G_U^2\le CL\).

The logarithmic term can be summed before any individual mass is
bounded. This is essential because an individual \(p_\alpha\) can
be arbitrarily small. Define the total averaged masses
\[
 \bar p\coloneqq \E_{V_0,G}\bar w(G),\qquad
 \bar b\coloneqq m^{k-1}\sum_j\E_{V_0,G}b_j(G).
\]
There are \(m^{k-1}\) tuples with each first face. The log-sum
inequality~\cite[Theorem~2.7.1]{CoverThomas2006} and the joint concavity
of \((x,y)\mapsto x\log(y/x)\)
therefore give
\[
 \E_{V_0}\sum_\alpha p_\alpha
          \log(p_{b_{\alpha_1}}/p_\alpha)
 \le\bar p\log(\bar b/\bar p).
\]
By~\eqref{appx-an:weight-concentration} and its single-face version,
\(\bar p\le Ch\sqrt L\) and \(\bar b\le Cm^kh\sqrt L\).
Writing \(\bar p=h\sqrt L\,x\), with \(0<x\le C\), yields
\[
 \bar p\log(\bar b/\bar p)
 \le h\sqrt L\,x\log(Cm^k/x)
 \le C(k+1)hL^{3/2}.
\]
The zero-mass case follows by continuity. Together with
\eqref{appx-an:entropy-sum}, this proves
\begin{equation}
 \E_{V_0}\sum_\alpha p_\alpha\norm{\mu_\alpha}^2
 \le C(k+1)hL^{3/2}.
 \label{appx-an:tilted-mean-bound}
\end{equation}
The reference masses have the same factor \(h\sqrt L\) as the
weights being bounded, so this factor cancels inside the logarithm.

It remains to bound the trace after summing the weights. A common
threshold shift turns \(F\theta^k\) into a derivative, which we can
interpret geometrically. If \(F_s,\theta_s\) are obtained by increasing
every threshold by \(s\), then
\[
 F\theta^k=-\frac hk\left.\partial_s\theta_s^k\right|_{s=0}.
\]
We turn this derivative into a surface integral by expressing powers
of \(F_s\) as mixtures of polyhedral indicators. For an integer
\(\ell\ge1\), let \(T_\ell\) have the Gamma law with shape \(\ell\)
and rate one. Take \(T_\ell\) and the mean-one exponentials
\(\omega_1,\ldots,\omega_m\) to be mutually independent and
independent of \(V_0,G\). The Gamma Laplace transform gives
\begin{equation}
 F_s(x)^\ell=\Pp\{u_j\cdot x\le
        a_j+s+h\log(\omega_j/T_\ell),\ j\le m\}.
 \label{appx-an:gamma-mixture}
\end{equation}
Indeed, conditional on \(T_\ell\), the probability of all the
inequalities is \(e^{-T_\ell e^{-s/h}T(x)}\). Its average over
\(T_\ell\) is \((1+e^{-s/h}T(x))^{-\ell}=F_s(x)^\ell\).
Expand \(\theta_s^k=\sum_{\ell=0}^k(-1)^\ell\binom k\ell  F_s^\ell\).
The absolute coefficients sum to \(2^k\), and the constant term has
zero derivative. For fixed random thresholds, the coarea
formula~\cite[Section~3.4]{EvansGariepy2015} gives the first variation:
increasing every unit-normal threshold at unit speed differentiates
the Gaussian integral into the boundary integral, with the factor
\(\sigma^{-1}\) after rescaling. This holds for almost every common
shift. Uniform bounds on Gaussian integrals over the finitely many
supporting hyperplanes justify differentiation under the mixture
average. The independent, absolutely continuous shift \(-h\log T_\ell\)
avoids the exceptional levels, including those at which a target
polyhedron is lower dimensional. Lemma~\ref{appx-an:lem:surface}
therefore gives
\begin{align*}
 \left|\E_{V_0,G}[H_2(G)F(U)\theta(U)^k]\right|
 &\le\frac{2^kh}{k\sigma}
       \sup_Q\E_{V_0}|\surf((Q-V_0)/\sigma)|\\
 &\le\frac{C2^kh}{k\sigma}
       (\sigma L^{3/2}+\delta L^2)
 \le C2^khL^{3/2},
\end{align*}
where the last inequality uses \(\sigma\ge\delta\sqrt L\).
Combining this trace bound with
\eqref{appx-an:matrix-trace-reduction} and
\eqref{appx-an:tilted-mean-bound} proves
\eqref{appx-an:matrix-bound}, after absorbing \(k+1\) and \(2^k\)
into \(C_1^{k+1}\).
\end{proof}

\subsection{An averaged Gaussian surface estimate}
\label{appx-an:sec:surface}

The trace term in Lemma~\ref{appx-an:lem:matrix} becomes a Gaussian
surface integral over a randomly translated polyhedron. We now show
that averaging over the concentrated translation gains the factor
\(\sigma\) needed to cancel the \(\sigma^{-1}\) in that argument.
We retain the notation
\[
 H_2(g)=\norm g^2-r,\qquad
 \surf(P)=\int_{\partial P}H_2(g)\varphi_r(g)\dd\mathcal H^{r-1}(g),
\]
where \(\varphi_r\) is the standard Gaussian density and
\(\mathcal H^{r-1}\) is Hausdorff measure, normalized to agree with
Euclidean surface measure on each facet~\cite[Chapter~2]{EvansGariepy2015}.
This is a signed integral: the absolute value is taken only after
integration over the boundary. Each facet is counted once. We define
it for proper, full-dimensional polyhedra, meaning those with nonempty
interior that are not the whole space.
We set \(\surf(P)\coloneqq 0\) for the empty set and the whole space. As in
Subsection~\ref{appx-an:sec:coords}, \(P^s\) means that all thresholds
in the defining system of inequalities with unit face normals are increased by
\(s\); it does not denote a Euclidean neighborhood.

\begin{lemma}[Averaged scalar surface]\label{appx-an:lem:surface}
Let \(V_0\) be a concentrated base, and let \(0<\sigma\le1\). Uniformly
over target polyhedra \(Q\) with face normals among \(u_1,\ldots,u_m\), empty
or full dimensional,
\begin{equation}
 \E\left|\surf\bigl((Q-V_0)/\sigma\bigr)\right|
 \le C\bigl(\sigma L^{3/2}+\delta L^2\bigr).
 \label{appx-an:surface-bound}
\end{equation}
In particular, if \(\sigma\ge\delta\sqrt L\), the right side is at most
\(C\sigma L^{3/2}\).
\end{lemma}

\begin{proof}
We first bound the surface integral of a fixed polyhedron by an integral
over short outward tubes, and then average over \(V_0\). The
facet--ridge identity follows from the classical Gauss--Green
theorem~\cite[Chapter~5]{EvansGariepy2015}. The outward facet construction
is due to Nazarov~\cite{Nazarov2003}; see the detailed proof
in~\cite[Lemmas~1--2]{CCK17N}, which explicitly traces this argument
to Nazarov. Below we derive the variable tube widths, ridge-sector
integrals, and weighted estimates used here. Choosing the width at each
boundary point reduces the subsequent averaging to two Gaussian moments.
For a related polyhedral application, see~\cite[Proposition~4]{GaoZhang}.

The facet contributions depend on one normal coordinate; the ridge
contributions depend on the plane spanned by two normals. We therefore
control only these projections, rather than the full Gaussian norm. For linearly independent \(u_j,u_k\), let \(\Pi_{jk}\)
be orthogonal projection onto \(\operatorname{span}(u_j,u_k)\), and put
\begin{equation}
 \mathcal M(g)\coloneqq \max\left\{\max_j|u_j\cdot g|,
       \max_{j,k:\,u_j,u_k\text{ independent}}\norm{\Pi_{jk}g}\right\},
 \qquad R(g)\coloneqq 1+\mathcal M(g).
 \label{appx-an:projection-radius}
\end{equation}
An empty second maximum is omitted. The function \(R\) is
\(1\)-Lipschitz, since each projection is a contraction. For
\(G\sim N(0,I_r)\), the maximum involves at most \(m+m^2\)
Gaussian projections of dimension at most two. Their tail bounds and
a union bound give
\[
 \Pp\{\mathcal M(G)>s\}\le\min\{1,Cm^2e^{-c_*s^2}\}.
\]
Integrating this bound, with the integral split at a constant times
\(\sqrt{\log(em)}\), and using \(m=2p\), yields
\begin{equation}
 \E R(G)^2\le CL,\qquad \E R(G)^3\le CL^{3/2}.
 \label{appx-an:projection-moments}
\end{equation}

For every proper, full-dimensional polyhedron \(P\) with face normals among
the \(u_j\), we will prove
\begin{equation}
 |\surf(P)|\le C\E_G\left[R(G)^3
       \ind\{G\in P^{1/R(G)}\setminus P\}\right].
 \label{appx-an:adaptive-tubes}
\end{equation}
The right side integrates over an outer layer whose width is smaller
where the normal projections are larger. This keeps the Gaussian
density comparable along each tube without using the full norm
\(\norm G\), whose moments would introduce the rank \(r\).

To prove~\eqref{appx-an:adaptive-tubes}, write the defining thresholds
of \(P\) as \(a'_j\). For the facet calculation, retain only the
inequalities that describe distinct actual facets; the outer layer
still uses the full defining system. We use the standard face geometry
of polyhedra; see~\cite[\S\S18--19]{Rockafellar1970}
and~\cite[Chapter~2]{Ziegler1995}. Write \(F_j\) for a retained
facet in \(u_j\cdot x=a'_j\), and let
\(s_j\coloneqq \int_{F_j}\varphi_r\dd\mathcal H^{r-1}\).
A ridge \(E_{jk}\) is a face of dimension \(r-2\) shared by exactly
two facets, \(F_j\) and \(F_k\); see~\cite[Chapter~2]{Ziegler1995}
for this standard face-incidence property. Let \(\alpha_{jk}\in(0,\pi)\) be the
angle between their face normals, and set
\[
 \rho_{jk}\coloneqq \int_{E_{jk}}\varphi_r\dd\mathcal H^{r-2},\qquad
 z_{jk}\coloneqq \Pi_{jk}x\quad(x\in E_{jk}).
\]
The vector \(z_{jk}\) is constant on the ridge: it is determined by
\(u_j\cdot z_{jk}=a'_j\) and \(u_k\cdot z_{jk}=a'_k\).
Applying the divergence theorem within each facet expresses the
surface integral through facets and their shared ridges:
\begin{equation}
 \surf(P)=\sum_j\bigl((a'_j)^2-1\bigr)s_j
       -\sum_{\{j,k\}}(d_{k\mid j}+d_{j\mid k})\rho_{jk},
 \label{appx-an:surface-identity}
\end{equation}
where the second sum runs over the ridges and
\[
 d_{k\mid j}\coloneqq \frac{a'_k-(u_j\cdot u_k)a'_j}{\sin\alpha_{jk}}.
\]
Within the hyperplane of facet \(j\), \(d_{k\mid j}\) is the signed
distance from \(a'_ju_j\) to the hyperplane containing the ridge,
measured in the outward unit normal direction of that ridge.
Writing \(x=a'_ju_j+y\) on facet \(j\), with \(y\perp u_j\),
splits \(H_2(x)\) into \((a'_j)^2-1\) and \(\norm y^2-(r-1)\).
The first term gives the facet contribution above. For the second,
\[
 (\norm y^2-(r-1))\varphi_{r-1}(y)
 =-\operatorname{div}_{u_j^\perp}\{y\varphi_{r-1}(y)\}.
\]
The classical Gauss--Green divergence theorem within the
facet~\cite[Chapter~5]{EvansGariepy2015} gives
\(-d_{k\mid j}\rho_{jk}\) on its ridge with facet \(k\), because
the outward unit normal to that ridge within the facet is
\((u_k-(u_j\cdot u_k)u_j)/\sin\alpha_{jk}\).
For an unbounded facet, apply the theorem after multiplying the
vector field by a smooth cutoff that is one on the ball of radius
\(M\), vanishes outside the ball of radius \(2M\), and has gradient
bounded by \(C/M\). The cutoff error tends to zero by Gaussian decay.
Dominated convergence for the facet and ridge integrals then gives
the same identity. This cutoff argument is our specialization of
Gauss--Green to the polynomial Gaussian integrands used here.
Summing over facets proves~\eqref{appx-an:surface-identity}.
For \(r=1\), there are no ridges and the identity is a sum over endpoints.

We pair the two contributions of each ridge before bounding them.
The elementary identity
\[
 d_{k\mid j}+d_{j\mid k}
 =2\sin(\alpha_{jk}/2)
       \frac{u_j+u_k}{\norm{u_j+u_k}}\cdot z_{jk}
\]
gives
\begin{equation}
 |d_{k\mid j}+d_{j\mid k}|\le\alpha_{jk}\norm{z_{jk}}.
 \label{appx-an:paired-ridge-bound}
\end{equation}
Thus the inverse sine in the separate terms disappears. This pairing
allows the estimate to remain uniform when facets are nearly parallel.

We now represent the absolute facet and paired-ridge contributions
by Gaussian integrals just outside \(P\). At a point \(x\) in the
relative interior of a facet or ridge, choose the tube width
\[
 \eta(x)\coloneqq \frac1{4R(x)}.
\]
Over a facet, use displacements \(y=s u_j\), with \(0<s<\eta(x)\).
Over a ridge, use \(0<\norm y<\eta(x)\) in the outward normal cone,
the sector of angle \(\alpha_{jk}\) spanned by \(u_j,u_k\), by the
normal-cone formula for polyhedra~\cite[Theorem~6.46]{RockafellarWets1998}.
By the projection--normal-cone characterization
\cite[Proposition~6.17]{RockafellarWets1998}, these normal displacements
satisfy \(x=\operatorname{proj}_P(x+y)\). Thus \(x\) is the nearest point
in \(P\) to \(x+y\), which will make the tubes over distinct faces
disjoint. In both cases,
\[
 |x\cdot y|\le\mathcal M(x)\norm y\le\tfrac14,\qquad
 \frac{\varphi_r(x+y)}{\varphi_r(x)}
       =e^{-x\cdot y-\norm y^2/2}\ge e^{-1/4-1/32}.
\]
For a ridge, the first bound uses \(x\cdot y=z_{jk}\cdot y\).
Since \(R\) is \(1\)-Lipschitz and \(R\ge1\), we also have
\(3R(x)/4\le R(x+y)\le5R(x)/4\).
Integrating the weight \(R^3\varphi_r\) along a facet's normal
interval therefore gives
\[
 \int_0^{\eta(x)} R(x+s u_j)^3\varphi_r(x+s u_j)\dd s
 \ge c_*R(x)^2\varphi_r(x)
 \ge c_*\{1+(a'_j)^2\}\varphi_r(x).
\]
The corresponding integral over a ridge's normal sector gives
\[
 \int_{\substack{y\in N_P(x)\\0<\norm y<\eta(x)}}
        R(x+y)^3\varphi_r(x+y)\dd y
 \ge c_*\alpha_{jk}R(x)\varphi_r(x)
 \ge c_*\alpha_{jk}\norm{z_{jk}}\varphi_r(x),
\]
where \(N_P(x)\) is the normal cone and \(\dd y\) is area in its
two-dimensional plane. The first integral contains one power of
\(\eta(x)\); the second contains two. This explains the common weight
\(R^3\) in~\eqref{appx-an:adaptive-tubes}.

The tubes over distinct faces are disjoint up to null sets, by
uniqueness of the nearest point in a closed convex
set~\cite[Example~2.25]{RockafellarWets1998}. Within each
face, the decomposition into tangential and normal coordinates has
Jacobian one. Fubini's theorem~\cite[Section~2.5]{Folland1999}
therefore applies even though
\(\eta(x)\) varies along the face. Every tube point \(g=x+y\) lies
outside \(P\), and, for every inequality in its original description,
\[
 u_j\cdot g-a'_j\le\norm y<\frac1{4R(x)}\le\frac1{R(g)}.
\]
Thus all tubes lie in \(\{g:g\in P^{1/R(g)}\setminus P\}\), including
when the defining system has redundant inequalities. Integrating the
two tube estimates over the facets and ridges, and applying
\eqref{appx-an:surface-identity} and
\eqref{appx-an:paired-ridge-bound}, proves
\eqref{appx-an:adaptive-tubes}.

We now apply this deterministic bound to \(P\coloneqq (Q-V_0)/\sigma\),
with \(G\) independent of \(V_0\). Empty targets and the whole space
contribute zero, so assume \(Q\) is proper and full dimensional.
Rescaling the boundary event gives
\[
 G\in P^{1/R(G)}\setminus P
 \quad\Longleftrightarrow\quad
 V_0+\sigma G\in Q^{\sigma/R(G)}\setminus Q.
\]
We average over \(V_0\) with \(G\) fixed, so the width
\(\sigma/R(G)\) and the threshold translations
\(-\sigma u_j\cdot G\) are deterministic. Enlarge the outer layer
to the symmetric layer in~\eqref{appx-an:concentration}. Concentration then
bounds its conditional probability by
\(C\{\sigma/R(G)+\delta\}\sqrt L\). Tonelli's
theorem~\cite[Section~2.5]{Folland1999} and
\eqref{appx-an:projection-moments} yield
\begin{align*}
 \E_{V_0}|\surf((Q-V_0)/\sigma)|
 &\le C\E_G\left[R(G)^3
       \Pp_{V_0}\{V_0+\sigma G\in Q^{\sigma/R(G)}\setminus Q\}\right]\\
 &\le C\sqrt L\{\sigma\E R(G)^2+\delta\E R(G)^3\}\\
 &\le C\{\sigma L^{3/2}+\delta L^2\}.
\end{align*}
This is~\eqref{appx-an:surface-bound}.
\end{proof}

% Source: concentration_proof
\section{\texorpdfstring{Proof of the anti-concentration inequality of Lemma~\ref{lem:concentration}}{Proof of the anti-concentration inequality of Lemma 1}}\label{appx-app:concentration}
\label{appx-conc:conditioning}

We represent each summand as a random choice between two independent
copies. Conditional on the copies, its \(m=2p\) face components
form a translated Rademacher sum. We then apply the anti-concentration
inequality of O'Donnell, Servedio, and Tan~\cite{OST}, in the form
recorded in~\cite[Lemma~J.4]{CCKK2022}. This proves the bounded-sum
estimate of~\cite[Lemma~2.1]{AL} for the face inequalities used here.

Use the notation \(V=\sum_iX_i\) of Lemma~\ref{lem:concentration}.
The bound \(\norm{X_i}_U\le\eta\) controls each face component,
and their total variances equal \(v\). We first prove the result
at variance one without the independent shift; scaling and conditioning
will then give the full assertion. Let \(X_i'\) be independent copies of the
summands, and let \(\varepsilon_i\) be independent uniform signs,
independent of all these vectors. Put
\[
 \bar X_i\coloneqq\frac{X_i+X_i'}2,\qquad
 \Delta_i\coloneqq\frac{X_i-X_i'}2,\qquad
 \mathcal G\coloneqq\sigma(X_i,X_i':1\le i\le n).
\]
Each \(\bar X_i+\varepsilon_i\Delta_i\) selects one of the two
copies, so
\begin{equation}
 V\stackrel d=\widetilde V\coloneqq
 \sum_i\bar X_i+\sum_i\varepsilon_i\Delta_i.
 \label{appx-eq:pair-representation}
\end{equation}
Conditional on \(\mathcal G\), the first sum is fixed,
\(\norm{\Delta_i}_U\le\eta\), and the variance in face direction
\(\alpha\) of the second sum is
\(V_\alpha\coloneqq\sum_i(u_\alpha\cdot\Delta_i)^2\).
To apply the Rademacher inequality, these conditional variances
must stay away from zero simultaneously. Independence and centering give
\[
 \E V_\alpha=\frac12,\qquad
 0\le(u_\alpha\cdot\Delta_i)^2\le\eta^2,\qquad
 \sum_i\Var\bigl((u_\alpha\cdot\Delta_i)^2\bigr)
 \le\frac{\eta^2}{2}.
\]
Bernstein's inequality~\cite[Chapter~2]{BLM2013}, applied to the
centered squared components, and a union bound imply
\begin{equation}
 \Pp(\mathcal E^c)\le m e^{-c_1/\eta^2},\qquad
 \mathcal E\coloneqq\{\min_{\alpha\le m}V_\alpha\ge1/4\},
 \label{appx-eq:conditional-variance}
\end{equation}
for an absolute constant \(c_1>0\).

On \(\mathcal E\), apply~\cite[Lemma~J.4]{CCKK2022} conditionally
with coefficient vectors
\(\sqrt n\,(u_\alpha\cdot\Delta_i)_{\alpha\le m}\).
Its normalized sum is exactly the random part of \(\widetilde V\);
its variance lower bound is \(1/4\), and its coefficient envelope
divided by \(\sqrt n\) is at most \(\eta\). Absorb the fixed
components of \(\sum_i\bar X_i\) into the thresholds. The outer and inner inequalities have thresholds
separated by \(2s\), so, since \(\log(em)\le CL\),
\[
 \Pp(\widetilde V\in A_a^s\setminus A_a^{-s}\mid\mathcal G)
 \le C(s+\eta)\sqrt L\qquad\text{on }\mathcal E.
\]
The result is uniform in all finite thresholds, including inconsistent
ones. Averaging and using~\eqref{appx-eq:conditional-variance} gives
\begin{equation}
 \Pp(V\in A_a^s\setminus A_a^{-s})
 \le C(s+\eta)\sqrt L+m e^{-c_1/\eta^2}.
 \label{appx-eq:concentration-pairs}
\end{equation}
The exceptional probability is harmless in the regime where the
claimed estimate is informative. If \(\eta\sqrt L\le c_2\),
for a sufficiently small absolute
\(c_2>0\), then \(\log m\le c_1/(2\eta^2)\). The last term is
at most \(e^{-c_1/(2\eta^2)}\le C\eta\le C\eta\sqrt L\).
Otherwise the claimed estimate follows from the bound by one.
This proves the bounded-sum assertion for \(v=1\).
For general \(v>0\), divide the summands and thresholds by
\(\sqrt v\); the width and summand bound become \(s/\sqrt v\)
and \(\eta/\sqrt v\). Conditioning permits the independent shift
\(R\), which translates the thresholds by \(-u_\alpha\cdot R\).

For a centered Gaussian vector \(V\), apply Nazarov's inequality
in~\cite[Lemma~A.1]{CCK2017}, also recorded
in~\cite[Lemma~J.3]{CCKK2022}, to
\(\bigl(u_\alpha\cdot V/\sqrt v\bigr)_{\alpha\le m}\). It gives
\[
 \Pp(V\in A_a^s\setminus A_a^{-s})\le Cs\sqrt{L/v}.
\]
Only positive marginal variances are required; the face components
may have singular covariance. Condition on the independent shift
and take the minimum with one to complete the proof.

% Source: subexponential_extension_proof
\section{\texorpdfstring{Sub-exponential transfer: proof of Corollary~\ref{cor:subexponential}}{Sub-exponential transfer: proof of Corollary 3}}\label{app:subexponential}
\begin{proof}[Proof of Corollary~\ref{cor:subexponential}]
We first prove the rectangle bounds by clipping each coordinate,
applying the bounded theorem after marginal standardization, and
transferring the result back to \(W_n\). This requires control of both
the discarded sum and the Gaussian covariance change. The final step
applies these bounds to the standardized face projections of a polyhedron.
Suppressing the subscript \(n\) on the parameters of the corollary,
write \(\ell=\log(en)\) and \(b=K\ell/\sqrt n\).

The norm condition implies \(\E e^{|X_{ij}|/K}\le2\). Indeed,
every scale larger than \(K\) is admissible, and monotone convergence
allows those scales to decrease to \(K>0\). We use this exponential
moment directly in the clipping argument.

The unrestricted bound follows from the bound by one, after adjusting its absolute
constant, when \(b\ge1/16\), so assume \(b<1/16\).

We choose a clipping level depending on \(\log n\); the dependence
on \(p\) will enter when we bound the residual sums. With \(T\coloneqq4K\ell\), define coordinatewise
\[
 C_T(x)\coloneqq \operatorname{sgn}(x)\min\{|x|,T\},\qquad
 Y_i\coloneqq C_T(X_i)-\E C_T(X_i),\qquad D_i\coloneqq X_i-C_T(X_i).
\]
Here \(Y_i\) is the centered clipped summand and \(D_i\) is the discarded
part. Their normalized sums, marked by the subscripts \(T\) and \(D\), satisfy
\[
 W_T\coloneqq n^{-1/2}\sum_iY_i,\qquad
 W_D\coloneqq n^{-1/2}\sum_i(D_i-\E D_i),\qquad W_n=W_T+W_D.
\]
The exponential-moment assumption gives, for \(s\ge0\),
\[
 \Pp(|D_{ij}|>s)\le2e^{-4\ell-s/K},\qquad
 \E|D_{ij}|^k\le2k!K^ke^{-4\ell}\quad(k\ge1).
\]
In particular, independence of the summands implies
\begin{equation}
 \max_j\Var(W_{D,j})\le4K^2e^{-4\ell}.
 \label{tail:residual-moments}
\end{equation}
The pointwise bound \(e^{|D_{ij}|/K}\le1+e^{-4\ell}e^{|X_{ij}|/K}\)
gives \(\E e^{|D_{ij}|/K}\le1+2e^{-4\ell}\). Together with
\(\E|D_{ij}|\le2Ke^{-4\ell}\), independence and the triangle inequality
therefore yield
\begin{align*}
 \E\exp\!\left(\frac{\sqrt n\,|W_{D,j}|}{K}\right)
 &\le\exp\!\left(\frac1K\sum_i\E|D_{ij}|\right)
          \prod_i\E e^{|D_{ij}|/K}\\
 &\le e^{2ne^{-4\ell}}(1+2e^{-4\ell})^n
 \le e^{4ne^{-4\ell}}.
\end{align*}
For \(\bar L\coloneqq \log(4np)\), Markov's inequality gives
\[
 \Pp\!\left(\sqrt n\,|W_{D,j}|\ge K[\bar L+4ne^{-4\ell}]\right)\le e^{-\bar L}.
\]
Since \(n,p\ge1\),
\[
 4ne^{-4\ell}=4e^{-4}n^{-3}\le\frac14\le\log4\le \bar L.
\]
A coordinate union bound now yields
\begin{equation}
 \Pp(\|W_D\|_\infty>\varrho_1)\le\frac1{4n}\le\frac1{2n},\qquad
 \varrho_1\coloneqq \frac{2K\bar L}{\sqrt n}.
 \label{tail:sum-coupling}
\end{equation}
Since the union bound is applied to the residual sums, the clipping
level can depend only on \(K\) and \(\log n\).

Before applying the bounded theorem, we verify that clipping has
changed each standard deviation only slightly. Set
\(\tau\coloneqq 2Ke^{-2\ell}\), the residual standard-deviation
bound from~\eqref{tail:residual-moments}, and write \(\operatorname{sd}(W_{T,j})\) for the
standard deviation of the \(j\)th coordinate of~\(W_T\).
Since each coordinate of \(W_n\) has variance one, the triangle
inequality in \(L^2\) gives
\[
 |\operatorname{sd}(W_{T,j})-1|\le\tau,
 \qquad \frac\tau b=\frac{2\sqrt n\,e^{-2\ell}}\ell
 \le\frac2{e^2}.
\]
Hence \(\operatorname{sd}(W_{T,j})\ge1/2\). Since \(\|Y_i\|_\infty\le8K\ell\), the
independent centered summands of the marginally standardized \(W_T\)
have envelope at most
\[
 \frac{8K\ell}{\sqrt n\min_j\operatorname{sd}(W_{T,j})}\le16b.
\]
Let \(Z_T\) be Gaussian with covariance \(\Cov(W_T)\).
Positive diagonal scaling preserves rectangle distance. Applying
Theorem~\ref{thm:unrestricted} with envelope \(16b\) gives
\begin{equation}
 \rh(W_T,Z_T)\le C\min\left\{1,
 bL^{3/2}\sqrt{1+\log_+\frac1{b\sqrt L}}\right\}.
 \label{tail:bounded-comparison}
\end{equation}
The factor 16 is absorbed into the absolute constant, since increasing
the envelope decreases the logarithmic factor.

The comparison just obtained has the clipped covariance. To recover
the original target, we reproduce the joint covariance of the clipped
and residual sums with a Gaussian pair. Construct a jointly centered Gaussian vector \((Z_T,Z_D)\) with the
block covariance of \((W_T,W_D)\). This covariance is positive
semidefinite, even when singular. Set \(Z_n\coloneqq Z_T+Z_D\) on this coupling.
Then \(Z_n\) has exactly the covariance of \(W_n\), while
\(\max_j\Var(Z_{D,j})\le\tau^2\). Gaussian tails therefore give
\begin{equation}
 \Pp(\|Z_D\|_\infty>\varrho_2)\le\frac1{2n},\qquad
 \varrho_2\coloneqq \sqrt8\,Ke^{-2\ell}\sqrt{\bar L}.
 \label{tail:Gaussian-coupling}
\end{equation}
The Gaussian components need not be independent.

For a bounded rectangle \(A=\prod_j[\ell_j,u_j]\), let
\(A^s\coloneqq\prod_j[\ell_j-s,u_j+s]\); inconsistent bounds
give the empty set. On the coupling events, membership is sandwiched
between these expanded and contracted rectangles. Thus
\begin{align*}
 \Pp(W_n\in A)
 &\le\Pp(W_T\in A^{\varrho_1})+\frac1{2n}\\
 &\le\Pp(Z_T\in A^{\varrho_1})+\rh(W_T,Z_T)+\frac1{2n}\\
 &\le\Pp(Z_n\in A^{\varrho_1+\varrho_2})+\rh(W_T,Z_T)+\frac1n.
\end{align*}
Contracting the thresholds gives the reverse bound
\(\Pp(W_n\in A)\ge\Pp(Z_n\in A^{-(\varrho_1+\varrho_2)})
-\rh(W_T,Z_T)-1/n\). Applying Gaussian rectangle concentration at
the unit-marginal target \(Z_n\), obtained from Lemma~\ref{lem:concentration}
using its support representation, yields
\begin{equation}
 \rh(W_n,Z_n)\le\rh(W_T,Z_T)
       +C(\varrho_1+\varrho_2)\sqrt L+\frac1n.
 \label{tail:transfer}
\end{equation}
The argument uses only the probabilities of the coupling events;
the components within each coupling may be dependent. Unbounded
rectangles follow by truncation and continuity from below.
The transfer error has order \(bL^{3/2}\): since
\(\bar L\le2\ell L\) and \(K\ge1\),
\[
 (\varrho_1+\varrho_2)\sqrt L+\frac1n
 \le CK\left[e^{-2\ell}\sqrt{L\bar L}
                     +\frac{\bar L\sqrt L}{\sqrt n}\right]+\frac1n
 \le CbL^{3/2}.
\]
Combine this with \eqref{tail:bounded-comparison} and the bound by one
to prove \eqref{eq:subexponential}.

For the sharper bound we must also check the correlations after
clipping and standardization. To prove~\eqref{eq:subexponential-gap},
it suffices to consider
\(b<\min\{1/16,c_0/4\}\); otherwise its right-hand side is at
least one after increasing \(C_{c_0}\).
Under the Gaussian coupling, both \(Z_n\) and the standardized clipped
vector \((Z_{T,j}/\operatorname{sd}(W_{T,j}))_{j\le p}\) have unit coordinate variances, and
\[
 \|Z_{T,j}/\operatorname{sd}(W_{T,j})-Z_{n,j}\|_{L^2}\le2\tau.
\]
For this calculation, abbreviate the standardized clipped Gaussian
vector by \(A\) and the target \(Z_n\) by \(B\). Both are centered
with unit coordinate variances, so Cauchy--Schwarz gives
\[
 \begin{aligned}
 |\E A_jA_k-\E B_jB_k|
 &\le\|A_j-B_j\|_{L^2}\|A_k\|_{L^2}
       +\|B_j\|_{L^2}\|A_k-B_k\|_{L^2}\\
 &\le4\tau.
 \end{aligned}
\]
Thus the covariance matrices differ entrywise by at most \(4\tau\). Since
\(4\tau\le(8/e^2)b<c_0/2\), the standardized clipped covariance
has correlation gap at least \(c_0/2\).
Theorem~\ref{thm:optimized} applied with envelope \(16b\) now gives
\[
 \rh(W_T,Z_T)\le C_{c_0}\min\left\{1,
 b\left[L^{3/2}+L\sqrt{1+\log_+\frac1{bL}}\right]\right\}.
\]
The transfer error in \eqref{tail:transfer} is bounded by the
\(CbL^{3/2}\) term. This proves \eqref{eq:subexponential-gap}.

The rectangle bounds now apply to the standardized face projections,
by the same reduction as in Subsection~\ref{app:polyhedral-transfer}.
We only need to verify the tail condition. Fix a polyhedron
\(P\in\mathcal P_m(J;\Sigma)\) with a defining representation
having \(q\le m\) inequalities, and put
\[
 \widetilde X_i\coloneqq
 \left(\frac{\nu_\alpha^\top X_i}{\sigma_\alpha}\right)_{\alpha\le q},
 \qquad \widetilde W_n\coloneqq n^{-1/2}\sum_i\widetilde X_i,
 \qquad \widetilde Z_n\coloneqq
 \left(\frac{\nu_\alpha^\top Z_n}{\sigma_\alpha}\right)_{\alpha\le q}.
\]
The transformed observations are independent and centered, and
\(n^{-1}\sum_i\E\widetilde X_{i\alpha}^2=1\).
The triangle inequality for the Orlicz norm, which follows from
H\"older's inequality applied to its defining exponential moment, gives
\[
 \|\widetilde X_{i\alpha}\|_{\psi_1}
 \le\sum_{j=1}^p\frac{|\nu_{\alpha j}|}{\sigma_\alpha}
                    \|X_{ij}\|_{\psi_1}
 \le\frac{K\norm{\nu_\alpha}_1}{\sigma_\alpha}
 \le JK\le K_{\mathcal P}.
\]
This step requires no independence among the coordinates of \(X_i\).
The vector \(\widetilde Z_n\) is Gaussian with the covariance of
\(\widetilde W_n\). For the rectangle
\(B\coloneqq\prod_{\alpha=1}^q(-\infty,a_\alpha/\sigma_\alpha]\),
\[
 \Pp(W_n\in P)=\Pp(\widetilde W_n\in B),\qquad
 \Pp(Z_n\in P)=\Pp(\widetilde Z_n\in B).
\]
Apply the rectangle bounds just proved, with tail scale
\(K_{\mathcal P}\) and dimension \(q\). Both right-hand sides are
nondecreasing in the dimension parameter \(\log(eq)\ge1\), so it
can be replaced by \(L_m\). In the fixed-gap case, the correlation
matrix has entries
\(\nu_\alpha^\top\Sigma\nu_\beta/(\sigma_\alpha\sigma_\beta)\),
which verifies the required hypothesis for the stated subclass.
Taking the supremum over \(P\) proves the polyhedral assertions.
The empty set and the whole space have zero error; zero-variance
inequalities are deterministic and are treated as in
Subsection~\ref{app:polyhedral-transfer}.
\end{proof}

The same proof does not automatically give a third-moment-cancellation
refinement. Odd coordinatewise clipping preserves central symmetry
of each summand, but it need not preserve aggregate cancellation
across different summands. Even under symmetry, the transfer error
in~\eqref{tail:transfer} still contributes \(bL^{3/2}\).

To compare the cost of allowing arbitrary covariance, consider the
companion paper~\cite[Corollary~1]{Local}. It assumes the same
coordinatewise \(\psi_1\) condition, an averaged fourth-moment bound
\(\max_jn^{-1}\sum_i\E|X_{ij}|^4\le B_4^2\), and uniformly positive
eigenvalues for principal covariance submatrices of order at most four.
When these principal-submatrix eigenvalues are bounded below by \(a>0\), its bound is
\[
 \frac{C_a}{\sqrt n}\left[B_4L^{3/2}
                   +K\log(enp)\sqrt L\right],
\]
capped at one. For fixed \(K,B_4,a\) and \(p\) growing polynomially in~\(n\),
this is \(O(n^{-1/2}(\log n)^{3/2})\), whereas
\eqref{eq:subexponential} is \(O(n^{-1/2}(\log n)^3)\).
The latter requires neither a fourth-moment hypothesis in addition to
the tail bound nor a local covariance assumption. If \(K\) varies,
the elementary fourth-moment consequence is \(B_4\le C K^2\);
the tail bound alone does not justify taking \(B_4=K\).

\section{\texorpdfstring{The symmetric lower bound: proof of Theorem~\ref{opt:symmetric-lower}}{The symmetric lower bound: proof of Theorem 6}}
\label{app:symmetric-lower}

We prove Theorem~\ref{opt:symmetric-lower} by constructing two maxima
with leading components close in law and small independent shifts.
The Gaussian shift has a larger mean; the density of the leading
maximum converts this mean difference into a probability difference.
We first bound the approximation and Taylor errors, including when the
dimension grows slowly. Throughout,
\(\Phi\) and \(\varphi\) denote the standard normal distribution
function and density, and \(\bar\Phi\coloneqq1-\Phi\).

\subsection{Two estimates for maxima}

We need a uniform approximation for a maximum of independent
Rademacher sums, keeping the lattice error explicit. Let
\(A_n\coloneqq n^{-1/2}\sum_{i=1}^n\eta_i\), where the
\(\eta_i\) are independent Rademacher signs, and let
\(A_{n,1},\ldots,A_{n,m}\) be independent copies. Put
\(\Lambda_m\coloneqq\log(em)\). There are absolute constants
\(c_0,C>0\) such that, for large \(n\) and
\(\Lambda_m\le c_0\sqrt n\), the Rademacher-maximum approximation is
\begin{equation}
 \sup_x\left|\Pp\{\max_{j\le m}A_{n,j}\le x\}-\Phi(x)^m\right|
 \le C\left(\frac{\sqrt{\Lambda_m}}{\sqrt n}
                 +\frac{\Lambda_m^2}{n}\right).
 \label{sym:eq:rademacher-max}
\end{equation}
The first term retains the effect of the lattice. A scalar absolute
error would accumulate over \(m\) coordinates, so we instead control
the relative error in the upper tail. Exponential tilting gives,
for a sufficiently small
absolute \(a>0\) and \(0\le x\le an^{1/4}\),
\[
 \left|\log\frac{\Pp(A_n>x)}{\bar\Phi(x)}\right|
 \le C\epsilon_n(x),\qquad
 \epsilon_n(x)\coloneqq\frac{1+x}{\sqrt n}+\frac{x^4}{n}.
\]
The tilt makes the threshold \(x\) the new mean, placing the tail
calculation in a region where the scalar central limit theorem applies.
The log moment-generating function is
\(K(\tau)\coloneqq n\log\cosh(\tau/\sqrt n)\). The solution to
\(K'(\tau)=x\) is \(\tau=\sqrt n\operatorname{arctanh}(x/\sqrt n)\).
Set \(v^2\coloneqq K''(\tau)=1-x^2/n\) and \(z\coloneqq\tau v\).
Taylor expansion gives
\[
 K(\tau)-\tau x=-x^2/2+O(x^4/n),
 \qquad z=x+O(x^3/n).
\]
Under the tilted law
\(d\Pp_\tau=e^{\tau A_n-K(\tau)}d\Pp\), the signs remain
independent, and \(T\coloneqq(A_n-x)/v\) has mean zero and variance
one. For \(Y_i\coloneqq(\eta_i-x/\sqrt n)/\sqrt n\),
\[
 \sum_i\E_\tau Y_i^2=v^2,
 \qquad \sum_i\E_\tau|Y_i|^3\le2v^2/\sqrt n.
\]
The scalar Berry--Esseen inequality~\cite{Berry1941} bounds the
distributional error by \(Cv^{-3}\sum_i\E_\tau|Y_i|^3\).
Since \(v\) is bounded away from zero, this gives
\(\sup_y|\Pp_\tau(T\le y)-\Phi(y)|\le C/\sqrt n\).
For \(z\ge0\), integration by parts in the tail yields
\[
 \E_\tau[e^{-zT}\ind\{T>0\}]
 =\Pp_\tau(T>0)-z\int_0^\infty e^{-zy}\Pp_\tau(T>y)\dd y
 =J(z)+O(n^{-1/2}),
\]
where \(J(z)\coloneqq\int_0^\infty e^{-zy}\varphi(y)\dd y\);
the integral multiplied by \(z\) is zero when \(z=0\).
The strict inequalities in this identity account for a possible atom
at zero. Elementary integration gives
\(J(z)\asymp(1+z)^{-1}\) and
\(|(\log J)'(z)|\le C/(1+z)\). Thus the error relative to \(J(z)\)
is \(O((1+x)/\sqrt n)\), and
\(\log J(z)-\log J(x)=O(x^3/[n(1+x)])\). Combining these facts with
\[
 \Pp(A_n>x)=e^{K(\tau)-\tau x}
                \E_\tau[e^{-zT}\ind\{T>0\}],
 \qquad \bar\Phi(x)=e^{-x^2/2}J(x),
\]
proves the relative-tail estimate. Taking \(a\) small also gives
uniform comparability of the two tails and
\(|\Pp(A_n>x)/\bar\Phi(x)-1|\le C\epsilon_n(x)\).

Independence expresses the maximum distribution as a power of the
scalar distribution function. We consider three threshold ranges.
For \(x\le1\),
Berry--Esseen and a fixed upper bound \(r_0<1\) on the two scalar
distribution functions give an error at most
\(Cm r_0^{m-1}/\sqrt n\le C/\sqrt n\). For
\(1<x\le an^{1/4}\), put \(p\coloneqq\Pp(A_n>x)\),
\(v_0\coloneqq\bar\Phi(x)\), and \(w\coloneqq mv_0\).
Tail comparability and the mean value theorem give
\[
 |(1-p)^m-(1-v_0)^m|\le C\epsilon_n(x)w e^{-c_*w}.
\]
If \(x\le2\sqrt{\Lambda_m}\), this is bounded by the asserted
maximum error. Otherwise \(w\le e^{-x^2/4}\), so the error is
at most \(C\epsilon_n(x)e^{-x^2/4}\le C/\sqrt n\).
Finally, choose \(c_0<a^2/4\). For \(x>an^{1/4}\), the
sub-Gaussian tails bound the error by
\(2m e^{-x^2/2}\le2e^{-a^2\sqrt n/4}=O(n^{-1/2})\).
This proves~\eqref{sym:eq:rademacher-max}.

The shift comparison requires a threshold where the Gaussian-maximum
density is of order \(\sqrt{\log m}\), and a global bound on its
derivative. Let
\(F_m(x)\coloneqq\Phi(x)^m\). For \(m\ge2\), choose \(u_m\) by
\(m\bar\Phi(u_m)=1\), and put \(H_*\coloneqq\max_{k\le m}H_k\)
for independent standard normals. The Gaussian-maximum estimates are,
as \(m\to\infty\),
\begin{equation}
 u_m\sim\sqrt{2\log m},\qquad
 F_m'(u_m)\sim\frac{\sqrt{2\log m}}e,\qquad
 \E H_*\sim\sqrt{2\log m},\qquad \E H_*^2=O(\log m).
 \label{sym:eq:gaussian-max}
\end{equation}
The accompanying smoothness bound is
\begin{equation}
 \|F_m''\|_\infty\le C\log(em),\qquad m\ge2.
 \label{sym:eq:gaussian-smoothness}
\end{equation}
The normal tail estimates
\(\bar\Phi(x)\asymp\varphi(x)/(1+x)\) for \(x\ge0\) and
\(\varphi(x)/\bar\Phi(x)\sim x\) imply the first two assertions,
using \(F_m'(u_m)=m\varphi(u_m)(1-1/m)^{m-1}\).
They also give \(H_*/\sqrt{2\log m}\to1\) in probability.
Integrating \(\Pp(\max_k|H_k|>y)\le\min\{1,2me^{-y^2/2}\}\)
gives the second-moment bound and uniform integrability needed for
the mean asymptotic. Finally,
\[
 F_m''(x)=m(m-1)\varphi(x)^2\Phi(x)^{m-2}
             -mx\varphi(x)\Phi(x)^{m-1}.
\]
For \(x<0\), the powers of \(\Phi(x)\le1/2\) bound both terms
by an absolute constant, treating \(m=2,3\) separately. For
\(0\le x\le2\sqrt{\log(em)}\), put \(w=m\bar\Phi(x)\) and use
\(\varphi(x)\le C(1+x)\bar\Phi(x)\). The resulting bound is
\(C(1+x)^2(w+w^2)e^{-c_*w}\le C\log(em)\), again treating
\(m=2,3\) separately. For larger \(x\), the bound
\(C(m^2e^{-x^2}+mx e^{-x^2/2})\le C\log(em)\) suffices.

\subsection{The symmetric construction and its discrepancy}

We now construct the observations in Theorem~\ref{opt:symmetric-lower},
modifying the common activation example in~\cite[Proposition~E.1]{FKLZ}.
Suppress the subscript \(n\): let \(m\to\infty\), put
\(\ell\coloneqq\log m\), and assume \(\ell^3/n\to0\).
With \(q\coloneqq\ell/(2n)\) and \(s\coloneqq\sqrt{1-q}\), define
\[
 X_{i,jk}\coloneqq s\eta_{ij}+I_i\varepsilon_{ik},\qquad 1\le j,k\le m.
\]
Here all \(\eta_{ij},\varepsilon_{ik}\) are independent Rademacher
signs, and the \(I_i\) are independent Bernoulli\((q)\) variables,
independent of the signs. All statements concern sufficiently large
\(n\), so \(0<q<1/2\).
The rows are independent and identically distributed and bounded by
\(s+1\le2\). Reversing all signs in one row preserves its law and
negates that row, so its distribution is centrally symmetric and
its third-moment tensor vanishes. The covariance matrix is
\[
 \Sigma_{(j,k),(j',k')}=(1-q)\ind\{j=j'\}+q\ind\{k=k'\}.
\]
For \(m\ge2\) and \(0<q<1/2\), its largest off-diagonal correlation
is \(1-q\), so the correlation gap is \(q\).
Its range consists of arrays \((a_j+b_k)_{j,k}\), a space of
dimension \(2m-1\); both components have positive variance, so this
is also its rank.

To compare the maxima, separate the always-present sign component
from the component contributed by the active rows. Define
\[
 A_j\coloneqq n^{-1/2}\sum_i\eta_{ij},\qquad
 T_k\coloneqq\sum_i I_i\varepsilon_{ik},\qquad
 T_*\coloneqq\max_{k\le m}T_k.
\]
For \(W=n^{-1/2}\sum_iX_i\) and a covariance-matched Gaussian
vector \(Z\), the maximum representations are
\[
 \max_{j,k}W_{jk}=s\max_jA_j+T_*/\sqrt n,\qquad
 \max_{j,k}Z_{jk}\stackrel d=s\max_jG_j+\sqrt q\,H_*,
\]
where \(G_j,H_k\) are mutually independent standard normals and
\(H_*=\max_kH_k\). In each representation the two components are
independent.

With many columns, one column is likely to be positive on every
active row, making \(T_*\) equal to the number \(N\coloneqq\sum_iI_i\)
of active rows. Here \(\E N=\ell/2\), \(\operatorname{Var}(N)\le\ell/2\),
and \(\E N^2=O(\ell^2)\). Conditional on the active rows, each
column is positive on all of them with probability \(2^{-N}\),
independently across columns. Thus, using Chebyshev's bound
\(\Pp(N>\ell)\le2/\ell\),
\begin{align*}
 \Pp(T_*\ne N\mid N=k)&=(1-2^{-k})^m\le e^{-m2^{-k}},\\
 \Pp(T_*\ne N)&\le\frac2\ell+e^{-m2^{-\ell}}=o(1),
\end{align*}
where \(m2^{-\ell}=\exp\{(1-\log2)\ell\}\to\infty\).
Since \(|T_*|\le N\),
Cauchy--Schwarz then yields
\[
 0\le\E(N-T_*)\le2(\E N^2)^{1/2}
                         \Pp(T_*\ne N)^{1/2}=o(\ell).
\]
Consequently, \(\E T_*=(1/2+o(1))\ell\) and
\(\E T_*^2=O(\ell^2)\).

We now replace the leading Rademacher maximum by a Gaussian maximum
at a cost smaller than the discrepancy we seek. Let \(F_{n,m}\) be
the distribution function of \(\max_jA_j\), and put
\(r_n\coloneqq\ell^{3/2}/\sqrt n\). The assumed condition
\(\ell^3/n\to0\), together with the Rademacher-maximum
estimate~\eqref{sym:eq:rademacher-max}, gives
\[
 \|F_{n,m}-F_m\|_\infty
 \le C\left(\frac{\sqrt\ell}{\sqrt n}+\frac{\ell^2}{n}\right)
 =o(r_n).
\]
At the threshold \(t_n\coloneqq s u_m\), define the shifts
\(V_W\coloneqq T_*/(s\sqrt n)\) and \(V_G\coloneqq\sqrt q\,H_*/s\).
Independence in the two maximum representations gives
\[
 \Pp(\max_{j,k}W_{jk}\le t_n)=\E F_m(u_m-V_W)+o(r_n),
 \qquad
 \Pp(\max_{j,k}Z_{jk}\le t_n)=\E F_m(u_m-V_G).
\]
Since \(s\to1\), the Gaussian moment estimates
in~\eqref{sym:eq:gaussian-max} and the bounds for \(T_*\) give
\[
 \E V_G-\E V_W=(1/2+o(1))\frac\ell{\sqrt n},
 \qquad \E V_G^2+\E V_W^2=O(\ell^2/n).
\]
The mean-shift difference is of order \(\ell/\sqrt n\); multiplying
it by \(F_m'(u_m)\asymp\sqrt\ell\) gives the desired order \(r_n\).
To control the remainder, use the global bound
\eqref{sym:eq:gaussian-smoothness}, which also covers negative shifts:
\[
 \E F_m(u_m-V)=F_m(u_m)-F_m'(u_m)\E V+O(\ell\E V^2).
\]
Subtracting the two expansions therefore gives
\[
 \Pp(\max_{j,k}W_{jk}\le t_n)-\Pp(\max_{j,k}Z_{jk}\le t_n)
 =\left(\frac1{e\sqrt2}+o(1)\right)r_n+O(r_n^2).
\]
Finally, recall that \(\delta_n=2/\sqrt n\) and
\(L_n=\log(em^2)=1+2\ell\). Since \(r_n\to0\) and
\(\delta_nL_n^{3/2}=2(1+2\ell)^{3/2}/\sqrt n\sim2^{5/2}r_n\),
the displayed lower orthant has discrepancy
\(\{1/(8e)+o(1)\}\delta_nL_n^{3/2}\), proving Theorem~\ref{opt:symmetric-lower}.


\end{document}

